% Leçon 40 — Expression écrite : caractères complexes liés à la rédaction d'un CV — Chinois Tle A — Cours signé OnBuch+
% Programme officiel MINESEC. Compiler avec : tectonic ZH40-ecrit-caracteres-complexes-lies-a-la-redaction-d-u.tex
\newcommand{\DOCMATIERE}{Chinois}
\newcommand{\DOCNIVEAU}{Tle A}
\newcommand{\DOCMODULE}{Module 5 : Mass médias et télécommunication}
\newcommand{\DOCLECON}{ZH40}
\newcommand{\DOCTITRE}{Leçon 40 — Expression écrite : caractères complexes liés à la rédaction d'un CV}
% OnBuch+ — préambule « cours signés » ENRICHI (compilé avec tectonic / XeLaTeX)
% Système visuel « Pop » d'OnBuch+ : fond crème (jamais blanc), bordures encre
% épaisses, ombres « dures » décalées, titres Archivo Black en capitales.
% Version MATHÉMATIQUES : graphes (pgfplots/tikz), boîtes pédagogiques
% (définition, propriété, méthode, attention, le savais-tu, exercice résolu,
% exercice, TP, bilan), page de garde et signature OnBuch+ ancrée en bas.
%
% Macros attendues dans main.tex AVANT \input{preamble} :
%   \DOCMATIERE  \DOCNIVEAU  \DOCMODULE  \DOCLECON (numéro)  \DOCTITRE
\documentclass[11pt]{article}

\usepackage{fontspec}
\usepackage[french]{babel} % français = langue principale ; le chinois via \zh{..} / textebox (xeCJK)
\usepackage{xeCJK}
\usepackage{pifont} % \ding{51} \ding{55} utilisés par les modèles
\usepackage[a4paper,margin=2cm,top=3.4cm,bottom=2.8cm,headheight=1.4cm,footskip=1.3cm]{geometry}
\usepackage{amsmath,amssymb,amsfonts}
\usepackage{mathtools} % \xmapsto, \coloneqq... souvent utilisés par les modèles
\usepackage{cancel} % simplifications barrées (bilans de forces)
% \abs{z} (module, valeur absolue) et \norm{u} (norme), utilisés par les modèles
\providecommand{\abs}{}\let\abs\relax\DeclarePairedDelimiter{\abs}{\lvert}{\rvert}
\providecommand{\norm}{}\let\norm\relax\DeclarePairedDelimiter{\norm}{\lVert}{\rVert}
\usepackage[table]{xcolor} % [table] : fournit aussi \rowcolors (alternance de lignes)
\usepackage{pagecolor}
\usepackage{tikz}
\usetikzlibrary{arrows.meta,positioning,calc,shapes.geometric,shapes.misc,shapes.arrows,decorations.pathmorphing,decorations.pathreplacing,decorations.markings,patterns,shadows,mindmap,trees,fit,backgrounds,matrix,intersections,angles,quotes}
\usepackage{pgfplots}
\usepackage[version=4]{mhchem} % équations nucléaires \ce{^{235}_{92}U}
\usepackage[european,siunitx]{circuitikz} % schémas électriques
\pgfplotsset{compat=1.18}
\usepgfplotslibrary{fillbetween,groupplots} % aires entre courbes (calcul intégral)
\usepackage{enumitem}
\usepackage{fancyhdr}
% [most] charge toutes les bibliothèques tcolorbox (listings, minted, raster,
% documentation...) et gonfle inutilement la mémoire TeX sur un document long
% (~300 boîtes) : on ne charge que ce qui est réellement utilisé.
\usepackage[skins,breakable]{tcolorbox}
\usepackage{graphicx}
\usepackage{colortbl}
\usepackage{booktabs}
\usepackage{tabularx}
\usepackage{diagbox} % case d'en-tête diagonale des tableaux à double entrée
% Colonnes extensibles centrée / gauche / droite (C, L, R), souvent utilisées par les modèles
\newcolumntype{C}{>{\centering\arraybackslash}X}
\newcolumntype{L}{>{\raggedright\arraybackslash}X}
\newcolumntype{R}{>{\raggedleft\arraybackslash}X}
\usepackage{array}
\usepackage{multicol}
\usepackage{multirow}
\usepackage{siunitx}
% parse-numbers=false : accepte aussi \SI{2,5 \times 10^{-3}}{...}
\sisetup{locale=FR,output-decimal-marker={,},group-minimum-digits=4,parse-numbers=false}
\DeclareSIUnit\ohm{\ensuremath{\symup{\Omega}}} % U+2126 absent de la police
\DeclareSIUnit\division{div} % réglages d'oscilloscope (ms/div, V/div)
\DeclareSIUnit\tr{tr} % tours (vitesse de rotation, ex. \SI{1500}{\tr\per\minute})
\DeclareSIUnit\molar{M} % concentration molaire (ex. \SI{3}{\molar}, \SI{1}{\micro\molar})
\DeclareSIUnit\an{an} % année (le modèle confond parfois avec \year, un compteur TeX) — ex. \SI{5}{\kilo\meter\per\an}
\DeclareSIUnit\FCFA{FCFA} % franc CFA (ex. \SI{500}{\FCFA\per\kilo\gram})
\DeclareSIUnit\degreeBrix{\textdegree Brix} % teneur en sucre d'un sirop/jus (réfractométrie)
\DeclareSIUnit\month{mois} % le modèle utilise parfois \month, un compteur TeX, comme unité de durée
\usepackage{qrcode}
% \degree : commande fréquemment utilisée par les modèles pour un angle ou une
% température en dehors de \SI{}{} (non fournie par les paquets chargés).
\providecommand{\degree}{^{\circ}}
% \qed (fin de démonstration), utilisé par les modèles sans amsthm
\providecommand{\qed}{\hfill$\square$}
% \barre : double barre des valeurs interdites dans un tableau de variations (tkz-tab)
\providecommand{\barre}{\ensuremath{\|}}
% environnement proof (amsthm non chargé) utilisé par les modèles
\ifdefined\proof\else\newenvironment{proof}[1][Démonstration]{\par\noindent\textit{#1.}\ }{\qed\par}\fi
\usepackage{lastpage}
\usepackage{hyperref}
\hypersetup{hidelinks}

% ============================================================
% Palette « Pop » OnBuch+ — mêmes hex que la classe P (plus_ui.dart)
% ============================================================
\definecolor{poporange}{HTML}{F59321}
\definecolor{popdark}{HTML}{E07A0C}
\definecolor{popcream}{HTML}{FFF6E8}
\definecolor{popink}{HTML}{1C1714}
\definecolor{poppurple}{HTML}{7A5AE0}
\definecolor{popgreen}{HTML}{2E9E6A}
\definecolor{poppink}{HTML}{E0457B}
\definecolor{popblue}{HTML}{2F7DE1}
\definecolor{popgold}{HTML}{C98A12}
\definecolor{popmuted}{HTML}{8A7F76}
\definecolor{popline}{HTML}{EADFD2}
\definecolor{popgray}{HTML}{8A7F76}
\colorlet{popgrey}{popgray}
% Alias tolérants (le modèle nomme parfois les couleurs différemment)
\colorlet{popwhite}{white}
\colorlet{popblack}{popink}
\colorlet{popred}{poppink}
\colorlet{popyellow}{popgold}
% Teintes claires (fonds de tableaux/encadrés) — le modèle nomme parfois
% les couleurs avec un suffixe L (ex. popblueL) sans les avoir définies.
\colorlet{poporangeL}{poporange!20}
\colorlet{popdarkL}{popdark!20}
\colorlet{poppurpleL}{poppurple!20}
\colorlet{popgreenL}{popgreen!20}
\colorlet{poppinkL}{poppink!20}
\colorlet{popblueL}{popblue!20}
\colorlet{popgoldL}{popgold!20}
\colorlet{popredL}{popred!20}
\colorlet{popgrayL}{popgray!20}
% Teintes claires pour fonds de tableaux / zones
\colorlet{popblueL}{popblue!10!white}
\colorlet{poporangeL}{poporange!14!white}
\colorlet{popgreenL}{popgreen!12!white}
\colorlet{poppurpleL}{poppurple!12!white}
\colorlet{poppinkL}{poppink!12!white}
\colorlet{popgoldL}{popgold!14!white}

\pagecolor{popcream}

% ============================================================
% Typographie
% ============================================================
\defaultfontfeatures{Ligatures=TeX}
\setmainfont{PlusJakartaSans-Medium.ttf}[Path=../../../fonts/,BoldFont=PlusJakartaSans-Bold.ttf]
% Symboles, API et emojis absents de la police principale : repli sur DejaVu Sans (caractères actifs)
\newfontface\symfont{DejaVuSans.ttf}[Path=../../../fonts/]
\makeatletter
\newcommand\symactive[1]{\catcode#1=13 \begingroup\lccode`\~=#1 \lowercase{\endgroup\def~}{{\symfont\char#1}}}
\newcommand\symblank[1]{\catcode#1=13 \begingroup\lccode`\~=#1 \lowercase{\endgroup\def~}{}}
\newcommand\symhyph[1]{\catcode#1=13 \begingroup\lccode`\~=#1 \lowercase{\endgroup\def~}{\mbox{-}}}
\newcount\symc
\newcommand\symrange[2]{\symc=#1 \@whilenum\symc<#2 \do{\expandafter\symactive\expandafter{\the\symc}\advance\symc by 1 }}
\newcommand\symblankrange[2]{\symc=#1 \@whilenum\symc<#2 \do{\expandafter\symblank\expandafter{\the\symc}\advance\symc by 1 }}
\makeatother
\symrange{"0250}{"02FF}\symactive{"2713}\symactive{"2714}\symactive{"2717}\symactive{"2718}\symactive{"2610}\symactive{"2611}\symactive{"2612}
\symrange{"2600}{"27BF}
\catcode"2705=13 \begingroup\lccode`\~="2705 \lowercase{\endgroup\def~}{{\symfont\char"2713}}\symblank{"26BD}\symblank{"26A1}
\symrange{"2460}{"257F}\symactive{"223C}\symactive{"2248}\symactive{"2260}
\symactive{"01F8}\symactive{"01F9}\symactive{"1E3E}\symactive{"1E3F}
\symhyph{"2011}\symhyph{"2E1A}\symblankrange{"1F300}{"1FAFF}\symblank{"FE0F}
% Caractères chinois : WenQuanYi Zen Hei (sous-ensemble GB2312 + pinyin), gras/italique simulés
\setCJKmainfont{WenQuanYiZenHei-subset.ttf}[Path=../../../fonts/,AutoFakeBold=3,AutoFakeSlant=0.2]
\xeCJKsetup{CJKspace=false}
% Pinyin : voyelles à caron (ǎ ǐ ǒ ǔ ǖ ǘ ǚ ǜ) absentes de la police principale -> caractères actifs
\newfontface\pyfont{WenQuanYiZenHei-subset.ttf}[Path=../../../fonts/]
\newcommand\pyactive[1]{\catcode#1=13 \begingroup\lccode`\~=#1 \lowercase{\endgroup\def~}{{\pyfont\char#1}}}
\pyactive{"01CD}\pyactive{"01CE}\pyactive{"01CF}\pyactive{"01D0}\pyactive{"01D1}\pyactive{"01D2}\pyactive{"01D3}\pyactive{"01D4}\pyactive{"01D5}\pyactive{"01D6}\pyactive{"01D7}\pyactive{"01D8}\pyactive{"01D9}\pyactive{"01DA}\pyactive{"01DB}\pyactive{"01DC}
% Maths en sans-serif (Fira Math) avec chiffres et lettres droites de Plus
% Jakarta, pour que les formules se fondent dans le texte.
\usepackage[math-style=ISO,bold-style=ISO]{unicode-math}
\setmathfont{FiraMath-Regular.otf}[Path=../../../fonts/]
\setmathfont{PlusJakartaSans-Medium.ttf}[Path=../../../fonts/,range={up/num,up/{Latin,latin},\mathup}]
\usepackage{icomma} % virgule décimale sans espace en mode math
% Caractères mathématiques Unicode absents des polices de texte (Plus Jakarta,
% Archivo Black) : rendus par leur équivalent mathématique, en texte comme en
% formule — sinon ils s'affichent en carré vide (ex. « Arithmétique dans ℤ »).
\usepackage{newunicodechar}
\newunicodechar{ℝ}{\ensuremath{\mathbb{R}}}
\newunicodechar{ℤ}{\ensuremath{\mathbb{Z}}}
\newunicodechar{ℂ}{\ensuremath{\mathbb{C}}}
\newunicodechar{ℕ}{\ensuremath{\mathbb{N}}}
\newunicodechar{ℚ}{\ensuremath{\mathbb{Q}}}
\newunicodechar{𝔻}{\ensuremath{\mathbb{D}}}
\newunicodechar{α}{\ensuremath{\alpha}}
\newunicodechar{β}{\ensuremath{\beta}}
\newunicodechar{γ}{\ensuremath{\gamma}}
\newunicodechar{δ}{\ensuremath{\delta}}
\newunicodechar{ε}{\ensuremath{\varepsilon}}
\newunicodechar{θ}{\ensuremath{\theta}}
\newunicodechar{λ}{\ensuremath{\lambda}}
\newunicodechar{μ}{\ensuremath{\mu}}
\newunicodechar{σ}{\ensuremath{\sigma}}
\newunicodechar{τ}{\ensuremath{\tau}}
\newunicodechar{φ}{\ensuremath{\varphi}}
\newunicodechar{ψ}{\ensuremath{\psi}}
\newunicodechar{ω}{\ensuremath{\omega}}
\newunicodechar{Σ}{\ensuremath{\Sigma}}
% majuscules produites par \MakeUppercase dans les titres (α-aminés -> Α)
\newunicodechar{Α}{\ensuremath{\alpha}}
\newunicodechar{Β}{\ensuremath{\beta}}
\newunicodechar{Γ}{\ensuremath{\Gamma}}
\newunicodechar{Δ}{\ensuremath{\Delta}}
\newunicodechar{Ε}{\ensuremath{\varepsilon}}
\newunicodechar{Θ}{\ensuremath{\Theta}}
\newunicodechar{Λ}{\ensuremath{\Lambda}}
\newunicodechar{Μ}{\ensuremath{\mu}}
\newunicodechar{Τ}{\ensuremath{\tau}}
\newunicodechar{Φ}{\ensuremath{\Phi}}
\newunicodechar{Ψ}{\ensuremath{\Psi}}
\newunicodechar{Ω}{\ensuremath{\Omega}}
\newunicodechar{↦}{\ensuremath{\mapsto}}
\newunicodechar{∈}{\ensuremath{\in}}
\newunicodechar{∉}{\ensuremath{\notin}}
\newunicodechar{∧}{\ensuremath{\wedge}}
\newunicodechar{⇔}{\ensuremath{\Leftrightarrow}}
\newunicodechar{⟺}{\ensuremath{\Longleftrightarrow}}
\newunicodechar{⇒}{\ensuremath{\Rightarrow}}
\newunicodechar{⟹}{\ensuremath{\Longrightarrow}}
\newunicodechar{∘}{\ensuremath{\circ}}
\newunicodechar{∪}{\ensuremath{\cup}}
\newunicodechar{∩}{\ensuremath{\cap}}
\newunicodechar{≡}{\ensuremath{\equiv}}
\newunicodechar{⊂}{\ensuremath{\subset}}
\newunicodechar{⊥}{\ensuremath{\perp}}
\newunicodechar{∃}{\ensuremath{\exists}}
\newunicodechar{∀}{\ensuremath{\forall}}
\newunicodechar{∛}{\ensuremath{\sqrt[3]{\,}}}
\newunicodechar{∜}{\ensuremath{\sqrt[4]{\,}}}
\newunicodechar{⃗}{\ensuremath{{}^{\to}}}
\newunicodechar{ⁿ}{\ifmmode^{n}\else\textsuperscript{\ensuremath{n}}\fi}
\newunicodechar{ˣ}{\ifmmode^{x}\else\textsuperscript{\ensuremath{x}}\fi}
\newunicodechar{ᵇ}{\ifmmode^{b}\else\textsuperscript{\ensuremath{b}}\fi}
\newunicodechar{ʳ}{\ifmmode^{r}\else\textsuperscript{\ensuremath{r}}\fi}
\newunicodechar{ᵘ}{\ifmmode^{u}\else\textsuperscript{\ensuremath{u}}\fi}
\newunicodechar{ʸ}{\ifmmode^{y}\else\textsuperscript{\ensuremath{y}}\fi}
\newunicodechar{ᵃ}{\ifmmode^{a}\else\textsuperscript{\ensuremath{a}}\fi}
\newunicodechar{ᵉ}{\ifmmode^{e}\else\textsuperscript{\ensuremath{e}}\fi}
\newunicodechar{ᵏ}{\ifmmode^{k}\else\textsuperscript{\ensuremath{k}}\fi}
\newunicodechar{ᵖ}{\ifmmode^{p}\else\textsuperscript{\ensuremath{p}}\fi}
\newunicodechar{ᴺ}{\ifmmode^{N}\else\textsuperscript{\ensuremath{N}}\fi}
\newunicodechar{ᶜ}{\ifmmode^{c}\else\textsuperscript{\ensuremath{c}}\fi}
\newunicodechar{ʰ}{\ifmmode^{h}\else\textsuperscript{\ensuremath{h}}\fi}
\newunicodechar{ᵠ}{\ifmmode^{q}\else\textsuperscript{\ensuremath{q}}\fi}
\newunicodechar{ₙ}{\ifmmode_{n}\else\textsubscript{\ensuremath{n}}\fi}
\newunicodechar{ₐ}{\ifmmode_{a}\else\textsubscript{\ensuremath{a}}\fi}
\newunicodechar{ᵢ}{\ifmmode_{i}\else\textsubscript{\ensuremath{i}}\fi}
\newunicodechar{ᵣ}{\ifmmode_{r}\else\textsubscript{\ensuremath{r}}\fi}
\newunicodechar{ₖ}{\ifmmode_{k}\else\textsubscript{\ensuremath{k}}\fi}
\newunicodechar{ᵦ}{\ifmmode_{\beta}\else\textsubscript{\ensuremath{\beta}}\fi}
\newunicodechar{ₓ}{\ifmmode_{x}\else\textsubscript{\ensuremath{x}}\fi}
\newfontfamily\popdisplay{ArchivoBlack-Regular.ttf}[Path=../../../fonts/]
\newfontfamily\poplogo{SpaceGrotesk-Bold.ttf}[Path=../../../fonts/]
\newfontfamily\popsemi{PlusJakartaSans-SemiBold.ttf}[Path=../../../fonts/]
\newfontfamily\popxbold{PlusJakartaSans-ExtraBold.ttf}[Path=../../../fonts/]
\newfontfamily\popxboldls{PlusJakartaSans-ExtraBold.ttf}[Path=../../../fonts/,LetterSpace=3.0]

\setlist[enumerate]{leftmargin=1.7em,itemsep=5pt,topsep=4pt}
\setlist[itemize]{leftmargin=1.7em,itemsep=4pt,topsep=4pt,label={\color{poporange}\textbullet}}
\setlength{\parindent}{0pt}
\setlength{\parskip}{5pt}
\renewcommand{\arraystretch}{1.25}

% pgfplots : style Pop par défaut
\pgfplotsset{
  every axis/.append style={axis line style={popink,line width=1pt},tick style={popink},
    grid style={popline,line width=0.5pt},label style={font=\small},tick label style={font=\footnotesize}},
  popcurve/.style={poporange,line width=1.8pt,no markers,smooth},
}
% Même style disponible dans un \draw TikZ simple (hors pgfplots)
\tikzset{popcurve/.style={poporange,line width=1.8pt}}
\tikzset{popgrid/.style={popline,very thin}}
\colorlet{popcurve}{poporange} % le nom du style est parfois utilisé comme couleur

% ============================================================
% Pill / squiggle
% ============================================================
\newcommand{\poppill}[3]{%
  \tikz[baseline=-0.55ex]\node[draw=popink,line width=1.2pt,rounded corners=2.6pt,%
    fill=#2,inner xsep=6.5pt,inner ysep=3pt]{%
    \popxboldls\fontsize{7}{7}\selectfont\color{#3}\MakeUppercase{#1}};%
}
\newcommand{\popsquiggle}[1]{%
  \tikz[baseline=-0.4ex]{
    \draw[#1,line width=2.6pt,line cap=round]
      plot [smooth] coordinates {(0,0.09) (0.34,0.01) (0.68,0.21) (1.02,0.01) (1.36,0.10)};
  }%
}
% Mot-clé surligné (marqueur orange sous le texte)
\newcommand{\cle}[1]{{\popxbold\color{popdark}#1}}

% ============================================================
% En-tête / pied
% ============================================================
\pagestyle{fancy}
\fancyhf{}
\renewcommand{\headrulewidth}{0pt}
\renewcommand{\footrulewidth}{0pt}
\lhead{\raisebox{-0.15ex}{\poplogo\fontsize{13}{13}\selectfont\color{popink}OnBuch\color{poporange}+}}
\rhead{\poppill{\DOCMATIERE\ \textbullet\ \DOCNIVEAU\ \textbullet\ Leçon \DOCLECON}{white}{popink}}
\lfoot{\popxbold\fontsize{7.6}{7.6}\selectfont\color{popmuted}\MakeUppercase{Cours signé OnBuch+}\ \textbullet\ {\popsemi\DOCTITRE}}
\rfoot{\tikz[baseline=-0.6ex]\node[rounded corners=8.5pt,fill=popink,minimum height=17pt,minimum width=17pt,inner xsep=5pt,inner sep=0]{\popxbold\fontsize{8}{8}\selectfont\color{popcream}\thepage\,/\,\pageref*{LastPage}};}

% ============================================================
% Titres
% ============================================================
\newcommand{\chaptitre}[2]{%
  \begin{tcolorbox}[enhanced,colback=poporange,colframe=popink,boxrule=2.6pt,arc=11pt,
      drop shadow=popink,left=15pt,right=15pt,top=12pt,bottom=11pt,before skip=3pt,after skip=18pt]
    \poppill{#2}{popink}{popcream}\par\vspace{9pt}
    {\raggedright\hyphenpenalty=10000\exhyphenpenalty=10000\popdisplay\fontsize{22}{24}\selectfont\color{popcream}\MakeUppercase{#1}\par}
  \end{tcolorbox}
}
% Section numérotée : pastille orange avec le numéro + titre Archivo
\newcounter{popsec}
\newcommand{\coursec}[1]{%
  \stepcounter{popsec}\setcounter{popsub}{0}%
  \par\vspace{16pt}\noindent
  \begin{minipage}{\linewidth}
  \tikz[baseline=-0.7ex]\node[draw=popink,line width=1.6pt,fill=poporange,rounded corners=4pt,minimum size=22pt,inner sep=0]{\popdisplay\fontsize{12}{12}\selectfont\color{popcream}\thepopsec};%
  \hspace{8pt}{\popdisplay\fontsize{15}{17}\selectfont\color{popink}\MakeUppercase{#1}}\par
  \vspace{4pt}\noindent{\color{poporange}\rule{36pt}{3.6pt}}
  \end{minipage}\par\nopagebreak\vspace{8pt}
}
\newcounter{popsub}
\newcommand{\courssub}[1]{%
  \stepcounter{popsub}%
  \par\vspace{9pt}\noindent{\popxbold\fontsize{11.5}{13}\selectfont\color{popink}{\color{poporange}\thepopsec.\thepopsub}\hspace{6pt}#1}\par\nopagebreak\vspace{4pt}
}

% ============================================================
% Boîtes pédagogiques — même recette : fond blanc, bordure épaisse colorée,
% ombre dure encre, badge plein. Titre optionnel [#1].
% ============================================================
\tcbset{popbase/.style={breakable,enhanced,colback=white,boxrule=2.3pt,arc=9pt,
  shadow={3pt}{-3pt}{0pt}{popink},left=14pt,right=14pt,top=11pt,bottom=11pt,before skip=11pt,after skip=15pt,
  fontupper=\popsemi\fontsize{10.6}{14.5}\selectfont\color{popink},
  fontlower=\popsemi\fontsize{10.6}{14.5}\selectfont\color{popink}}}
% Coche dessinée (le glyphe ✓ n'existe pas dans les polices du document)
\newcommand{\popcheck}[1]{\tikz[baseline=-0.5ex]\draw[#1,line width=1.6pt,line cap=round,line join=round](-0.13,0.0)--(-0.04,-0.09)--(0.13,0.1);}
\providecommand{\coche}{\popcheck{popgreen}}
\providecommand{\resultat}[1]{\par\smallskip\noindent\fcolorbox{popgreen}{popgreenL}{\parbox{\dimexpr\linewidth-2\fboxsep-2\fboxrule\relax}{\textbf{Résultat.} #1}}\par\smallskip}
\newcommand{\popboxhead}[3]{\poppill{#1}{#2}{white}\ifx\relax#3\relax\else\hspace{6pt}{\popxbold\fontsize{10.6}{12}\selectfont\color{popink}#3}\fi\par\vspace{7pt}}

\newtcolorbox{aretenir}[1][]{popbase,colframe=popgreen,before upper={\popboxhead{À retenir}{popgreen}{#1}}}
% Texte en chinois (document de lecture, dialogue, extrait) : cadre
% d'étude. Un titre optionnel (source, auteur) s'affiche dans l'en-tête.
\newtcolorbox{textebox}[1][]{popbase,colframe=popblue,before upper={\popboxhead{文本 · Texte}{popblue}{#1}},after upper={}}
% Marques ✓ / ✗ (forme correcte / fautive) souvent utilisées par les modèles
\providecommand{\cross}{\textcolor{poppink}{\ensuremath{\times}}}
\providecommand{\xmark}{\cross}\providecommand{\crossmark}{\cross}\providecommand{\cmark}{\ensuremath{\checkmark}}
% Ligne de réponse / justification à compléter (inventée par les modèles)
\providecommand{\justification}[1]{\par\noindent\textit{Justification~:} \dotfill\par}
% \textipa (package tipa non chargé) : la police ne couvre pas l'API -> simple passage
\providecommand{\textipa}[1]{#1}
% Soulignement (ulem non chargé) : \uline, \ul
\providecommand{\uline}[1]{\underline{#1}}\providecommand{\ul}[1]{\underline{#1}}
% \glossaire{\item ...} inventée par les modèles : liste de glossaire
\providecommand{\glossaire}[1]{\begin{itemize}#1\end{itemize}}
% \alert (beamer) inventée par les modèles : texte mis en évidence
\providecommand{\alert}[1]{\textbf{\textcolor{poppink}{#1}}}
% variantes ASCII d'ordinaux écrites en macro
\providecommand{\iere}{\textsuperscript{re}}\providecommand{\ieme}{\textsuperscript{e}}\providecommand{\ere}{\textsuperscript{re}}\providecommand{\eme}{\textsuperscript{e}}\providecommand{\er}{\textsuperscript{er}}\providecommand{\ie}{\textsuperscript{e}}
% Ordinaux français écrits en macro par les modèles : 1\ière, 3\ième
\newcommand{\ière}{\textsuperscript{re}}\newcommand{\ième}{\textsuperscript{e}}
% \exemple (inventée par les modèles) : étiquette « Exemple : »
\providecommand{\exemple}{\textbf{Exemple~:}\ }
% \square utilisable aussi hors mode math (cases à cocher)
\AtBeginDocument{\let\squareMath\square\renewcommand{\square}{\ensuremath{\squareMath}}}
% Mot ou phrase en chinois à l'intérieur d'un paragraphe français
\newcommand{\zh}[1]{\ifmmode\text{#1}\else{#1}\fi}
\let\eng\zh \let\esp\zh \let\all\zh % alias tolérés
% flèches d'intonation inventées par les modèles
\providecommand{\textsearrow}{}\renewcommand{\textsearrow}{\ensuremath{\searrow}}\providecommand{\textnearrow}{}\renewcommand{\textnearrow}{\ensuremath{\nearrow}}
% \fr{..} : mot français (inventé par les modèles)
\providecommand{\fr}[1]{\textit{#1}}
% zhbox : encadré de texte chinois inventé par les modèles
\newenvironment{zhbox}{\begin{quote}}{\end{quote}}
% \courssubsub : titre de niveau 3 inventé par les modèles
\providecommand{\courssubsub}[1]{\par\medskip\noindent\textbf{#1}\par\smallskip}
% \zhbf : texte chinois en gras inventé par les modèles
\providecommand{\zhbf}[1]{\textbf{#1}}
% \vs : « versus » inventé par les modèles
\providecommand{\vs}{\textit{vs}\ }
% \blank : case à compléter inventée par les modèles
\providecommand{\blank}[1][]{\underline{\hspace{1.6em}}}
\newtcolorbox{exemplebox}[1][]{popbase,colframe=poporange,before upper={\popboxhead{Exemple}{poporange}{#1}}}
\newtcolorbox{definition}[1][]{popbase,colframe=popblue,before upper={\popboxhead{Définition}{popblue}{#1}}}
\newtcolorbox{propriete}[1][]{popbase,colframe=poppurple,before upper={\popboxhead{Propriété}{poppurple}{#1}}}
\newtcolorbox{methode}[1][]{popbase,colframe=popgold,before upper={\popboxhead{Méthode}{popgold}{#1}}}
\newtcolorbox{attention}[1][]{popbase,colframe=poppink,before upper={\popboxhead{Attention}{poppink}{#1}}}
\newtcolorbox{experience}[1][]{popbase,colframe=popblue,colback=popblueL,before upper={\popboxhead{Expérience}{popblue}{#1}}}
% « Le savais-tu ? » : carte violette pleine (contexte, culture, Cameroun)
\newtcolorbox{savaistu}[1][]{popbase,colback=poppurple,colframe=popink,
  fontupper=\popsemi\fontsize{10.4}{14.2}\selectfont\color{white},
  before upper={\poppill{Le savais-tu\,?}{white}{poppurple}\ifx\relax#1\relax\else\hspace{6pt}{\popxbold\color{white}#1}\fi\par\vspace{7pt}}}
% Exercice résolu : énoncé en haut, \tcblower puis la solution sur fond crème
\newcounter{popexo}\newcounter{popexr}
\newtcolorbox{exoresolu}[1][]{popbase,colframe=poporange,colbacklower=popcream,
  segmentation style={popink,line width=1.2pt,dashed},
  before upper={\stepcounter{popexr}\popboxhead{Exercice résolu \thepopexr}{poporange}{#1}},
  before lower={\poppill{Solution}{popink}{popcream}\par\vspace{6pt}}}
% Exercice (non résolu en place) — la correction est dans la partie Corrigés
\newtcolorbox{exercice}[1][]{popbase,colframe=popink,
  before upper={\stepcounter{popexo}\popboxhead{Exercice \thepopexo}{popink}{#1}}}
% Corrigé
\newtcolorbox{corrige}[1][]{popbase,colframe=popgreen,colback=popgreenL,
  before upper={\popboxhead{Corrigé}{popgreen}{#1}}}
% Objectifs / prérequis / bilan
\newtcolorbox{objectifs}{popbase,colframe=popink,colback=poporangeL,
  before upper={\popboxhead{Objectifs de la leçon}{popink}{}}}
\newtcolorbox{prerequis}{popbase,colframe=popink,colback=popblueL,
  before upper={\popboxhead{Prérequis}{popblue}{}}}
\newtcolorbox{bilan}{popbase,colframe=popink,colback=white,boxrule=2.8pt,
  before upper={\popboxhead{Fiche bilan}{poporange}{L'essentiel à connaître pour l'examen}}}
% Figure encadrée avec légende
\newtcolorbox{popfigure}[1][]{enhanced,breakable=false,colback=white,colframe=popline,boxrule=1.4pt,arc=8pt,
  left=10pt,right=10pt,top=10pt,bottom=8pt,before skip=10pt,after skip=12pt,halign=center,
  lower separated=false,halign lower=center,
  fontlower=\popsemi\fontsize{9}{11.5}\selectfont\color{popmuted},
  #1}
\newcommand{\legende}[1]{\par\vspace{6pt}{\popsemi\fontsize{9}{11.5}\selectfont\color{popmuted}#1}}

% ============================================================
% Signature OnBuch+
% ============================================================
\newcommand{\poplogoinline}[1]{{\poplogo\fontsize{#1}{#1}\selectfont\color{popink}OnBuch\color{poporange}+}}
\newcommand{\signatureonbuch}{%
  \begin{tcolorbox}[enhanced,colback=popink,colframe=popink,boxrule=2.2pt,arc=10pt,
      drop shadow=poporange,left=14pt,right=14pt,top=10pt,bottom=10pt,before skip=0pt,after skip=0pt]
    \tikz[baseline=-0.7ex]\node[circle,fill=popgreen,draw=popcream,line width=1.3pt,minimum size=22pt,inner sep=0]{\popcheck{white}};%
    \hspace{9pt}\begin{minipage}[c]{0.62\linewidth}
      {\popxbold\fontsize{12}{13}\selectfont\color{popcream}Signé\ \textbullet\ {\poplogo OnBuch\color{poporange}+}}\par\vspace{2pt}
      {\raggedright\popsemi\fontsize{8}{10.5}\selectfont\color{popline}\MakeUppercase{Équipe pédagogique OnBuch+}\\\MakeUppercase{Conforme au programme officiel MINESEC}\par}
    \end{minipage}\hfill
    {\poplogo\fontsize{22}{22}\selectfont\color{popcream}OnBuch\color{poporange}+}
  \end{tcolorbox}%
}

% ============================================================
% Page de garde
% #1 titre · #2 matière · #3 niveau · #4 module · #5 n° leçon · #6 durée ·
% #7 description / objectifs (texte ou itemize)
% ============================================================
\newcommand{\pagegarde}[7]{%
  \thispagestyle{empty}
  \newgeometry{a4paper,margin=1.7cm,top=1.6cm,bottom=1.4cm}%
  \begin{tikzpicture}[remember picture,overlay]
    \fill[poporange!18] ([xshift=-3.2cm,yshift=-3.6cm]current page.north east) circle (4.6cm);
    \fill[poppurple!14] ([xshift=2.4cm,yshift=7.6cm]current page.south west) circle (3.8cm);
  \end{tikzpicture}%
  {\poplogo\fontsize{30}{30}\selectfont\color{popink}OnBuch\color{poporange}+}\hfill
  \raisebox{0.6ex}{\poppill{Cours signé \textbullet\ Premium}{popink}{popcream}}\par
  \vspace{1pt}\popsquiggle{poppurple}\par
  \vspace{8pt}
  \begin{tcolorbox}[enhanced,colback=poporange,colframe=popink,boxrule=2.8pt,arc=13pt,
      drop shadow=popink,left=18pt,right=18pt,top=12pt,bottom=13pt,after skip=10pt]
    \poppill{#2 \textbullet\ #3}{popink}{popcream}\hspace{5pt}\poppill{#4}{white}{popink}\par\vspace{8pt}
    {\popdisplay\fontsize{14}{14}\selectfont\color{popink}LEÇON #5}\par\vspace{4pt}
    {\raggedright\hyphenpenalty=10000\exhyphenpenalty=10000\popdisplay\fontsize{25}{27}\selectfont\color{popcream}\MakeUppercase{#1}\par}
    \vspace{8pt}
    {\popxbold\fontsize{9.5}{9.5}\selectfont\color{popink}\MakeUppercase{Durée conseillée : #6}}
  \end{tcolorbox}
  \begin{tcolorbox}[enhanced,colback=white,colframe=popink,boxrule=2.2pt,arc=10pt,
      drop shadow=popink,left=16pt,right=16pt,top=10pt,bottom=10pt,after skip=8pt]
    \poppill{Dans cette leçon}{poppurple}{white}\par\vspace{5pt}
    \popsemi\fontsize{9.3}{12.6}\selectfont\color{popink}#7
  \end{tcolorbox}
  \noindent\begin{minipage}[t]{0.487\linewidth}
    \begin{tcolorbox}[enhanced,colback=popgreen,colframe=popink,boxrule=2.2pt,arc=10pt,
        drop shadow=popink,left=11pt,right=11pt,top=10pt,bottom=10pt,height=3.1cm,valign=center]
      \tcbox[colback=white,colframe=popink,boxrule=1.3pt,arc=5pt,boxsep=0pt,nobeforeafter,
        left=3pt,right=3pt,top=3pt,bottom=3pt,valign=center]{\qrcode[height=1.9cm]{https://play.google.com/store/apps/details?id=com.luvvix.onbuch&hl=fr}}%
      \hspace{8pt}\begin{minipage}[c]{0.5\linewidth}
        {\popdisplay\fontsize{11}{12.5}\selectfont\color{white}\MakeUppercase{Télécharge OnBuch}}\par\vspace{4pt}
        {\raggedright\popsemi\fontsize{8.4}{11}\selectfont\color{white}Gratuit \textbullet\ Google Play\\Tuteur IA, annales, concours, résultats\par}
      \end{minipage}
    \end{tcolorbox}
  \end{minipage}\hfill
  \begin{minipage}[t]{0.487\linewidth}
    \begin{tcolorbox}[enhanced,colback=poppurple,colframe=popink,boxrule=2.2pt,arc=10pt,
        drop shadow=popink,left=11pt,right=11pt,top=10pt,bottom=10pt,height=3.1cm,valign=center]
      \tikz[baseline=-0.7ex]\node[circle,fill=white,draw=popink,line width=1.3pt,minimum size=1.6cm,inner sep=0]{\popdisplay\fontsize{24}{24}\selectfont\color{poppurple}+};%
      \hspace{8pt}\begin{minipage}[c]{0.6\linewidth}
        {\popdisplay\fontsize{11}{12.5}\selectfont\color{white}\MakeUppercase{Passe à OnBuch+}}\par\vspace{4pt}
        {\raggedright\popsemi\fontsize{8.4}{11}\selectfont\color{white}Tous les cours signés, Tuteur IA illimité, fiches PDF — onglet OnBuch+ de l'app\par}
      \end{minipage}
    \end{tcolorbox}
  \end{minipage}
  \vspace{8pt}
  \popcontenu
  \vfill
  \signatureonbuch
  \restoregeometry
  \newpage
}

% Rangée « Ce cours contient » (page de garde)
\newcommand{\poptile}[3]{% #1 couleur #2 gros texte #3 légende
  \begin{tcolorbox}[enhanced,colback=#1,colframe=popink,boxrule=2pt,arc=9pt,shadow={2.5pt}{-2.5pt}{0pt}{popink},
      left=6pt,right=6pt,top=9pt,bottom=9pt,halign=center,valign=center,height=1.75cm,width=\linewidth,nobeforeafter]
    {\popdisplay\fontsize{12.5}{13}\selectfont\color{white}\MakeUppercase{#2}}\par\vspace{3pt}
    {\centering\popsemi\fontsize{7.8}{9.5}\selectfont\color{white}#3\par}
  \end{tcolorbox}}
\newcommand{\popcontenu}{%
  \par\noindent{\popxboldls\fontsize{7.5}{7.5}\selectfont\color{popmuted}CE COURS SIGNÉ CONTIENT}\par\vspace{7pt}
  \noindent\begin{minipage}[t]{0.235\linewidth}\poptile{popblue}{Cours}{complet, illustré, pas à pas}\end{minipage}\hfill
  \begin{minipage}[t]{0.235\linewidth}\poptile{poporange}{Méthodes}{et exercices résolus}\end{minipage}\hfill
  \begin{minipage}[t]{0.235\linewidth}\poptile{poppink}{Exercices}{type examen + corrigés}\end{minipage}\hfill
  \begin{minipage}[t]{0.235\linewidth}\poptile{popgreen}{Fiche bilan}{l'essentiel à retenir}\end{minipage}\par}

% Sommaire
\newtcolorbox{sommaire}{enhanced,colback=white,colframe=popink,boxrule=2.2pt,arc=10pt,
  drop shadow=popink,left=18pt,right=18pt,top=13pt,bottom=13pt,before skip=6pt,after skip=20pt,
  before upper={\poppill{Sommaire}{poppurple}{white}\par\vspace{9pt}\popsemi\fontsize{10.6}{14.5}\selectfont\color{popink}}}

% Signature de fin de cours — poussée tout en bas de la dernière page
\newcommand{\findecours}{%
  \par\vfill\nopagebreak
  \begin{center}\popsquiggle{poporange}\end{center}\vspace{4pt}
  \signatureonbuch
}
\AtBeginDocument{\def\llbracket{\mathopen{[\mkern-3.5mu[}}\def\rrbracket{\mathclose{]\mkern-3.5mu]}}} % intervalles d'entiers (glyphe ⟦ absent de la police)
% Glyphes absents de Fira Math : redéfinis à partir de glyphes disponibles
\AtBeginDocument{%
  \def\setminus{\mathbin{\backslash}}\let\smallsetminus\setminus
  \def\vdots{\mathord{\vbox{\baselineskip3.2pt\lineskiplimit0pt\kern4pt\hbox{.}\hbox{.}\hbox{.}}}}%
  \def\ddots{\mathinner{\mkern1mu\raise6.5pt\hbox{.}\mkern2mu\raise3.6pt\hbox{.}\mkern2mu\raise0.7pt\hbox{.}\mkern1mu}}%
  \def\triangle{\mathord{\scalebox{0.9}{$\symup{\Delta}$}}}%
  \def\nabla{\mathord{\raisebox{\depth}{\scalebox{1}[-1]{$\symup{\Delta}$}}}}%
  \def\checkmark{\popcheck{popdark}}%
  \def\star{\mathbin{\ast}}\def\bigstar{\mathord{\ast}}%
  \def\diamond{\mathbin{\scalebox{0.6}{$\Diamond$}}}%
  \def\bigcup{\mathop{\mathchoice{\raisebox{-0.3em}{\scalebox{1.7}{$\cup$}}}{\scalebox{1.25}{$\cup$}}{\cup}{\cup}}\displaylimits}%
  \def\bigcap{\mathop{\mathchoice{\raisebox{-0.3em}{\scalebox{1.7}{$\cap$}}}{\scalebox{1.25}{$\cap$}}{\cap}{\cap}}\displaylimits}%
  \def\lesseqqgtr{\mathrel{\lessgtr}}\def\bigoplus{\mathop{\oplus}\displaylimits}%
}
\providecommand{\ssi}{si et seulement si} % abréviation fréquente
\AtBeginDocument{\providecommand{\wideparen}{\overparen}} % arc de cercle

%% FIGFIX-BEGIN v5 — correctifs globaux des figures (docs/onbuch-generation/tools/figfix)
\usepackage{adjustbox}\usepackage{etoolbox}\tcbuselibrary{raster}
% toute figure TikZ/pgfplots plus large que la ligne est réduite à la largeur disponible
\BeforeBeginEnvironment{tikzpicture}{\begin{adjustbox}{max width=\linewidth}}
\AfterEndEnvironment{tikzpicture}{\end{adjustbox}}
% légendes pgfplots lisibles par-dessus les courbes
\pgfplotsset{every axis/.append style={legend style={fill=white,fill opacity=0.92,text opacity=1,draw=black!15,inner sep=3pt}}}
% étiquettes d'arêtes (syntaxe edge["..."]) avec fond blanc
\tikzset{every edge quotes/.append style={fill=white,inner sep=1.2pt}}
%% FIGFIX-END
\begin{document}

\pagegarde{Leçon 40 — Expression écrite : caractères complexes liés à la rédaction d'un CV}{Chinois}{Tle A}{Module 5 : Mass médias et télécommunication}{ZH40}{2 h}{Cette leçon d'expression écrite apprend aux élèves de Terminale A à rédiger un CV simple en chinois simplifié. Elle insiste sur l'écriture correcte des caractères complexes du thème (traits, radicaux, caractères proches) et sur la traduction thème/version de phrases professionnelles.
\vspace{4pt}
\begin{multicols}{2}\small
\begin{itemize}
\item À la fin de cette leçon, je sais identifier les rubriques d'un CV chinois (个人信息, 教育, 经验, 技能)
\item je sais écrire correctement les caractères complexes du thème en respectant l'ordre des traits et les radicaux
\item je sais distinguer les caractères proches faciles à confondre (简/间, 历/厉, 验/检)
\item je sais utiliser les connecteurs utiles pour structurer un texte de présentation personnelle
\item je sais rédiger un mini-CV simple en chinois à partir d'un modèle commenté
\item je sais traduire du français vers le chinois (thème) des phrases de CV
\end{itemize}\end{multicols}}

\begin{sommaire}
\begin{multicols}{2}
\begin{enumerate}
\item Découverte : le CV chinois et ses rubriques
\item Écriture des caractères complexes (1) : 简历, 经验
\item Écriture des caractères complexes (2) : 教育, 技能, 聘
\item Modèle de CV commenté, structure et connecteurs
\item Traduction : thème et version
\item Grille d'évaluation et production autonome
\item Activité d'intégration
\item Méthodes et exercices résolus
\item Exercices
\item Corrigés des exercices
\item Fiche bilan
\end{enumerate}
\end{multicols}
\end{sommaire}

\begin{objectifs}
\begin{itemize}
\item À la fin de cette leçon, je sais identifier les rubriques d'un CV chinois (个人信息, 教育, 经验, 技能)
\item je sais écrire correctement les caractères complexes du thème en respectant l'ordre des traits et les radicaux
\item je sais distinguer les caractères proches faciles à confondre (简/间, 历/厉, 验/检)
\item je sais utiliser les connecteurs utiles pour structurer un texte de présentation personnelle
\item je sais rédiger un mini-CV simple en chinois à partir d'un modèle commenté
\item je sais traduire du français vers le chinois (thème) des phrases de CV
\item je sais traduire du chinois vers le français (version) un extrait de CV
\item je sais auto-évaluer ma production écrite à l'aide d'une grille de critères
\end{itemize}
\end{objectifs}

\begin{prerequis}
\begin{itemize}
\item Se présenter et présenter sa famille (姓名, 年龄, 职业) — notions de Seconde/Première
\item Les radicaux de base et l'ordre général des traits (de haut en bas, de gauche à droite)
\item Les verbes fréquents : 学习, 工作, 会, 喜欢
\item Les structures de date et de lieu (在 + lieu, année + 年)
\item Le vocabulaire scolaire : 学校, 学生, 老师, 专业 (Première)
\end{itemize}
\end{prerequis}

\newpage

\coursec{Découverte : le CV chinois et ses rubriques}

\courssub{Situation d'accroche et problématique}

À Yaoundé, Marie-Claire, élève de Terminale A, rêve de décrocher une bourse pour étudier à l'Université de Pékin. Le formulaire de candidature exige un CV rédigé en chinois : nom, âge, formation, expérience. Comme beaucoup de jeunes Camerounais postulant à des stages dans les entreprises chinoises installées à Douala ou à Kribi, elle doit maîtriser les caractères complexes du monde professionnel : \zh{简历}, \zh{经验}, \zh{教育}. Devant sa feuille blanche, elle se pose mille questions : par où commencer ? Quelles rubriques un employeur chinois attend-il ? Comment écrire « date de naissance » sans faute ?

\textbf{Problématique :} comment écrire correctement les caractères complexes du monde professionnel et rédiger un CV simple en chinois ? C'est tout l'objet de cette leçon. Dans cette première partie, nous allons d'abord \emph{observer} un vrai modèle de CV chinois, \emph{découvrir} le mot-clé \zh{简历} et ses rubriques, puis \emph{comparer} les habitudes chinoises et camerouno-françaises.

\courssub{Première observation : un CV chinois authentique}

Avant toute explication, lisons le CV que Li Ming, un étudiant camerounais de Yaoundé, a préparé pour postuler à un stage dans une entreprise de télécommunications.

\begin{textebox}[个人简历 — Le CV de Li Ming]
个人简历\\
\emph{Gèrén jiǎnlì.}\\
« Curriculum vitae personnel. »\\[4pt]
姓名：李明。\\
\emph{Xìngmíng : Lǐ Míng.}\\
« Nom et prénom : Li Ming. »\\[4pt]
出生日期：2006年5月10日。\\
\emph{Chūshēng rìqī : 2006 nián 5 yuè 10 rì.}\\
« Date de naissance : 10 mai 2006. »\\[4pt]
地址：雅温得。\\
\emph{Dìzhǐ : Yǎwēndé.}\\
« Adresse : Yaoundé. »\\[4pt]
教育：2020–2024年在雅温得中学学习。\\
\emph{Jiàoyù : 2020–2024 nián zài Yǎwēndé zhōngxué xuéxí.}\\
« Formation : de 2020 à 2024, études au collège-lycée de Yaoundé. »\\[4pt]
技能：我会说中文、法语和英语。\\
\emph{Jìnéng : Wǒ huì shuō Zhōngwén, Fǎyǔ hé Yīngyǔ.}\\
« Compétences : je sais parler chinois, français et anglais. »\\[4pt]
爱好：我喜欢足球和音乐。\\
\emph{Àihào : Wǒ xǐhuan zúqiú hé yīnyuè.}\\
« Loisirs : j'aime le football et la musique. »
\end{textebox}

\textbf{Questions de compréhension (observation guidée) :}
\begin{enumerate}
\item Combien de rubriques compte ce CV ? Relevez-les en caractères.
\item Quelle information personnelle apparaît en deuxième position, juste après le nom ?
\item Comment Li Ming exprime-t-il « je sais parler chinois » ? Soulignez le verbe modal utilisé.
\item Dans quelle langue la date est-elle écrite ? Observez l'ordre : année, mois, jour.
\end{enumerate}

\textbf{Glossaire du texte :}

\begin{center}
\begin{tabularx}{\linewidth}{|X|X|X|X|}
\hline
\rowcolor{popblueL} Caractères & Pinyin & Nature & Traduction \\
\hline
简历 & jiǎnlì & nom & curriculum vitae (CV) \\
\hline
姓名 & xìngmíng & nom & nom et prénom \\
\hline
出生日期 & chūshēng rìqī & expression nominale & date de naissance \\
\hline
地址 & dìzhǐ & nom & adresse \\
\hline
教育 & jiàoyù & nom & éducation, formation \\
\hline
中学 & zhōngxué & nom & école secondaire, collège-lycée \\
\hline
技能 & jìnéng & nom & compétence, savoir-faire \\
\hline
爱好 & àihào & nom & loisir, passe-temps \\
\hline
会 & huì & verbe modal & savoir (faire), être capable de \\
\hline
\end{tabularx}
\end{center}

\courssub{Le mot-clé : 简历 (jiǎnlì), le curriculum vitae}

Observons le mot \zh{简历}. Il est composé de deux caractères que nous pouvons décomposer pour mieux les mémoriser :

\begin{itemize}
\item \zh{简} (jiǎn) = « simple, simplifié ». On le retrouve dans \zh{简单} (jiǎndān, « simple ») et dans \zh{简体字} (jiǎntǐzì, « caractères simplifiés »). Sa clé est le \cle{radical du bambou} \zh{⺮} (zhúzìtóu) en haut : autrefois, on écrivait sur des lamelles de bambou, et « simplifier » l'écriture était lié à ce support.
\item \zh{历} (lì) = « parcours, expérience, histoire ». On le retrouve dans \zh{历史} (lìshǐ, « histoire ») et \zh{经历} (jīnglì, « expérience vécue »). Sa clé en chinois simplifié est \zh{厂} (hǎn, « falaise »).
\end{itemize}

\begin{definition}[简历 (jiǎnlì)]
Le mot \zh{简历} (jiǎnlì), littéralement « parcours simplifié », désigne le \cle{curriculum vitae} : un document écrit qui résume le parcours d'une personne (identité, formation, expérience, compétences) présenté à un employeur ou à une institution. On dit aussi \zh{个人简历} (gèrén jiǎnlì, « CV personnel ») pour le titre complet du document.
\end{definition}

\begin{exemplebox}[Le mot 简历 en contexte]
\zh{这是我的简历。}\\
\emph{Zhè shì wǒ de jiǎnlì.}\\
« Voici mon CV. »\\[4pt]
\zh{请把你的简历发给我们。}\\
\emph{Qǐng bǎ nǐ de jiǎnlì fā gěi wǒmen.}\\
« Veuillez nous envoyer votre CV. »\\[4pt]
\zh{找工作需要一份简历。}\\
\emph{Zhǎo gōngzuò xūyào yí fèn jiǎnlì.}\\
« Pour chercher un emploi, il faut un CV. » (Notez le classificateur \zh{份} fèn, utilisé pour les documents.)
\end{exemplebox}

\courssub{Les cinq rubriques d'un CV chinois}

En observant le CV de Li Ming, nous avons relevé ses rubriques. Un CV chinois simple, au niveau HSK 3-4, comporte typiquement \cle{cinq rubriques} :

\begin{center}
\begin{tabularx}{\linewidth}{|X|X|X|X|}
\hline
\rowcolor{popblueL} Caractères & Pinyin & Nature & Traduction \\
\hline
个人信息 & gèrén xìnxī & expression nominale & informations personnelles \\
\hline
教育 & jiàoyù & nom & formation, éducation \\
\hline
工作经验 & gōngzuò jīngyàn & expression nominale & expérience professionnelle \\
\hline
技能 & jìnéng & nom & compétences \\
\hline
爱好 & àihào & nom & loisirs, centres d'intérêt \\
\hline
\end{tabularx}
\end{center}

Décomposons les expressions composées, car chaque syllabe a un sens :

\begin{itemize}
\item \zh{个人信息} (gèrén xìnxī) : \zh{个人} = « individuel, personnel » (\zh{个} + \zh{人}, « personne ») ; \zh{信息} = « information ». C'est la rubrique de l'identité : nom, date de naissance, adresse, téléphone.
\item \zh{工作经验} (gōngzuò jīngyàn) : \zh{工作} = « travail » ; \zh{经验} = « expérience ». Le caractère \zh{经} porte la clé de la soie \zh{纟} ; \zh{验} porte la clé du cheval \zh{马} (autrefois, « éprouver, vérifier » un cheval). Nous étudierons l'écriture de ces caractères complexes dans la section suivante.
\item \zh{技能} (jìnéng) : \zh{技} = « technique » (clé de la main \zh{扌}) ; \zh{能} = « capacité » (clé \zh{月}).
\item \zh{爱好} (àihào) : \zh{爱} = « aimer » ; \zh{好} se lit ici au quatrième ton \emph{hào} = « goût, penchant » (et non \emph{hǎo}, « bon »).
\end{itemize}

\begin{attention}[Piège de prononciation : 好 dans 爱好]
Dans \zh{爱好}, le caractère \zh{好} se prononce \emph{hào} (4\textsuperscript{e} ton) et signifie « aimer, avoir le goût de ». Ne dites pas *\emph{àihǎo}. C'est un cas classique de caractère à plusieurs lectures : \zh{好} \emph{hǎo} = « bon » ; \zh{好} \emph{hào} = « aimer » (comme dans \zh{好奇} hàoqí, « curieux »).
\end{attention}

\begin{popfigure}
\begin{tikzpicture}[mindmap, grow cyclic, every node/.style={concept}, concept color=popblue!60,
root concept/.append style={font=\Large, concept color=popblue},
level 1/.append style={level distance=3.2cm, sibling angle=72},
level 1 concept/.append style={font=\small}]
\node[root concept, align=center]{简历\\ jiǎnlì\\ le CV}
  child[concept color=poporange!70] { node[align=center]{个人信息\\ gèrén xìnxī\\ identité} }
  child[concept color=popgreen!70] { node[align=center]{教育\\ jiàoyù\\ formation} }
  child[concept color=poppurple!60] { node[align=center]{经验\\ jīngyàn\\ expérience} }
  child[concept color=poppink!70] { node[align=center]{技能\\ jìnéng\\ compétences} }
  child[concept color=popgold!80] { node[align=center]{爱好\\ àihào\\ loisirs} };
\end{tikzpicture}
\legende{Carte mentale : les cinq rubriques du CV chinois 简历.}
\end{popfigure}

\courssub{CV chinois et CV franco-camerounais : quelles différences ?}

Le CV n'est pas rédigé de la même manière partout. Voici les différences essentielles à connaître pour ne pas commettre d'impair culturel.

\begin{center}
\begin{tabularx}{\linewidth}{|X|X|X|}
\hline
\rowcolor{popblueL} Élément & CV français / camerounais & CV chinois \\
\hline
Photo & facultative, parfois déconseillée & souvent jointe, très courante \\
\hline
Date de naissance & souvent omise (lutte contre les discriminations) & régulièrement indiquée : \zh{出生日期} \\
\hline
Ordre des expériences & chronologique inverse (la plus récente d'abord) & chronologique inverse également \\
\hline
Longueur & une page recommandée & une page, concision très valorisée \\
\hline
État civil & parfois mentionné au Cameroun & peut apparaître dans \zh{个人信息} \\
\hline
\end{tabularx}
\end{center}

\begin{popfigure}
\begin{tikzpicture}[
box/.style={draw=popline, rounded corners=3pt, fill=#1, text width=4.6cm, align=left, inner sep=6pt, font=\small},
title/.style={font=\bfseries\small, align=center}]
\node[title] (t1) at (0,0) {CV français / camerounais};
\node[title] (t2) at (6.2,0) {CV chinois 简历};
\node[box=poporangeL, anchor=north] (a1) at (0,-0.5) {Photo facultative};
\node[box=popblueL, anchor=north] (b1) at (6.2,-0.5) {Photo souvent jointe};
\node[box=poporangeL, anchor=north] (a2) at ($(a1.south)+(0,-0.25)$) {Date de naissance souvent omise};
\node[box=popblueL, anchor=north] (b2) at ($(b1.south)+(0,-0.25)$) {出生日期 indiquée};
\node[box=poporangeL, anchor=north] (a3) at ($(a2.south)+(0,-0.25)$) {Ordre chronologique inverse};
\node[box=popblueL, anchor=north] (b3) at ($(b2.south)+(0,-0.25)$) {Ordre chronologique inverse aussi};
\node[box=poporangeL, anchor=north] (a4) at ($(a3.south)+(0,-0.25)$) {Rubriques : formation, expérience, compétences};
\node[box=popblueL, anchor=north] (b4) at ($(b3.south)+(0,-0.25)$) {个人信息 / 教育 / 经验 / 技能 / 爱好};
\draw[-{Stealth[length=2.5mm]}, thick, popmuted] (a2.east) -- (b2.west);
\draw[-{Stealth[length=2.5mm]}, thick, popmuted] (a4.east) -- (b4.west);
\end{tikzpicture}
\legende{Schéma comparatif : le CV franco-camerounais face au CV chinois.}
\end{popfigure}

\begin{savaistu}[Une seule page, et pas une de plus !]
En Chine, le CV tient presque toujours sur \cle{une seule page}. La concision est très valorisée : un recruteur chinois reçoit parfois des centaines de candidatures pour un seul poste, et un document long est mal perçu. Le mot \zh{简历} lui-même porte cette idée : un parcours « simplifié », réduit à l'essentiel. Au Cameroun comme en Chine, soigner la présentation (papier propre, rubriques claires, aucune rature) est une marque de respect envers le lecteur.
\end{savaistu}

\courssub{Les formules de base pour se présenter}

Pour remplir les rubriques \zh{个人信息} et \zh{技能}, trois structures fondamentales suffisent au niveau HSK 2-3. Observons-les d'abord dans le CV de Li Ming, puis énonçons la règle.

\begin{propriete}[Structure 1 : 我叫… — se nommer]
\textbf{Schéma :} 我 + 叫 + nom.\\
Le verbe \zh{叫} (jiào, « s'appeler ») introduit le nom ou le prénom. Pas de verbe « être » ici.\\
\zh{我叫李明。} \emph{Wǒ jiào Lǐ Míng.} « Je m'appelle Li Ming. »\\
Variante plus formelle pour un CV : \zh{姓名：李明。} \emph{Xìngmíng : Lǐ Míng.} « Nom : Li Ming. »
\end{propriete}

\begin{propriete}[Structure 2 : 我出生于… — donner sa naissance]
\textbf{Schéma :} 我 + 出生于 + date / lieu.\\
\zh{出生} (chūshēng) = « naître » ; \zh{于} (yú) = préposition formelle « en, à ». La date se donne dans l'ordre \cle{du plus grand au plus petit} : année + 年, mois + 月, jour + 日.\\
\zh{我出生于2006年5月10日。} \emph{Wǒ chūshēng yú 2006 nián 5 yuè 10 rì.} « Je suis né(e) le 10 mai 2006. »\\
\zh{我出生于雅温得。} \emph{Wǒ chūshēng yú Yǎwēndé.} « Je suis né(e) à Yaoundé. »
\end{propriete}

\begin{propriete}[Structure 3 : 我会说… — déclarer une compétence linguistique]
\textbf{Schéma :} 我 + 会 + 说 + langue (+ 和 + langue).\\
Le modal \zh{会} (huì) exprime la \cle{compétence acquise} (« savoir faire »), et non le futur. La coordination se fait avec \zh{和} (hé, « et ») ou l'énumération avec la virgule chinoise \zh{、}.\\
\zh{我会说中文。} \emph{Wǒ huì shuō Zhōngwén.} « Je sais parler chinois. »\\
\zh{我会说中文和法语。} \emph{Wǒ huì shuō Zhōngwén hé Fǎyǔ.} « Je sais parler chinois et français. »\\
\zh{我会说中文、法语和英语。} \emph{Wǒ huì shuō Zhōngwén, Fǎyǔ hé Yīngyǔ.} « Je sais parler chinois, français et anglais. »
\end{propriete}

\begin{attention}[Ne traduisez pas « date de naissance » par 生日 dans un CV formel]
Les francophones traduisent souvent « date de naissance » par \zh{生日} (shēngrì). Or \zh{生日} signifie « anniversaire » (la fête, le jour qu'on célèbre chaque année) : \zh{祝你生日快乐！} « Joyeux anniversaire ! ». Dans un document administratif ou un CV, on utilise l'expression formelle \zh{出生日期} (chūshēng rìqī), littéralement « date de naissance ». Écrire \zh{生日：2006年5月10日} dans un CV ferait très informel, presque enfantin. Retenez : \cle{生日 = anniversaire (fête) ; 出生日期 = date de naissance (administratif)}.
\end{attention}

\begin{attention}[L'ordre de la date : du plus grand au plus petit]
En français, on dit « le 10 mai 2006 » (jour → mois → année). En chinois, c'est l'inverse : \zh{2006年5月10日} (année → mois → jour). C'est une règle générale de la langue chinoise : \emph{du plus grand au plus petit}, aussi pour les adresses (\zh{喀麦隆雅温得} « Cameroun, Yaoundé » : pays puis ville). Ne dites jamais *\zh{10日5月2006年}.
\end{attention}

\begin{exoresolu}[Compléter les rubriques d'un mini-CV]
Énoncé : Aminatou, élève de Terminale à Douala, rédige son CV en chinois. Complétez les phrases avec les mots suivants : \zh{出生日期}, \zh{叫}, \zh{会}, \zh{爱好}.\\
1. 我\_\_\_\_Aminatou。\\
2. \_\_\_\_：2007年3月15日。\\
3. 我\_\_\_\_说法语和英语。\\
4. \_\_\_\_：我喜欢音乐和篮球。
\tcblower
Solution détaillée :\\
1. \zh{我叫Aminatou。} — Le verbe « s'appeler » est \zh{叫} ; aucun verbe « être » n'est nécessaire.\\
2. \zh{出生日期：2007年3月15日。} — Dans une rubrique formelle de CV, on utilise \zh{出生日期} et non \zh{生日} (anniversaire). L'ordre année → mois → jour est respecté.\\
3. \zh{我会说法语和英语。} — Le modal \zh{会} exprime la compétence acquise « savoir parler » ; il se place avant le verbe \zh{说}.\\
4. \zh{爱好：我喜欢音乐和篮球。} — La rubrique des loisirs est \zh{爱好} (attention : \zh{好} se lit \emph{hào}, 4\textsuperscript{e} ton).
\end{exoresolu}

\begin{aretenir}[L'essentiel de la section]
\begin{itemize}
\item \zh{简历} (jiǎnlì) = le CV : \zh{简} « simple » + \zh{历} « parcours ».
\item Cinq rubriques : \zh{个人信息} (identité), \zh{教育} (formation), \zh{工作经验} (expérience), \zh{技能} (compétences), \zh{爱好} (loisirs).
\item En Chine : photo courante, \zh{出生日期} indiquée, une seule page, ordre chronologique inverse.
\item Trois formules clés : \zh{我叫…} (je m'appelle), \zh{我出生于…} (je suis né(e) le/à), \zh{我会说…} (je sais parler).
\item Pièges : \zh{生日} ≠ \zh{出生日期} ; date dans l'ordre année → mois → jour ; \zh{爱好} se lit \emph{àihào}.
\end{itemize}
\end{aretenir}

\coursec{Écriture des caractères complexes (1) : 简历、经验}

\courssub{Analyse graphique et ordre des traits}

Maîtriser l'écriture des caractères du thème « CV » demande de la rigueur. Nous appliquons les règles fondamentales : \emph{haut → bas}, \emph{gauche → droite}, \emph{horizontal avant vertical}, \emph{traits centraux avant les ailes}, \emph{fermeture en dernier}.

\begin{center}
\begin{tabularx}{\linewidth}{|X|X|X|X|X|}
\hline
\rowcolor{popblueL} Caractère & Pinyin & Radical (Clé) & Nombre de traits & Ordre des traits (description) \\
\hline
简 & jiǎn & \zh{⺮} (bambou, zhú) & 13 & 1. Haut : \zh{⺮} (6 traits : points, horizontales, crochet). 2. Bas : \zh{间} (porte \zh{门} + \zh{日}) : \zh{门} (3 traits) puis \zh{日} (4 traits). \\
\hline
历 & lì & \zh{厂} (falaise, hǎn) & 5 & 1. \zh{厂} (2 traits : horizontale pliée, verticale). 2. Intérieur : \zh{彐} (jǐ, 3 traits : horizontale, courbe, horizontale). \\
\hline
经 & jīng & \zh{纟} (soie, sī) & 8 & 1. Clé \zh{纟} (3 traits : diagonales, point). 2. Phonétique \zh{巠} (jīng) : \zh{工} (3 traits) + \zh{彐} (3 traits) simplifié ici en 5 traits au total pour la partie droite. \\
\hline
验 & yàn & \zh{马} (cheval, mǎ) & 10 & 1. Haut : \zh{ㄉ} (2 traits). 2. \zh{马} (3 traits : horizontale, verticale pliée, point). 3. Bas : \zh{ㄗ} (2 traits) — forme simplifiée de \zh{騐}. \\
\hline
\end{tabularx}
\end{center}

\courssub{Caractères proches : ne pas confondre}

\begin{center}
\begin{tabularx}{\linewidth}{|X|X|X|}
\hline
\rowcolor{popblueL} Caractère cible & Caractère proche & Différence clé \\
\hline
\zh{简} (jiǎn, simple) & \zh{检} (jiǎn, inspecter) & \zh{简} a \zh{⺮} (bambou) en haut ; \zh{检} a \zh{木} (bois) à gauche. \\
\hline
\zh{历} (lì, parcours) & \zh{厉} (lì, sévère) & \zh{历} a \zh{彐} inside \zh{厂} ; \zh{厉} a \zh{历} (traditionnel) ou \zh{厂} + \zh{丽} (simplifié : \zh{厉} = \zh{厂} + \zh{丽}). \\
\hline
\zh{经} (jīng, trame/expérience) & \zh{径} (jìng, diamètre/chemin) & \zh{经} : clé \zh{纟} + \zh{巠} ; \zh{径} : clé \zh{纟} + \zh{巠} (même phonétique) mais sens « chemin direct ». \zh{经} est pour « expérience » (\zh{经验}). \\
\hline
\zh{验} (yàn, vérifier) & \zh{马} (mǎ, cheval) & \zh{验} contient \zh{马} en bas à droite, mais avec un toit \zh{ㄉ} et une base \zh{ㄗ}. Ne pas écrire \zh{马} seul pour \zh{验}. \\
\hline
\end{tabularx}
\end{center}

\begin{methode}[Mémoriser 经验 (jīngyàn, expérience)]
\begin{enumerate}
\item \textbf{Décomposer} : \zh{经} (fil de la trame → parcours constant) + \zh{验} (éprouver un cheval → vérifier).
\item \textbf{Visualiser} : Le fil (\zh{纟}) qui traverse l'épreuve du cheval (\zh{马}) = l'expérience forgée par la pratique.
\item \textbf{Écrire} : 3 fois en respectant l'ordre des traits, en disant le nom des traits à voix haute (ex: « horizontale, verticale, pliée… »).
\item \textbf{Contextualiser} : Écrire \zh{工作经验} (gōngzuò jīngyàn) et \zh{经历} (jīnglì).
\end{enumerate}
\end{methode}

\begin{exercice}[Entraînement à l'écriture (1)]
Écrivez chaque caractère 5 fois dans votre cahier d'écriture (纸格子 zhǐ gézi), en respectant l'ordre des traits et les proportions. Puis écrivez les mots composés : \zh{简历}, \zh{经验}, \zh{工作经验}, \zh{经历}.
\end{exercice}

\coursec{Écriture des caractères complexes (2) : 教育、技能、聘}

\courssub{Analyse graphique et ordre des traits}

\begin{center}
\begin{tabularx}{\linewidth}{|X|X|X|X|X|}
\hline
\rowcolor{popblueL} Caractère & Pinyin & Radical (Clé) & Nombre de traits & Ordre des traits (description) \\
\hline
教 & jiào & \zh{攵} (pū, action de taper) & 11 & 1. Haut : \zh{孝} (xiào) : \zh{老} (6 traits) sans le trait final + \zh{子} (3 traits). 2. Bas : \zh{攵} (4 traits : horizontale, verticale, diagonale gauche, point droit). \\
\hline
育 & yù & \zh{月} (viande/chair, ròu) & 8 & 1. Haut : \zh{厶} (sī, 2 traits). 2. Milieu : \zh{月} (4 traits : deux verticales, deux horizontales). 3. Bas : \zh{ㄅ} (2 traits) — forme simplifiée de \zh{月} inversé ? Non, \zh{育} = \zh{厶} + \zh{月} + \zh{ㄅ} (variante). Ordre standard : \zh{厶} (2), \zh{月} (4), \zh{ㄅ} (2) = 8 traits. \\
\hline
技 & jì & \zh{扌} (main, shǒu) & 7 & 1. Clé \zh{扌} (3 traits : horizontale, crochet vertical, point). 2. Phonétique \zh{枝} (zhī) simplifiée : \zh{木} (4 traits) + \zh{支} (partie droite) → ici \zh{技} = \zh{扌} + \zh{枝} (simplifié en 4 traits pour la droite). Total 7. \\
\hline
能 & néng & \zh{月} (viande/chair, ròu) & 10 & 1. Haut : \zh{厶} (2 traits). 2. \zh{匕} (bǐ, 2 traits). 3. \zh{月} (4 traits). 4. Bas : \zh{匕} (2 traits) — structure complexe « double lune ». \\
\hline
聘 & pìn & \zh{讠} (parole, yán) & 13 & 1. Clé \zh{讠} (2 traits : point, courbe). 2. Phonétique \zh{贲} (bēn/bì) : \zh{贝} (4 traits) + \zh{贲} (haut : \zh{ㄉ} + \zh{田} + \zh{ㄗ}) — simplifié en \zh{贲} (11 traits). Total 13. \\
\hline
\end{tabularx}
\end{center}

\courssub{Caractères proches : vigilance requise}

\begin{center}
\begin{tabularx}{\linewidth}{|X|X|X|}
\hline
\rowcolor{popblueL} Caractère cible & Caractère proche & Différence clé \\
\hline
\zh{教} (jiào, enseigner) & \zh{孝} (xiào, piété filiale) & \zh{教} = \zh{孝} + \zh{攵} (action). \zh{孝} seul = piété. \\
\hline
\zh{育} (yù, élever/éduquer) & \zh{毓} (yù, naître/élever - littér.) & \zh{育} est le forme courante simplifiée ; \zh{毓} est plus littéraire, même sens, structure différente (haut \zh{毋}). \\
\hline
\zh{技} (jì, technique) & \zh{枝} (zhī, branche) & \zh{技} a \zh{扌} (main) → action manuelle ; \zh{枝} a \zh{木} (bois) → partie d'arbre. \\
\hline
\zh{能} (néng, capacité) & \zh{态} (tài, attitude) & \zh{能} : double \zh{月} + \zh{匕} ; \zh{态} : \zh{心} (cœur) + \zh{能} (phonétique) → état d'esprit. \\
\hline
\zh{聘} (pìn, engager) & \zh{请} (qǐng, inviter/prier) & \zh{聘} : \zh{讠} + \zh{贲} (embauche formelle) ; \zh{请} : \zh{讠} + \zh{青} (demander poliment). \\
\hline
\end{tabularx}
\end{center}

\begin{attention}[Piège d'écriture : 能 (néng) et 态 (tài)]
\zh{能} s'écrit avec \emph{deux} éléments \zh{月} (lunes/viande) superposés (un en haut à gauche via \zh{厶}+\zh{匕}, un au milieu) et un \zh{匕} en bas. \zh{态} (attitude) s'écrit \zh{心} (cœur, clé de fond) + \zh{能} (phonétique). Ne pas oublier le point final du \zh{心} (le point bas du \zh{心} devient une horizontale courte en position de fond).
\end{attention}

\begin{methode}[Mémoriser 教育 (jiàoyù) et 技能 (jìnéng)]
\begin{enumerate}
\item \textbf{教育} : \zh{教} (action \zh{攵} de \zh{孝} transmettre) + \zh{育} (faire grandir \zh{月} le corps). « Enseigner et élever ».
\item \textbf{技能} : \zh{技} (main \zh{扌} habile \zh{枝}) + \zh{能} (capacité innée/acquise). « Savoir-faire technique ».
\item \textbf{聘用} (pìnyòng, embaucher) : \zh{聘} (parole \zh{讠} formelle \zh{贲}) + \zh{用} (utiliser). Verbe clé pour la rubrique \zh{工作经验}.
\end{enumerate}
\end{methode}

\begin{exercice}[Entraînement à l'écriture (2)]
Écrivez chaque caractère 5 fois : \zh{教}, \zh{育}, \zh{技}, \zh{能}, \zh{聘}. Puis les mots : \zh{教育}, \zh{技能}, \zh{聘用}, \zh{应聘} (yìngpìn, postuler).
\end{exercice}

\coursec{Modèle de CV commenté, structure et connecteurs}

\courssub{Texte modèle complet : CV de stage (niveau HSK 3-4)}

Voici un CV plus complet, structuré avec des connecteurs logiques, adapté à une candidature pour un stage en entreprise.

\begin{textebox}[个人简历 — Modèle complet pour stage]
个人简历\\
\emph{Gèrén jiǎnlì.}\\[4pt]
\textbf{个人信息}\\
姓名：王伟 (Wáng Wěi)。\\
出生日期：2005年10月20日。\\
联系电话：+237 6XX XX XX XX。\\
电子邮箱：wangwei@email.com。\\
地址：喀麦隆杜阿拉。\\[4pt]
\textbf{教育背景} (Jiàoyù bèijǐng)\\
2021年9月 – 2024年6月 \quad 雅温得第一中学 (Yǎwēndé Dìyī Zhōngxué)\\
\emph{期间} (qíjiān), \emph{主修} (zhǔxiū) 科学系列, \emph{获得} (huòdé) 优异成绩。\\[4pt]
\textbf{工作经验}\\
2023年7月 – 2023年8月 \quad 华信科技有限公司 (Huáxìn Kējì Yǒuxiàn Gōngsī, Douala)\\
\emph{实习生} (shíxīshēng), \emph{负责} (fùzé) 数据录入 \emph{和} 客户接待。\\
\emph{协助} (xiēzhù) 经理 \emph{完成} (wánchéng) 月度报告。\\[4pt]
\textbf{技能特长}\\
语言: 中文 (HSK 4), 法语 (langue maternelle), 英语 (B2)。\\
计算机: 办公软件 (Word, Excel, PPT), 基础编程 (Python)。\\
\emph{掌握} (zhǎngwò) 汉字输入法 (拼音/五笔)。\\[4pt]
\textbf{爱好与自我评价}\\
爱好: 篮球、阅读、中非文化交流。\\
自我评价: 我\emph{踏实肯干} (tàshí kěngàn), \emph{有团队精神} (yǒu tuánduì jīngshén), \emph{渴望} (kěwàng) 学习新知识。
\end{textebox}

\courssub{Connecteurs et mots-outils du CV professionnel}

Pour structurer un CV, on utilise des connecteurs temporels et logiques précis. Voici les plus utiles au niveau Terminale :

\begin{center}
\begin{tabularx}{\linewidth}{|X|X|X|}
\hline
\rowcolor{popblueL} Connecteur (Caractères) & Pinyin & Emploi / Traduction \\
\hline
期间 & qíjiān & « pendant ce temps, durant cette période » (suivi d'une action). Ex: 2021–2024 \zh{期间}… \\
\hline
随后 & suíhòu & « ensuite, par la suite » (chronologie). Ex: \zh{毕业后，随后找工作。} \\
\hline
此外 & cǐwài & « de plus, par ailleurs » (ajout d'information). Ex: \zh{此外，我会英语。} \\
\hline
负责 & fùzé & « être responsable de, gérer » (verbe d'action pour expérience). Ex: \zh{负责销售。} \\
\hline
协助 & xiēzhù & « assister, aider » (verbe d'action). Ex: \zh{协助经理。} \\
\hline
参与 & cānyù & « participer à ». Ex: \zh{参与项目开发。} \\
\hline
掌握 & zhǎngwò & « maîtriser (une compétence) ». Ex: \zh{掌握Python。} \\
\hline
熟练 & shúliàn & « maîtriser couramment, être rompu à » (adjectif). Ex: \zh{熟练使用Office。} \\
\hline
\end{tabularx}
\end{center}

\begin{propriete}[Structure de la rubrique 工作经验]
\textbf{Schéma :} Période + Lieu/Entreprise + Poste + \zh{负责/协助/参与} + Mission.\\
L'ordre est \emph{antéchronologique} (le plus récent en premier). Les phrases sont nominales ou verbales sans sujet (le sujet « je » est sous-entendu).
\end{propriete}

\begin{exemplebox}[Exemple rubrique Expérience]
2023年6月–8月 \quad 杜阿拉港务局\\
实习生 \zh{负责} 集装箱数据录入，\zh{协助} 翻译中法文件。
\end{exemplebox}


\begin{attention}[Le piège du « Je » (我) dans le CV chinois]
Dans un CV formel chinois, on \emph{omet le sujet 我} dans les rubriques \zh{工作经验} et \zh{技能特长}. On écrit directement le verbe d'action : \zh{负责…}, \zh{掌握…}, \zh{熟练使用…}. Écrire \zh{我负责…} \zh{我掌握…} alourdit le texte et fait « élève » plutôt que « professionnel ».
\end{attention}

\coursec{Traduction : thème et version}

\courssub{Thème (Français → Chinois)}

Traduisez les phrases suivantes en chinois (caractères + pinyin). Utilisez le vocabulaire et les structures de la leçon.

\begin{exercice}[Thème n°1 : Informations personnelles]
\begin{enumerate}
\item Je m'appelle Paul Biya, je suis né le 13 février 1933 à Mvomeka'a.
\item Mon adresse est : B.P. 123, Yaoundé, Cameroun.
\item Je sais parler chinois, français et anglais.
\item Mes loisirs sont la lecture et le football.
\end{enumerate}
\end{exercice}

\begin{corrige}[Thème n°1]
\begin{enumerate}
\item \zh{我叫保罗·比亚。我出生于1933年2月13日，出生地是姆沃梅卡。} (Wǒ jiào Bǎoluó·Bǐyà. Wǒ chūshēng yú 1933 nián 2 yuè 13 rì, chūshēng dì shì Mǔwòméikǎ.)
\item \zh{地址：喀麦隆雅温得邮政信箱123号。} (Dìzhǐ: Kāmàilóng Yǎwēndé yóuzhèng xìnxiāng 123 hào.)
\item \zh{我会说中文、法语和英语。} (Wǒ huì shuō Zhōngwén, Fǎyǔ hé Yīngyǔ.)
\item \zh{爱好：阅读和足球。} (Àihào: yuèdú hé zúqiú.)
\end{enumerate}
\end{corrige}

\begin{exercice}[Thème n°2 : Formation et Expérience]
\begin{enumerate}
\item De 2020 à 2023, j'ai étudié au Lycée de Douala.
\item En 2023, j'ai fait un stage à l'entreprise Huawei à Douala. J'étais responsable de la saisie de données.
\item J'ai participé à un projet de traduction.
\item Je maîtrise les logiciels de bureautique (Word, Excel).
\end{enumerate}
\end{exercice}

\begin{corrige}[Thème n°2]
\begin{enumerate}
\item \zh{2020年至2023年在杜阿拉中学学习。} (2020 nián zhì 2023 nián zài Dūālā zhōngxué xuéxí.)
\item \zh{2023年在杜阿拉华为公司实习，负责数据录入。} (2023 nián zài Dūālā Huáwéi gōngsī shíxí, fùzé shùjù lùrù.)
\item \zh{参与了翻译项目。} (Cānyù le fānyì xiàngmù.) \emph{Note :} \zh{了} marque l'accomplissement de l'action passée.
\item \zh{熟练使用办公软件 (Word, Excel)。} (Shúliàn shǐyòng bàngōng ruǎnjiàn (Word, Excel).)
\end{enumerate}
\end{corrige}

\courssub{Version (Chinois → Français)}

Traduisez les extraits de CV suivants en français naturel.

\begin{exercice}[Version n°1]
\begin{enumerate}
\item 个人信息：姓名：玛丽·克莱尔。出生日期：2006年8月15日。联系电话：+237 6XX XX XX XX。
\item 教育背景：2021年–2024年 雅温得女子中学。期间主修文学系列。
\item 工作经验：2023年7月 杜阿拉中非友谊医院 实习生 协助护士整理病历。
\item 技能特长：语言：中文HSK3，法语母语。计算机：熟练使用WPS Office。
\item 爱好：喜欢唱歌、跳舞，热衷中非文化交流活动。
\end{enumerate}
\end{exercice}

\begin{corrige}[Version n°1]
\begin{enumerate}
\item Informations personnelles : Nom : Marie-Claire. Date de naissance : 15 août 2006. Téléphone : +237 6XX XX XX XX.
\item Formation : 2021–2024 Lycée des jeunes filles de Yaoundé. Durant cette période, spécialisation en série littéraire.
\item Expérience professionnelle : Juillet 2023, Hôpital de l'amitié sino-camerounaise de Douala. Stagiaire. A assisté les infirmières pour classer les dossiers médicaux.
\item Compétences : Langues : Chinois HSK3, français langue maternelle. Informatique : Maîtrise courante de WPS Office.
\item Loisirs : Aime le chant, la danse, passionnée par les activités d'échange culturel sino-camerounais.
\end{enumerate}
\end{corrige}

\coursec{Grille d'évaluation et production autonome}

\courssub{Grille d'évaluation du CV (sur 20 points)}

Utilisez cette grille pour vous auto-évaluer ou évaluer un camarade.

\begin{center}
\begin{tabularx}{\linewidth}{|X|X|X|}
\hline
\rowcolor{popblueL} Critère & Indicateurs de réussite & Points \\
\hline
\textbf{Caractères \& Écriture} (5 pts) & Caractères cibles (\zh{简历经验教育技能聘} etc.) corrects, ordre des traits respecté, proportions équilibrées, pas de confusion avec caractères proches. & /5 \\
\hline
\textbf{Vocabulaire} (4 pts) & Rubriques complètes (\zh{个人信息}, \zh{教育}, \zh{工作经验}, \zh{技能}, \zh{爱好}). Vocabulaire précis (\zh{出生日期} pas \zh{生日}, \zh{熟练}, \zh{负责}). Classificateurs corrects (\zh{份}简历). & /4 \\
\hline
\textbf{Grammaire \& Structures} (5 pts) & Structures de base correctes (\zh{我叫}, \zh{我出生于}, \zh{我会说}). Omission du sujet \zh{我} dans \zh{工作经验}. Connecteurs utilisés (\zh{期间}, \zh{负责}, \zh{掌握}). Ordre date/adresse respecté. & /5 \\
\hline
\textbf{Structure \& Connecteurs} (3 pts) & Rubriques clairement séparées, ordre antéchronologique respecté, mise en page aérée (une page), connecteurs logiques fluides. & /3 \\
\hline
\textbf{Contenu \& Socioculturel} (3 pts) & Infos pertinentes et véridiques. Photo (optionnelle mais mentionnée). Adaptation aux attentes chinoises (concision, formalisme). Pas de fautes de registre (ex: \zh{生日} dans CV). & /3 \\
\hline
\textbf{TOTAL} & & \textbf{/20} \\
\hline
\end{tabularx}
\end{center}

\courssub{Production autonome : Rédigez votre CV}

\begin{experience}[Atelier : Mon CV en chinois (Durée : 30 min)]
\textbf{Consigne :} Rédigez votre propre \zh{个人简历} sur une feuille A4 (ou format numérique) en vous inspirant du modèle de Wang Wei.
\begin{enumerate}
\item \textbf{Préparation (5 min) :} Remplissez un brouillon en français : Nom, Date/lieu naissance, Adresse/Tél/Email, Formation (depuis la 3ème), Expériences (stages, jobs d'été, bénévolat), Compétences (Langues + Informatique + Autres), Loisirs.
\item \textbf{Rédaction (15 min) :} Transposez en chinois. \emph{Attention :} Utilisez \zh{出生日期}, omettez \zh{我} dans l'expérience, utilisez \zh{负责/协助/掌握/熟练}. Vérifiez l'ordre date (année-mois-jour) et adresse (pays-ville).
\item \textbf{Relecture croisée (10 min) :} Échangez avec un binôme. Vérifiez avec la grille d'évaluation : caractères cibles, structures, connecteurs, mise en page. Surlignez les erreurs suspectes au stylo rouge.
\item \textbf{Remise :} Version finale propre pour le portfolio de la classe.
\end{enumerate}
\end{experience}

\begin{aretenir}[Synthèse de la leçon ZH40 — Le CV chinois]
\begin{itemize}
\item \textbf{Vocabulaire clé :} \zh{简历} (CV), \zh{个人信息}, \zh{教育背景}, \zh{工作经验}, \zh{技能特长}, \zh{爱好}, \zh{出生日期}, \zh{实习生}, \zh{负责}, \zh{协助}, \zh{掌握}, \zh{熟练}, \zh{聘用}.
\item \textbf{Caractères complexes à maîtriser :} \zh{简历}, \zh{经验}, \zh{教育}, \zh{技能}, \zh{聘} (traits, radicaux, confusions).
\item \textbf{Structures types :} \zh{我叫…} / \zh{姓名：…} ; \zh{我出生于…年…月…日} ; \zh{我会说…} ; \zh{期间 + action} ; \zh{负责/协助/参与 + mission}.
\item \textbf{Règles d'or :} Une page max. Photo acceptée. \zh{出生日期} (pas \zh{生日}). Date : Année-Mois-Jour. Adresse : Pays-Ville. Pas de \zh{我} dans l'expérience. Connecteurs pro (\zh{期间}, \zh{此外}).
\item \textbf{Culture :} Concision, formalisme, respect de la hiérarchie (formules \zh{贵公司} guì gōngsī « votre honorable entreprise » en lettre de motivation).
\end{itemize}
\end{aretenir}

\coursec{Écriture des caractères complexes (1) : 简历, 经验}

Après avoir découvert la structure générale d'un CV chinois (\zh{简历}) et ses rubriques principales dans la section précédente, nous allons maintenant nous concentrer sur l'écriture précise des deux mots-clés qui forment le cœur de ce document : \zh{简历} (jiǎnlì, CV) et \zh{经验} (jīngyàn, expérience). Maîtriser ces caractères complexes — leurs radicaux, leur ordre des traits et leurs pièges visuels — est indispensable pour produire un document soigné, lisible et professionnel. Un trait manqué ou un radical confondu change le sens du mot et nuit à la crédibilité du candidat.

\courssub{Le caractère 简 (jiǎn) : « simple, bref »}

Observons d'abord le caractère \zh{简} dans son contexte lexical : \zh{简历} (jiǎnlì, CV), \zh{简单} (jiǎndān, simple), \zh{简介} (jiǎnjiè, résumé, brève introduction). Le sens commun est la simplicité, la concision, l'absence de superflu.

\begin{definition}[Analyse du caractère 简]
\textbf{Caractère :} \zh{简} (jiǎn) \\
\textbf{Structure :} Haut-bas (上下结构, shàngxià jiégòu) — le radical en haut, le composant en bas. \\
\textbf{Radical (clé) :} \zh{⺮} (zhú zì tóu, « bambou »), 6 traits. Il évoque l'origine du mot : les tablettes de bambou (\zh{简牍}, jiǎndú) utilisées pour écrire avant l'invention du papier. Un texte « simple » tient sur peu de tablettes. \\
\textbf{Composant du bas :} \zh{间} (jiān, intervalle, espace), qui donne aussi la prononciation. \\
\textbf{Nombre total de traits :} 13 (6 + 7).
\end{definition}

\begin{propriete}[Ordre des traits de 简 (13 traits) — Règle : haut avant bas]
Pour une structure haut-bas, on écrit d'abord la partie supérieure dans son intégralité, puis la partie inférieure.
\begin{enumerate}
    \item \textbf{Radical ⺮ (bambou, 6 traits) :} ① 撇 (piě) — ② 横 (héng) — ③ 点 (diǎn) — ④ 撇 (piě) — ⑤ 横 (héng) — ⑥ 点 (diǎn). Le radical bambou s'écrit comme deux petits « 个 » côte à côte : chaque moitié = 撇 + 横 + 点.
    \item \textbf{Partie basse 间 (7 traits) :} d'abord la porte \zh{门} : ⑦ 点 (diǎn) — ⑧ 竖 (shù) — ⑨ 横折钩 (héngzhégōu) ; puis le soleil \zh{日} à l'intérieur : ⑩ 竖 (shù) — (11) 横折 (héngzhé) — (12) 横 (héng) — (13) 横 (héng).
\end{enumerate}
\emph{Astuce :} pour \zh{门}, on commence toujours par le point en haut à gauche, puis le vertical gauche, puis le 横折钩 qui ferme à droite. Pour \zh{日}, on ferme la boîte par le trait horizontal du bas en dernier.
\end{propriete}

\begin{center}
\begin{tabularx}{\linewidth}{|X|X|X|X|}
\hline
\rowcolor{popblueL} \textbf{Caractères} & \textbf{Pinyin} & \textbf{Nature} & \textbf{Traduction} \\ \hline
简 & jiǎn & adj. & simple, bref \\ \hline
简历 & jiǎnlì & n. & CV (curriculum vitæ) \\ \hline
简单 & jiǎndān & adj. & simple, facile \\ \hline
简介 & jiǎnjiè & n. & résumé, brève présentation \\ \hline
\end{tabularx}
\end{center}

\courssub{Le caractère 历 (lì) : « parcours, histoire, calendrier »}

Dans \zh{简历}, le second caractère \zh{历} indique le « parcours », la « succession d'étapes » (d'où \zh{履历} lǚlì, autre mot pour CV, et \zh{历史} lìshǐ, histoire). C'est un caractère très compact.

\begin{definition}[Analyse du caractère 历]
\textbf{Caractère :} \zh{历} (lì) \\
\textbf{Structure :} Demi-encadrement supérieur gauche (左上包围结构, zuǒshàng bāowéi jiégòu) — le radical \zh{厂} enveloppe le composant \zh{力}. \\
\textbf{Radical :} \zh{厂} (hǎn, « falaise, rocher »), 2 traits. Sens originel : abri sous roche, passage. \\
\textbf{Composant interne :} \zh{力} (lì, « force, effort »), 2 traits. L'effort fourni au fil du temps forge le parcours. \\
\textbf{Nombre total de traits :} 4 (2 + 2). \emph{Attention : c'est un caractère à très peu de traits, chaque trait compte double.}
\end{definition}

\begin{propriete}[Ordre des traits de 历 (4 traits) — Règle : cadre avant contenu]
Pour une structure en demi-encadrement (\zh{厂}, \zh{广}, \zh{匚}…), on écrit d'abord le cadre extérieur, puis ce qu'il contient.
\begin{enumerate}
    \item \textbf{Radical 厂 (2 traits) :} ① 横 (héng, le trait horizontal du haut) — ② 撇 (piě, le long trait oblique qui descend vers la gauche).
    \item \textbf{Intérieur 力 (2 traits) :} ③ 横折钩 (héngzhégōu) — ④ 撇 (piě).
\end{enumerate}
\emph{Point de vigilance :} le 2ᵉ trait de \zh{厂} est un \textbf{撇} (oblique vers la gauche), pas un vertical. Le 4ᵉ trait (le 撇 de \zh{力}) part du centre du 横折钩 et descend vers la gauche sans toucher le cadre.
\end{propriete}

\courssub{Le caractère 经 (jīng) : « passer par, trame, expérience »}

Le mot \zh{经验} (jīngyàn) signifie littéralement « ce que l'on a traversé et vérifié ». \zh{经} porte l'idée de fil conducteur, de passage dans le temps (trame du tissu, classique bouddhique, gérer : \zh{已经} yǐjing, déjà ; \zh{经过} jīngguò, passer par).

\begin{definition}[Analyse du caractère 经]
\textbf{Caractère :} \zh{经} (jīng) \\
\textbf{Structure :} Gauche-droite (左右结构, zuǒyòu jiégòu). \\
\textbf{Radical :} \zh{纟} (sī, « soie, fil »), 3 traits. Version simplifiée de \zh{糸}. Il indique le domaine du fil, du textile, de la continuité. \\
\textbf{Composant de droite (phonétique) :} \zh{巠} (jīng), 5 traits, qui donne la prononciation. \\
\textbf{Nombre total de traits :} 8 (3 + 5).
\end{definition}

\begin{propriete}[Ordre des traits de 经 (8 traits) — Règle : gauche avant droite]
Pour une structure gauche-droite, on écrit d'abord la partie gauche (le radical) en entier, puis la partie droite.
\begin{enumerate}
    \item \textbf{Radical 纟 (3 traits) :} ① 撇折 (piězhé) — ② 撇折 (piězhé) — ③ 提 (tí, trait montant vers la droite).
    \item \textbf{Partie droite 巠 (5 traits) :} ④ 横撇 (héngpiě) — ⑤ 点 (diǎn) — ⑥ 横 (héng) — ⑦ 竖 (shù) — ⑧ 横 (héng, trait de base).
\end{enumerate}
\emph{Astuce :} les deux premiers traits du radical 纟 sont des \textbf{撇折} (oblique + cassure), et non des points ; le 3ᵉ trait 提 « lance » vers la partie droite. La partie droite se termine comme un petit \zh{工} (横 — 竖 — 横).
\end{propriete}

\courssub{Le caractère 验 (yàn) : « vérifier, éprouver, expérience »}

Second caractère de \zh{经验}, \zh{验} apporte la nuance de « test », de « preuve par la pratique ». C'est le plus complexe des quatre (10 traits).

\begin{definition}[Analyse du caractère 验]
\textbf{Caractère :} \zh{验} (yàn) \\
\textbf{Structure :} Gauche-droite (左右结构). \\
\textbf{Radical :} \zh{马} (mǎ, « cheval »), 3 traits (forme simplifiée). Le cheval sert à aller loin, à tester le terrain. \\
\textbf{Composant de droite (phonétique) :} \zh{佥} (qiān, « tous, ensemble »), 7 traits. \\
\textbf{Nombre total de traits :} 10 (3 + 7).
\end{definition}

\begin{propriete}[Ordre des traits de 验 (10 traits) — Règle : gauche avant droite (radical cheval d'abord)]
\begin{enumerate}
    \item \textbf{Radical 马 (3 traits) :} ① 横折 (héngzhé) — ② 竖折折钩 (shùzhézhégōu, le grand trait avec deux cassures et un crochet) — ③ 横 (héng, le trait horizontal du bas).
    \item \textbf{Partie droite 佥 (7 traits) :} ④ 撇 (piě) — ⑤ 捺 (nà, le toit \zh{人} s'écrit 撇 puis 捺) — ⑥ 横 (héng) — ⑦ 点 (diǎn) — ⑧ 点 (diǎn) — ⑨ 撇 (piě) — ⑩ 横 (héng, trait final de fermeture).
\end{enumerate}
\emph{Point de vigilance :} dans \zh{马}, le trait ③ 横 ne touche pas le bord droit ; il reste « suspendu » à gauche. Dans \zh{佥}, le toit \zh{人} doit couvrir toute la partie inférieure.
\end{propriete}

\begin{popfigure}
\begin{tikzpicture}[
    node distance=0.3cm and 0.5cm,
    strokebox/.style={draw=popdark, thick, rectangle, rounded corners=2pt, minimum width=1.8cm, minimum height=1cm, align=center, font=\small},
    arrow/.style={-Stealth, thick, popblue}
]
\node[strokebox, fill=popblueL, align=center] (jian){\zh{简} (jiǎn)\\13 traits};
\node[strokebox, fill=poporangeL, right=of jian, align=center] (li){\zh{历} (lì)\\4 traits};
\node[strokebox, fill=popgreenL, below=of jian, align=center] (jing){\zh{经} (jīng)\\8 traits};
\node[strokebox, fill=poppinkL, right=of jing, align=center] (yan){\zh{验} (yàn)\\10 traits};

\node[strokebox, fill=popcream, below left=0.5cm and -1.5cm of jian, align=center] (jian_rad){Radical : \zh{⺮} (bambou)\\6 traits};
\node[strokebox, fill=popcream, below right=0.5cm and -1.5cm of jian, align=center] (jian_ph){Bas : \zh{间} (jiān)\\7 traits};
\draw[arrow] (jian_rad.north) -- (jian.south west);
\draw[arrow] (jian_ph.north) -- (jian.south east);

\node[strokebox, fill=popcream, below left=0.5cm and -1.5cm of li, align=center] (li_rad){Radical : \zh{厂} (falaise)\\2 traits};
\node[strokebox, fill=popcream, below right=0.5cm and -1.5cm of li, align=center] (li_in){Intérieur : \zh{力} (force)\\2 traits};
\draw[arrow] (li_rad.north) -- (li.south west);
\draw[arrow] (li_in.north) -- (li.south east);

\node[strokebox, fill=popcream, below left=0.5cm and -1.5cm of jing, align=center] (jing_rad){Radical : \zh{纟} (soie)\\3 traits};
\node[strokebox, fill=popcream, below right=0.5cm and -1.5cm of jing, align=center] (jing_ph){Droite : \zh{巠} (jīng)\\5 traits};
\draw[arrow] (jing_rad.north) -- (jing.south west);
\draw[arrow] (jing_ph.north) -- (jing.south east);

\node[strokebox, fill=popcream, below left=0.5cm and -1.5cm of yan, align=center] (yan_rad){Radical : \zh{马} (cheval)\\3 traits};
\node[strokebox, fill=popcream, below right=0.5cm and -1.5cm of yan, align=center] (yan_ph){Droite : \zh{佥} (qiān)\\7 traits};
\draw[arrow] (yan_rad.north) -- (yan.south west);
\draw[arrow] (yan_ph.north) -- (yan.south east);

\node[above=0.2cm of jian, font=\bfseries\large, text=popdark] {Décomposition des quatre caractères cibles};
\end{tikzpicture}
\legende{Tableau de décomposition des quatre caractères cibles : radical (sémantique) + composant phonétique ou structurel = nombre total de traits.}
\end{popfigure}

\begin{popfigure}
\begin{tikzpicture}[
    stepnode/.style={draw=popdark, thick, rectangle, rounded corners=3pt, minimum width=2.4cm, minimum height=1.1cm, align=center, font=\small, fill=popcream},
    num/.style={circle, fill=popblue, text=white, font=\bfseries\tiny, minimum size=0.45cm, inner sep=0pt, anchor=south east}
]
\foreach \i/\desc in {
    1/横折 (马),
    2/竖折折钩 (马),
    3/横 (马),
    4/撇 (佥),
    5/捺 (佥),
    6/横 (佥),
    7/点 (佥),
    8/点 (佥),
    9/撇 (佥),
    10/横 (佥)
} {
    \pgfmathsetmacro{\col}{mod(\i-1,5)}
    \pgfmathsetmacro{\row}{-floor((\i-1)/5)}
    \node[stepnode] (n\i) at (\col*2.6cm, \row*1.5cm) {\desc};
    \node[num] at (n\i.north east) {\i};
}
\node[above=0.3cm of n1, font=\bfseries\large, text=popdark] {Frise d'écriture de \zh{验} (yàn) — 10 traits};
\end{tikzpicture}
\legende{Ordre des 10 traits de \zh{验} : d'abord le radical \zh{马} (traits 1 à 3), puis le composant phonétique \zh{佥} (traits 4 à 10), en commençant par le toit \zh{人} (撇 puis 捺).}
\end{popfigure}

\courssub{Caractères proches : distinctions cruciales}

Les francophones confondent souvent des caractères qui se ressemblent visuellement mais diffèrent par un radical ou un trait. Voici les trois paires les plus dangereuses pour cette leçon.

\begin{attention}[Pièges visuels : 简 / 间 ; 历 / 厉 ; 验 / 检]
\begin{itemize}
    \item \textbf{简 (jiǎn) vs 间 (jiān)} : même partie basse \zh{间}, mais radical différent. \zh{简} a \zh{⺮} (bambou) en haut $\rightarrow$ tablettes de bambou $\rightarrow$ « simple ». \zh{间} seul a \zh{门} (porte) autour de \zh{日} (soleil) $\rightarrow$ la lumière qui passe par la porte $\rightarrow$ « intervalle, pièce, entre ». \textbf{Ne pas écrire 间 à la place de 简 dans 简历 !}
    \item \textbf{历 (lì) vs 厉 (lì)} : même radical \zh{厂}, même prononciation, même ton ! \zh{历} contient \zh{力} (force, 2 traits : 横折钩 + 撇) $\rightarrow$ l'effort dans le temps $\rightarrow$ parcours. \zh{厉} contient \zh{万} (dix mille, 3 traits : 横 + 横折钩 + 撇) $\rightarrow$ « sévère, strict » (ex. : \zh{严厉} yánlì, sévère). \textbf{Le trait horizontal supplémentaire de 万 est le piège.}
    \item \textbf{验 (yàn) vs 检 (jiǎn)} : tous deux signifient « vérifier, examiner ». \zh{验} a le radical \zh{马} (cheval : tester sur le terrain). \zh{检} a le radical \zh{木} (bois). Ex. : \zh{体检} tǐjiǎn (bilan médical), \zh{检查} jiǎnchá (vérifier), \zh{验收} yànshōu (réceptionner après contrôle).
\end{itemize}
\end{attention}

\begin{exemplebox}[Mots-clés du CV : 简历 et 经验]
\begin{center}
\begin{tabularx}{\linewidth}{|X|X|X|}
\hline
\rowcolor{popblueL} \textbf{Chinois (caractères + pinyin)} & \textbf{Structure} & \textbf{Traduction française} \\ \hline
\zh{简历} jiǎnlì & 简 (simple) + 历 (parcours) & CV, curriculum vitæ \\ \hline
\zh{工作经验} gōngzuò jīngyàn & 工作 (travail) + 经验 (expérience) & expérience professionnelle \\ \hline
\zh{项目经验} xiàngmù jīngyàn & 项目 (projet) + 经验 & expérience de projet \\ \hline
\zh{实习经验} shíxí jīngyàn & 实习 (stage) + 经验 & expérience de stage \\ \hline
\zh{求职简历} qiúzhí jiǎnlì & 求职 (chercher un emploi) + 简历 & CV de candidature \\ \hline
\end{tabularx}
\end{center}
\emph{Exemples de phrases :}
\begin{itemize}
    \item \zh{我的简历很简短。} (Wǒ de jiǎnlì hěn jiǎnduǎn.) — Mon CV est très court.
    \item \zh{他有丰富的工作经验。} (Tā yǒu fēngfù de gōngzuò jīngyàn.) — Il a une riche expérience professionnelle.
    \item \zh{请把您的简历发给我。} (Qǐng bǎ nín de jiǎnlì fā gěi wǒ.) — Veuillez m'envoyer votre CV.
\end{itemize}
\end{exemplebox}

\begin{methode}[Mémoriser 经验 : « le fil et le cheval de l'expérience »]
Pour ne plus confondre l'ordre des caractères ni leurs radicaux, construisez cette image mentale en 3 temps :
\begin{enumerate}
    \item \textbf{纟 (soie, fil) :} l'expérience est un \textbf{fil conducteur} qui relie vos étapes passées à votre poste futur. C'est le radical de \zh{经} (jīng).
    \item \textbf{马 (cheval) :} pour acquérir de l'expérience, il faut \textbf{aller sur le terrain}, galoper, tester la route. C'est le radical de \zh{验} (yàn).
    \item \textbf{Synthèse :} \zh{经验} = \zh{纟} (le fil du temps) + \zh{马} (l'action sur le terrain). \emph{« L'expérience, c'est le fil que l'on tire en chevauchant. »}
\end{enumerate}
Écrivez trois fois \zh{经验} en prononçant : \emph{« Fil (jīng) — Cheval (yàn) »}.
\end{methode}

\begin{savaistu}[Le radical 纟 (sī) : la soie qui lie le langage]
Le radical \zh{纟} (forme simplifiée de \zh{糸}, « fil de soie fin ») apparaît dans plus de 200 caractères courants. Il marque invariablement une idée de \textbf{fil, tissu, continuité, lien, couleur} ou \textbf{document écrit} (autrefois sur soie).
\begin{center}
\begin{tabularx}{\linewidth}{|X|X|X|}
\hline
\rowcolor{poppurpleL} \textbf{Caractère} & \textbf{Pinyin} & \textbf{Lien sémantique avec « fil / soie »} \\ \hline
\zh{红} & hóng & rouge (couleur de la soie teinte, couleur porte-bonheur) \\ \hline
\zh{纸} & zhǐ & papier (qui a remplacé la soie pour l'écriture) \\ \hline
\zh{级} & jí & niveau, grade (les fils de la hiérarchie) \\ \hline
\zh{约} & yuē & accord, rendez-vous (nouer un fil, lien social) \\ \hline
\zh{纪} & jì & discipline, chronique (fil conducteur du temps) \\ \hline
\zh{练} & liàn & pratiquer, s'entraîner (travailler la soie, répéter le geste) \\ \hline
\zh{经} & jīng & trame, passer par, classique (le fil longitudinal du tissu) \\ \hline
\end{tabularx}
\end{center}
\emph{Repérer 纟 vous aide à deviner le sens d'un caractère inconnu : il s'agit souvent d'une matière souple, d'un lien, d'une suite ou d'un écrit.}
\end{savaistu}

\begin{exoresolu}[Thème / Version : mise en application immédiate]
\textbf{Énoncé :} Traduisez les phrases suivantes en chinois (caractères + pinyin) ou en français.
\begin{enumerate}
    \item \textit{Thème :} « Mon CV est simple et clair. »
    \item \textit{Thème :} « Elle a trois ans d'expérience professionnelle. »
    \item \textit{Version :} \zh{他的简历没有工作经验。}
    \item \textit{Version :} \zh{请检查我的简历。}
\end{enumerate}

\tcblower
\textbf{Solution détaillée :}
\begin{enumerate}
    \item \zh{我的简历简单明了。} (Wǒ de jiǎnlì jiǎndān míngliǎo.) \\
    \emph{Analyse :} « mon » = \zh{我的} (wǒ de) ; « CV » = \zh{简历} (jiǎnlì) ; « simple et clair » = \zh{简单明了} (jiǎndān míngliǎo, expression en 4 caractères : simple + clair). \textbf{Attention :} n'écrivez pas \zh{间历}.
    \item \zh{她有三年的工作经验。} (Tā yǒu sān nián de gōngzuò jīngyàn.) \\
    \emph{Analyse :} « elle » = \zh{她} (tā) ; « a » = \zh{有} (yǒu) ; « trois ans » = \zh{三年} (sān nián) + \zh{的} (de) ; « expérience professionnelle » = \zh{工作经验} (gōngzuò jīngyàn). \textbf{Ordre des mots :} durée + 的 + nom composé.
    \item \textbf{« Son CV ne mentionne pas d'expérience professionnelle. »} \\
    \emph{Analyse :} \zh{他} (tā, il) + \zh{的} (de) + \zh{简历} (jiǎnlì) + \zh{没有} (méiyǒu, ne pas avoir) + \zh{工作经验} (gōngzuò jīngyàn). \zh{没有} nie la possession ou l'existence (jamais \zh{不} devant \zh{有}).
    \item \textbf{« Veuillez vérifier mon CV. »} \\
    \emph{Analyse :} \zh{请} (qǐng, s'il vous plaît) + \zh{检查} (jiǎnchá, vérifier, inspecter — radical 木 !) + \zh{我的} (wǒ de) + \zh{简历} (jiǎnlì). \emph{Note :} on utilise le verbe composé \zh{检查} et non \zh{验} seul, qui ne s'emploie pas comme verbe libre.
\end{enumerate}
\end{exoresolu}

\begin{exercice}[Entraînement à l'écriture et à la distinction]
\begin{enumerate}
    \item \textbf{Écriture guidée :} sur votre cahier, écrivez 5 fois chaque caractère en respectant scrupuleusement l'ordre des traits : \zh{简}, \zh{历}, \zh{经}, \zh{验}. Numérotez chaque trait au fur et à mesure (1 à 13, 1 à 4, 1 à 8, 1 à 10).
    \item \textbf{Choix du bon caractère :} complétez par \zh{简} ou \zh{间} ; \zh{历} ou \zh{厉} ; \zh{验} ou \zh{检}.
    \begin{enumerate}
        \item 这道题很\_\_\_\_单。(Cet exercice est très simple.)
        \item 老师很\_\_\_\_害。(Le professeur est très sévère.)
        \item 我要\_\_\_\_收这个包裹。(Je dois réceptionner ce colis après contrôle.)
        \item 中喀两国之\_\_\_\_的关系很好。(Les relations entre la Chine et le Cameroun sont bonnes.)
    \end{enumerate}
    \item \textbf{Production courte :} écrivez 3 phrases en chinois (caractères + pinyin) pour vous présenter en vue d'un stage : « Je m'appelle... Je suis élève en Terminale A. J'ai fait un stage dans... Mon CV est joint. » (Utilisez \zh{简历}, \zh{经验}, \zh{实习}.)
\end{enumerate}
\end{exercice}

\begin{corrige}[Exercice : écriture et distinction]
\begin{enumerate}
    \item \emph{Autocorrection :} vérifiez votre cahier avec les frises de la leçon. Points clés : \zh{简} — le radical ⺮ = 撇、横、点、撇、横、点, et \zh{门} commence par le point ; \zh{历} — le 2ᵉ trait de \zh{厂} est un 撇, et \zh{力} = 横折钩 + 撇 ; \zh{经} — le radical 纟 = 撇折、撇折、提 ; \zh{验} — le radical 马 = 横折、竖折折钩、横, et \zh{佥} commence par le toit 撇 + 捺.
    \item \textbf{Choix :}
    \begin{enumerate}
        \item \zh{简} (jiǎn) — radical bambou pour « simple » : \zh{这道题很简单。}
        \item \zh{厉} (lì) — \zh{厂} + \zh{万} (3 traits) pour « sévère » : \zh{老师很厉害。}
        \item \zh{验} (yàn) — radical cheval pour « réceptionner après contrôle » : \zh{我要验收这个包裹。}
        \item \zh{间} (jiān) — radical porte pour « entre » : \zh{中喀两国之间的关系很好。}
    \end{enumerate}
    \item \textbf{Production (modèle) :}
    \begin{itemize}
        \item \zh{我叫杜邦，是高三A班的学生。} (Wǒ jiào Dùbāng, shì gāosān A bān de xuéshēng.) — Je m'appelle Dupont, je suis élève de Terminale A.
        \item \zh{我在雅温得的一家公司实习过。} (Wǒ zài Yǎwēndé de yì jiā gōngsī shíxí guo.) — J'ai fait un stage dans une entreprise de Yaoundé.
        \item \zh{我的简历在附件里，请查收。} (Wǒ de jiǎnlì zài fùjiàn lǐ, qǐng cháshōu.) — Mon CV est en pièce jointe, veuillez le réceptionner.
    \end{itemize}
\end{enumerate}
\end{corrige}

\coursec{Écriture des caractères complexes (2) : 教育, 技能, 聘}

\courssub{Transition et objectifs}

Après avoir maîtrisé les caractères de \zh{简历} (jiǎnlì, CV) et \zh{经验} (jīngyàn, expérience) dans la section précédente, nous abordons maintenant le cœur du contenu d'un CV : la formation (\zh{教育} jiàoyù), les compétences (\zh{技能} jìnéng) et le vocabulaire du recrutement (\zh{聘} pìn). Ces caractères sont graphiquement plus complexes ; ils mobilisent des radicaux productifs (\zh{攵}, \zh{扌}, \zh{耳}) et présentent des pièges de confusion fréquents pour les apprenants francophones (notamment \zh{技}/\zh{枝} et \zh{能}/\zh{熊}).

\begin{aretenir}[Objectifs de cette section]
\begin{itemize}
    \item Écrire correctement les 5 caractères : \zh{教}, \zh{育}, \zh{技}, \zh{能}, \zh{聘} (ordre des traits, nombre de traits, structure).
    \item Identifier et différencier les paires de caractères proches : \zh{教}/\zh{孝}, \zh{技}/\zh{枝}, \zh{能}/\zh{熊}.
    \item Comprendre la logique de composition (phonétique + sémantique) pour mémoriser durablement.
    \item Utiliser le vocabulaire associé : \zh{教育} (éducation, formation), \zh{技能} (compétence), \zh{招聘} (recruter).
\end{itemize}
\end{aretenir}

\courssub{Analyse graphique et étymologique des cinq caractères}

Nous procédons par ordre d'apparition dans les rubriques types d'un CV chinois : \zh{教育背景} (jiàoyù bèijǐng, Formation), \zh{专业技能} (zhuānyè jìnéng, Compétences professionnelles), \zh{求职意向} (qiúzhí yìxiàng, Poste recherché).

\begin{definition}[Caractère \zh{教} (jiào) — Enseigner]
\begin{itemize}
    \item \textbf{Structure} : Gauche-droite (\zh{左右结构}). Composé de \zh{孝} (xiào, piété filiale, 7 traits) à gauche et de \zh{攵} (pū, radical « taper / guider », 4 traits) à droite.
    \item \textbf{Nombre de traits} : 11 traits (7 + 4).
    \item \textbf{Logique étymologique} : \zh{孝} évoque l'enfant qui respecte l'aîné ; \zh{攵} représente une main tenant une baguette, symbole de la guidance et de la correction. « Guider l'enfant par l'exemple et la règle » → enseigner.
    \item \textbf{Radical} : \zh{攵} (radical n°66, 4 traits), souvent appelé \zh{反文旁} (fǎnwénpáng).
    \item \textbf{Attention à la prononciation} : \zh{教} se lit \emph{jiào} dans \zh{教育} (éducation), \zh{教师} (professeur), mais \emph{jiāo} (1\textsuperscript{er} ton) dans \zh{教书} (jiāoshū, enseigner, donner des cours) et \zh{我教你} (wǒ jiāo nǐ, je t'apprends).
\end{itemize}
\end{definition}

\begin{definition}[Caractère \zh{育} (yù) — Élever, éduquer, former]
\begin{itemize}
    \item \textbf{Structure} : Haut-bas (\zh{上下结构}). Partie haute : \zh{亠} (tóu, « couvercle », 2 traits) + \zh{厶} (sī, 2 traits) ; partie basse : \zh{月} (4 traits), qui est ici la forme compressée de \zh{肉} (ròu, « chair »).
    \item \textbf{Nombre de traits} : 8 traits.
    \item \textbf{Logique étymologique} : L'image ancienne évoque une mère qui nourrit et fait grandir son enfant (le \zh{月}/\zh{肉} suggère le corps). Sens premier : « faire grandir le corps », d'où « élever, éduquer, former l'esprit ».
    \item \textbf{Radical} : \zh{月} (radical n°74, 4 traits). Attention : ici c'est le radical « chair » (\zh{肉}) et non « lune », bien que la forme écrite soit identique.
\end{itemize}
\end{definition}

\begin{propriete}[Règle d'ordre des traits : structures gauche-droite et haut-bas]
\begin{enumerate}
    \item \textbf{Structure gauche-droite} (\zh{左右结构}) : on écrit d'abord la partie \textbf{gauche}, puis la partie \textbf{droite}. Exemple : \zh{教} → d'abord \zh{孝} (7 traits), puis \zh{攵} (4 traits).
    \item \textbf{Structure haut-bas} (\zh{上下结构}) : on écrit d'abord la partie \textbf{haute}, puis la partie \textbf{basse}. Exemple : \zh{育} → d'abord \zh{亠} + \zh{厶}, puis \zh{月}.
    \item À l'intérieur de chaque composant, les règles générales restent valables : horizontal avant vertical (\zh{横} avant \zh{竖}), haut avant bas, gauche avant droite.
\end{enumerate}
\end{propriete}

\begin{definition}[Caractère \zh{技} (jì) — Technique, habileté]
\begin{itemize}
    \item \textbf{Structure} : Gauche-droite (\zh{左右结构}). Gauche : \zh{扌} (shǒu, radical « main », 3 traits). Droite : \zh{支} (zhī, phonétique, 4 traits).
    \item \textbf{Nombre de traits} : 7 traits (3 + 4).
    \item \textbf{Logique} : La main (\zh{扌}) associée au son \zh{支} → l'habileté manuelle, la technique, le savoir-faire.
    \item \textbf{Radical} : \zh{扌} (shǒu, variante du radical n°64 \zh{手} « main », 3 traits).
    \item \textbf{Ordre des traits de \zh{支}} : \zh{横} (héng) → \zh{竖} (shù) → \zh{横撇} (héngpiě) → \zh{捺} (nà).
\end{itemize}
\end{definition}

\begin{definition}[Caractère \zh{能} (néng) — Capacité, pouvoir, être capable de]
\begin{itemize}
    \item \textbf{Structure} : Gauche-droite (\zh{左右结构}). Partie gauche : \zh{厶} (sī, 2 traits) au-dessus de \zh{月} (4 traits) ; partie droite : deux \zh{匕} (bǐ, 2 traits chacun) superposés.
    \item \textbf{Nombre de traits} : 10 traits (2 + 4 + 2 + 2).
    \item \textbf{Logique étymologique} : Le caractère ancien représentait un ours, animal symbole de force ; le sens s'est ensuite abstrait en « force → énergie → capacité ». Pour dire « ours », le chinois moderne a ajouté quatre points \zh{灬} (les pattes) sous \zh{能}, créant \zh{熊} (xióng).
    \item \textbf{Emploi grammatical} : \zh{能} est un verbe modal : \textbf{Sujet + 能 + Verbe} (« être capable de, savoir, pouvoir »). Exemple : \zh{我能说法语。} (Wǒ néng shuō Fǎyǔ.) — Je sais parler français.
\end{itemize}
\end{definition}

\begin{attention}[Piège majeur : \zh{能} (néng) vs \zh{熊} (xióng)]
\begin{itemize}
    \item \zh{能} (néng) : \textbf{10 traits}. Pas de points en dessous. Signification : capacité, pouvoir, énergie.
    \item \zh{熊} (xióng) : \textbf{14 traits}. C'est \zh{能} + \zh{灬} (quatre points en bas, qui représentent les pattes de l'animal). Signification : ours.
    \item \textbf{Astuce mnémotechnique} : l'ours (\zh{熊}) marche à quatre pattes → il a \textbf{quatre points} \zh{灬} sous le corps. La capacité (\zh{能}) est abstraite → elle n'a pas de pattes.
    \item \textbf{Erreur fréquente des francophones} : écrire \zh{我熊说中文。} au lieu de \zh{我能说中文。} (Wǒ néng shuō Zhōngwén. — Je sais parler chinois). En dictée, vérifier toujours la présence ou l'absence des quatre points.
\end{itemize}
\end{attention}

\begin{definition}[Caractère \zh{聘} (pìn) — Engager, recruter (formel)]
\begin{itemize}
    \item \textbf{Structure} : Gauche-droite (\zh{左右结构}). Gauche : \zh{耳} (ěr, radical « oreille », 6 traits). Droite : \zh{甹} (pìn, composant phonétique, 7 traits).
    \item \textbf{Nombre de traits} : 13 traits (6 + 7).
    \item \textbf{Logique} : L'oreille (\zh{耳}) qui écoute une proposition → recevoir et accepter un engagement → recruter, embaucher. Le composant droit \zh{甹} donne le son \emph{pìn}.
    \item \textbf{Radical} : \zh{耳} (ěr, radical n°128, 6 traits).
    \item \textbf{Ordre des traits de \zh{耳}} : \zh{横} → \zh{竖} → \zh{竖} → \zh{横} → \zh{横} → \zh{横} (6 traits).
    \item \textbf{Vocabulaire clé} : \zh{招聘} (zhāopìn, recruter), \zh{聘用} (pìnyòng, embaucher), \zh{聘请} (pìnqǐng, inviter/engager un expert).
\end{itemize}
\end{definition}

\courssub{Règle d'or : le radical \zh{扌} (main) dans les caractères gauche-droite}

\begin{propriete}[Règle — Écriture du radical \zh{扌} (shǒuzìpáng)]
Dans les caractères de structure gauche-droite ayant le radical \zh{扌} (main) à gauche :
\begin{enumerate}
    \item Le radical s'écrit \textbf{en premier} (avant la partie droite).
    \item \textbf{Ordre des 3 traits du radical \zh{扌}} :
        \begin{enumerate}
            \item \zh{横} (héng) — trait horizontal court ;
            \item \zh{竖钩} (shùgōu) — trait vertical avec crochet ;
            \item \zh{提} (tí) — trait montant vers la droite.
        \end{enumerate}
    \item La partie droite s'écrit ensuite, selon ses propres règles internes.
\end{enumerate}
\emph{Exemples} : \zh{技} (jì), \zh{打} (dǎ), \zh{提} (tí), \zh{招} (zhāo), \zh{持} (chí).
\end{propriete}

\begin{methode}[Méthode pour compter et écrire les traits sans erreur]
Pour tout nouveau caractère complexe, suivez ce rituel en 4 étapes :
\begin{enumerate}
    \item \textbf{Décomposer} : identifier la structure (gauche-droite, haut-bas, encadrement) et les composants (radical + phonétique).
    \item \textbf{Nommer} : écrire le caractère très lentement sur papier quadrillé (\zh{格子纸} gézǐzhǐ), en disant à voix haute le nom de chaque trait dans l'ordre standard : \zh{横} (héng), \zh{竖} (shù), \zh{撇} (piě), \zh{捺} (nà), \zh{点} (diǎn), \zh{提} (tí), \zh{折} (zhé), \zh{钩} (gōu).
    \item \textbf{Compter} : cocher chaque trait prononcé et vérifier le total (ex. : 7 pour \zh{技}, 10 pour \zh{能}, 13 pour \zh{聘}).
    \item \textbf{Valider} : comparer avec le modèle standard, puis réécrire 5 fois de suite sans regarder le modèle.
\end{enumerate}
\emph{Conseil} : utilisez un stylo fin pour bien sentir la pression (plein / délié) des traits \zh{撇} et \zh{捺}.
\end{methode}

\courssub{Tableau récapitulatif des 5 caractères cibles}

\begin{center}
\begin{tabularx}{\linewidth}{|X|X|X|X|X|X|}
\hline
\rowcolor{popblueL} \textbf{Caractère} & \textbf{Pinyin} & \textbf{Structure} & \textbf{Radical (clé)} & \textbf{Nb traits} & \textbf{Sens principal} \\
\hline
\zh{教} & jiào / jiāo & Gauche-droite (\zh{孝} + \zh{攵}) & \zh{攵} (pū) & 11 & Enseigner, éducation \\
\hline
\zh{育} & yù & Haut-bas (\zh{亠} + \zh{厶} + \zh{月}) & \zh{月} (ròu) & 8 & Élever, éduquer, former \\
\hline
\zh{技} & jì & Gauche-droite (\zh{扌} + \zh{支}) & \zh{扌} (shǒu) & 7 & Technique, habileté \\
\hline
\zh{能} & néng & Gauche-droite (\zh{厶} + \zh{月} + \zh{匕}×2) & \zh{月} (ròu) & 10 & Capacité, pouvoir \\
\hline
\zh{聘} & pìn & Gauche-droite (\zh{耳} + \zh{甹}) & \zh{耳} (ěr) & 13 & Engager, recruter (formel) \\
\hline
\end{tabularx}
\end{center}

\courssub{Focus sur les paires confusibles : analyse contrastive}

Les trois paires ci-dessous sont sources d'erreurs systématiques à l'examen (écrit et dictée). Le schéma ci-dessous synthétise les différences graphiques et sémantiques.

\begin{popfigure}
\begin{tikzpicture}[
    node distance=0.8cm and 1.2cm,
    box/.style={draw=popblue, thick, rounded corners, fill=popblueL, text width=3.4cm, align=center, minimum height=2.6cm},
    boxtitle/.style={draw=popdark, thick, rounded corners, fill=popgold, text width=3.4cm, align=center, font=\bfseries},
    vs/.style={font=\bfseries\large, text=poppink}
]
% Paire 1: 教 vs 孝
\node[boxtitle] (t1) at (0,0) {\zh{教} (jiào)};
\node[box, align=center] (c1) [below of=t1]{Gauche : \zh{孝}\\Droite : \zh{攵} (guider)\\\textbf{11 traits}\\\emph{Enseigner}};
\node[boxtitle] (t2) at (5,0) {\zh{孝} (xiào)};
\node[box, align=center] (c2) [below of=t2]{Haut : \zh{耂} (vieux)\\Bas : \zh{子} (enfant)\\\textbf{7 traits}\\\emph{Piété filiale}};
\node[vs] at (2.5,-1.4) {VS};

% Paire 2: 技 vs 枝
\node[boxtitle] (t3) at (0,-5) {\zh{技} (jì)};
\node[box, align=center] (c3) [below of=t3]{Gauche : \zh{扌} (main)\\3 traits\\Droite : \zh{支}\\\textbf{7 traits}\\\emph{Technique}};
\node[boxtitle] (t4) at (5,-5) {\zh{枝} (zhī)};
\node[box, align=center] (c4) [below of=t4]{Gauche : \zh{木} (bois)\\4 traits\\Droite : \zh{支}\\\textbf{8 traits}\\\emph{Branche}};
\node[vs] at (2.5,-6.4) {VS};

% Paire 3: 能 vs 熊
\node[boxtitle] (t5) at (0,-10) {\zh{能} (néng)};
\node[box, align=center] (c5) [below of=t5]{\zh{厶} + \zh{月} + \zh{匕}×2\\\textbf{Sans les 4 points}\\\textbf{10 traits}\\\emph{Capacité, pouvoir}};
\node[boxtitle] (t6) at (5,-10) {\zh{熊} (xióng)};
\node[box, align=center] (c6) [below of=t6]{\zh{能} + \zh{灬} (4 points = pattes)\\\textbf{14 traits}\\\emph{Ours}};
\node[vs] at (2.5,-11.4) {VS};
\end{tikzpicture}
\legende{Tableau comparatif des paires confusibles : \zh{教}/\zh{孝}, \zh{技}/\zh{枝}, \zh{能}/\zh{熊}. Noter la différence de radical (main \zh{扌} vs bois \zh{木}) et la présence ou l'absence des quatre points \zh{灬}.}
\end{popfigure}

\begin{exemplebox}[Exemples contextualisés : paires confusibles]
\begin{enumerate}
    \item \zh{教} vs \zh{孝} :
    \begin{itemize}
        \item \zh{他教中文。} (Tā jiāo Zhōngwén.) — Il enseigne le chinois.
        \item \zh{中国人重视孝道。} (Zhōngguórén zhòngshì xiàodào.) — Les Chinois accordent de l'importance à la piété filiale.
    \end{itemize}
    \item \zh{技} vs \zh{枝} :
    \begin{itemize}
        \item \zh{我的专业技能很强。} (Wǒ de zhuānyè jìnéng hěn qiáng.) — Mes compétences professionnelles sont solides.
        \item \zh{这棵树的枝很粗壮。} (Zhè kē shù de zhī hěn cūzhuàng.) — Les branches de cet arbre sont très épaisses.
    \end{itemize}
    \item \zh{能} vs \zh{熊} :
    \begin{itemize}
        \item \zh{我能说法语和汉语。} (Wǒ néng shuō Fǎyǔ hé Hànyǔ.) — Je sais parler français et chinois.
        \item \zh{大熊猫是中国的国宝。} (Dàxióngmāo shì Zhōngguó de guóbǎo.) — Le panda géant est un trésor national de la Chine.
    \end{itemize}
\end{enumerate}
\end{exemplebox}

\courssub{Schémas de structure : \zh{技} et \zh{聘} (gauche-droite)}

Ces deux caractères partagent la structure gauche-droite mais avec des radicaux différents (\zh{扌} vs \zh{耳}). Maîtriser l'alignement horizontal (hauteur équilibrée des deux parties) est crucial pour la lisibilité.

\begin{popfigure}
\begin{tikzpicture}[
    node distance=0.5cm and 0.3cm,
    rad/.style={draw=popblue, thick, fill=popblueL, rounded corners, minimum width=2cm, minimum height=1.4cm, align=center, font=\large},
    phon/.style={draw=popgreen, thick, fill=popgreenL, rounded corners, minimum width=2cm, minimum height=1.4cm, align=center, font=\large},
    lab/.style={font=\small, text=popmuted, align=center},
    title/.style={font=\bfseries\large, text=popdark, align=center}
]
% Caractère 技
\node[title] at (1.25, 1.6) {\zh{技} (jì) — Gauche-Droite};
\node[rad, align=center] (jl) at (0, 0){\zh{扌}\\{\small main, 3 traits}};
\node[phon, align=center] (jr) at (2.5, 0){\zh{支}\\{\small phonétique, 4 traits}};
\node[lab] [below of=jl] {Écrit en 1\textsuperscript{er}};
\node[lab] [below of=jr] {Écrit en 2\textsuperscript{e}};
\draw[thick, poporange, -Stealth] (jl.east) -- (jr.west);

% Caractère 聘
\node[title] at (7.25, 1.6) {\zh{聘} (pìn) — Gauche-Droite};
\node[rad, align=center] (pl) at (6, 0){\zh{耳}\\{\small oreille, 6 traits}};
\node[phon, align=center] (pr) at (8.5, 0){\zh{甹}\\{\small phonétique, 7 traits}};
\node[lab] [below of=pl] {Écrit en 1\textsuperscript{er}};
\node[lab] [below of=pr] {Écrit en 2\textsuperscript{e}};
\draw[thick, poporange, -Stealth] (pl.east) -- (pr.west);

% Détails traits
\node[lab, text width=4.2cm] at (1.25, -2.2) {Ordre de \zh{扌} : \zh{横} → \zh{竖钩} → \zh{提}\\Ordre de \zh{支} : \zh{横} → \zh{竖} → \zh{横撇} → \zh{捺}};
\node[lab, text width=4.2cm] at (7.25, -2.2) {Ordre de \zh{耳} : \zh{横} → \zh{竖} → \zh{竖} → \zh{横} → \zh{横} → \zh{横}\\Puis \zh{甹} (7 traits) à droite};
\end{tikzpicture}
\legende{Schéma de structure gauche-droite pour \zh{技} et \zh{聘}. Le radical (bleu) s'écrit toujours avant le composant phonétique (vert) ; les deux parties doivent être visuellement équilibrées en hauteur.}
\end{popfigure}

\courssub{Vocabulaire clé du CV : mots composés et traduction}

Les caractères étudiés forment les mots « pivots » des rubriques de CV. Maîtrisez ces associations (collocations).

\begin{center}
\begin{tabularx}{\linewidth}{|X|X|X|X|}
\hline
\rowcolor{popblueL} \textbf{Caractères} & \textbf{Pinyin} & \textbf{Nature} & \textbf{Traduction / Contexte CV} \\
\hline
教育 & jiàoyù & Nom / Verbe & Éducation, formation. Rubrique : \zh{教育背景} (jiàoyù bèijǐng) = Formation, parcours scolaire. \\
\hline
技能 & jìnéng & Nom & Compétence, savoir-faire. Rubrique : \zh{专业技能} (zhuānyè jìnéng) = Compétences professionnelles. \\
\hline
招聘 & zhāopìn & Verbe & Recruter. \zh{招聘会} (zhāopìnhuì) = salon de l'emploi ; \zh{招聘广告} (zhāopìn guǎnggào) = offre d'emploi. \\
\hline
聘用 & pìnyòng & Verbe & Embaucher (acte formel). \zh{被聘用} (bèi pìnyòng) = être embauché. \\
\hline
教师 & jiàoshī & Nom & Professeur, enseignant. \zh{语言教师} (yǔyán jiàoshī) = professeur de langue. \\
\hline
培养 & péiyǎng & Verbe & Former, cultiver (un talent). \zh{培养目标} (péiyǎng mùbiāo) = objectif de formation. \\
\hline
能力 & nénglì & Nom & Capacité, aptitude. \zh{语言能力} (yǔyán nénglì) = compétences linguistiques ; \zh{沟通能力} (gōutōng nénglì) = capacités de communication. \\
\hline
\end{tabularx}
\end{center}

\begin{savaistu}[Étymologie culturelle : \zh{教育} (jiàoyù) — « Guider et faire grandir »]
Le mot \zh{教育} combine deux verbes complémentaires :
\begin{itemize}
    \item \zh{教} (jiào) : \textbf{guider, montrer, instruire}. C'est l'action extérieure de l'enseignant (le \zh{攵}, main qui guide).
    \item \zh{育} (yù) : \textbf{élever, nourrir, faire mûrir}. C'est le processus intérieur de l'apprenant (le \zh{月}/\zh{肉}, le corps qui grandit).
\end{itemize}
Dans la pensée chinoise classique, transmettre sans faire grandir est de l'endoctrinement ; laisser grandir sans guider est de l'abandon. L'éducation véritable (\zh{教育}) unit la transmission du savoir et le développement de la personne. Un proverbe chinois célèbre le dit : \zh{十年树木，百年树人。} (Shí nián shù mù, bǎi nián shù rén.) — « Dix ans pour faire pousser un arbre, cent ans pour former un homme. » La rubrique \zh{教育背景} d'un CV témoigne de ce long processus de formation.
\end{savaistu}

\courssub{Entraînement guidé à la copie (5 lignes par caractère)}

\begin{methode}[Consignes pour l'entraînement graphique sur \zh{格子纸}]
\begin{enumerate}
    \item Utilisez une feuille à carreaux chinois (\zh{格子纸} gézǐzhǐ) : un grand carré divisé en 4 (ou 9) petits carrés.
    \item Écrivez \textbf{un seul caractère par grand carré}, bien centré.
    \item Respectez l'\textbf{ordre des traits} officiel (voir les tableaux ci-dessus).
    \item Faites \textbf{5 répétitions} par caractère ; pour chaque répétition, prononcez le pinyin et le sens.
    \item \textbf{Auto-vérification} : après les 5 lignes, entourez votre plus belle occurrence et vérifiez : nombre de traits ? proportions ? position du radical ?
\end{enumerate}
\end{methode}

\begin{exercice}[Entraînement à la copie — sur cahier ou feuille à carreaux]
Recopiez chaque caractère ci-dessous \textbf{5 fois} en respectant l'ordre des traits.
\begin{enumerate}
    \item \zh{教} (11 traits) — \_\_\_ \_\_\_ \_\_\_ \_\_\_ \_\_\_
    \item \zh{育} (8 traits) — \_\_\_ \_\_\_ \_\_\_ \_\_\_ \_\_\_
    \item \zh{技} (7 traits) — \_\_\_ \_\_\_ \_\_\_ \_\_\_ \_\_\_
    \item \zh{能} (10 traits) — \_\_\_ \_\_\_ \_\_\_ \_\_\_ \_\_\_
    \item \zh{聘} (13 traits) — \_\_\_ \_\_\_ \_\_\_ \_\_\_ \_\_\_
\end{enumerate}
\emph{Défi} : écrivez les deux mots composés \zh{教育} et \zh{技能} 3 fois chacun, en respectant l'espacement entre les caractères.
\end{exercice}

\courssub{Exercice résolu : discrimination graphique et choix lexical}

\begin{exoresolu}[Compléter avec le bon caractère : \zh{技}/\zh{枝}, \zh{能}/\zh{熊}, \zh{教}/\zh{孝}]
\textbf{Énoncé :} Choisissez le caractère correct entre parenthèses pour compléter les phrases. Écrivez le caractère, le pinyin et traduisez.

1. 这棵树的 \_\_\_ (技/枝) 很粗壮。

2. 他 \_\_\_ (能/熊) 说流利的英语。

3. 我们要尊敬老人，懂得 \_\_\_ (教/孝) 道。

4. 简历上要写清楚专业 \_\_\_ \_\_\_ (技能/枝熊)。

5. 签了劳动合同以后，公司正式 \_\_\_ 用他。(聘/请)

\tcblower

\textbf{Solution détaillée :}

1. \zh{枝} (zhī) — \emph{branche}. Contexte : arbre (\zh{树}) + description physique (\zh{粗壮}, épais et solide). Le radical bois \zh{木} est exigé ; \zh{技} a le radical main.
    \zh{这棵树的枝很粗壮。} (Zhè kē shù de zhī hěn cūzhuàng.) — Les branches de cet arbre sont très épaisses.

2. \zh{能} (néng) — \emph{pouvoir, être capable de}. Structure : verbe modal + verbe (\zh{说}). \zh{熊} est un nom (ours) et ne peut pas précéder un verbe.
    \zh{他能说流利的英语。} (Tā néng shuō liúlì de Yīngyǔ.) — Il sait parler anglais couramment.

3. \zh{孝} (xiào) — \emph{piété filiale}. Mot composé figé : \zh{孝道} (xiàodào, le devoir filial).
    \zh{我们要尊敬老人，懂得孝道。} (Wǒmen yào zūnjìng lǎorén, dǒngde xiàodào.) — Nous devons respecter les personnes âgées et comprendre le devoir filial.

4. \zh{技能} (jìnéng) — \emph{compétences}. Mot composé standard du CV : \zh{技} (technique) + \zh{能} (capacité). Ni \zh{枝} (branche) ni \zh{熊} (ours) ne conviennent.
    \zh{简历上要写清楚专业技能。} (Jiǎnlì shang yào xiě qīngchu zhuānyè jìnéng.) — Sur le CV, il faut écrire clairement les compétences professionnelles.

5. \zh{聘} (pìn) — \emph{engager formellement}. Collocation : \zh{聘用} (pìnyòng, embaucher). \zh{请用} (qǐngyòng) signifierait « je vous invite à (manger, utiliser) », jamais « embaucher ».
    \zh{签了劳动合同以后，公司正式聘用他。} (Qiān le láodòng hétong yǐhòu, gōngsī zhèngshì pìnyòng tā.) — Après la signature du contrat de travail, l'entreprise l'embauche officiellement.
\end{exoresolu}

\courssub{Exercices d'application (travail en autonomie)}

\begin{exercice}[Discrimination et production écrite]
\begin{enumerate}
    \item \textbf{Dictée de mots} : écrivez en caractères, avec pinyin et traduction : \zh{教育}, \zh{技能}, \zh{招聘}, \zh{能力}, \zh{教师}, \zh{培养}.
    \item \textbf{Recherche d'erreurs} : identifiez l'erreur de caractère dans chaque phrase et corrigez-la.
        \begin{enumerate}
            \item 我的枝能很强。
            \item 他是一位优秀的孝师。
            \item 我熊说中文。
            \item 这家公司正在招骋新员工。(attention au caractère droit de \zh{聘})
        \end{enumerate}
    \item \textbf{Mini-rédaction CV (5 lignes)} : rédigez les rubriques \zh{教育背景} et \zh{专业技能} pour un étudiant camerounais fictif (ou pour vous-même), en chinois simplifié avec pinyin.
        \emph{Modèles} :
        \zh{教育背景：2020年—2023年，雅温得大学，法语专业，学士学位。}
        \zh{专业技能：法语流利，汉语HSK四级，办公软件熟练，沟通能力强。}
\end{enumerate}
\end{exercice}

\begin{corrige}[Exercice : Discrimination et production écrite]
\textbf{1. Dictée (corrigé) :}
\begin{itemize}
    \item \zh{教育} (jiàoyù) — éducation, formation.
    \item \zh{技能} (jìnéng) — compétence, savoir-faire.
    \item \zh{招聘} (zhāopìn) — recruter.
    \item \zh{能力} (nénglì) — capacité, aptitude.
    \item \zh{教师} (jiàoshī) — professeur, enseignant.
    \item \zh{培养} (péiyǎng) — former, cultiver.
\end{itemize}

\textbf{2. Corrections :}
\begin{enumerate}
    \item \zh{枝} → \zh{技} : \zh{我的\textbf{技}能很强。} (Wǒ de jìnéng hěn qiáng.) — Mes compétences sont solides.
    \item \zh{孝} → \zh{教} : \zh{他是一位优秀的\textbf{教}师。} (Tā shì yí wèi yōuxiù de jiàoshī.) — C'est un excellent professeur.
    \item \zh{熊} → \zh{能} : \zh{我\textbf{能}说中文。} (Wǒ néng shuō Zhōngwén.) — Je sais parler chinois.
    \item \zh{骋} → \zh{聘} : \zh{这家公司正在招\textbf{聘}新员工。} (Zhè jiā gōngsī zhèngzài zhāopìn xīn yuángōng.) — Cette entreprise recrute de nouveaux employés. \emph{Attention} : \zh{骋} (chěng, « galoper », radical cheval \zh{马}) n'existe que dans \zh{驰骋} ; le recrutement exige \zh{聘} (radical oreille \zh{耳}).
\end{enumerate}

\textbf{3. Proposition de rédaction (exemple corrigé) :}

\zh{教育背景：2021年—2024年，布埃亚大学，国际贸易专业，学士学位。}

(Jiàoyù bèijǐng : 2021 nián — 2024 nián, Bù'āiyà Dàxué, guójì màoyì zhuānyè, xuéshì xuéwèi.)

« Formation : 2021-2024, Université de Buea, spécialité Commerce international, licence. »

\zh{专业技能：英语流利，汉语HSK四级，办公软件熟练，沟通能力强。}

(Zhuānyè jìnéng : Yīngyǔ liúlì, Hànyǔ HSK sì jí, bàngōng ruǎnjiàn shúliàn, gōutōng nénglì qiáng.)

« Compétences professionnelles : anglais courant, chinois HSK niveau 4, maîtrise des logiciels de bureautique, fortes capacités de communication. »
\end{corrige}

\coursec{Modèle de CV commenté, structure et connecteurs}

Dans la section précédente, nous avons appris à écrire, trait par trait, les caractères complexes \zh{教育} (jiàoyù, « éducation, formation »), \zh{技能} (jìnéng, « compétence ») et \zh{聘} (pìn, « embaucher »). Nous savons maintenant tracer les mots-clés d'un CV. Il reste une étape essentielle : savoir \textbf{assembler} ces mots en phrases correctes et organiser ces phrases en un texte complet, clair et bien ponctué. C'est l'objet de cette section : lire, comprendre et commenter un modèle de CV chinois simple, puis en extraire la structure et les connecteurs réutilisables.

\courssub{Lecture guidée du modèle complet}

\begin{textebox}[Un autoportrait-CV : 李明 se présente]
\zh{我叫李明。首先，我介绍我的教育：2020年到2024年，我在雅温得中学学习。然后，我说我的技能：我会说中文、法语和英语，而且我会用电脑。我也喜欢足球和音乐。}\\[4pt]
\emph{Wǒ jiào Lǐ Míng. Shǒuxiān, wǒ jièshào wǒ de jiàoyù: èr-líng-èr-líng nián dào èr-líng-èr-sì nián, wǒ zài Yǎwēndé zhōngxué xuéxí. Ránhòu, wǒ shuō wǒ de jìnéng: wǒ huì shuō Zhōngwén, Fǎyǔ hé Yīngyǔ, érqiě wǒ huì yòng diànnǎo. Wǒ yě xǐhuan zúqiú hé yīnyuè.}\\[4pt]
« Je m'appelle Li Ming. D'abord, je présente ma formation : de 2020 à 2024, j'ai étudié au lycée de Yaoundé. Ensuite, je parle de mes compétences : je sais parler chinois, français et anglais, et de plus je sais utiliser l'ordinateur. J'aime aussi le football et la musique. »
\end{textebox}

\textbf{Questions de compréhension (oral ou écrit) :}
\begin{enumerate}
\item \zh{他叫什么名字？} \emph{Tā jiào shénme míngzi?} — Comment s'appelle-t-il ?
\item \zh{他在哪里学习？} \emph{Tā zài nǎlǐ xuéxí?} — Où a-t-il étudié ?
\item \zh{他会说什么语言？} \emph{Tā huì shuō shénme yǔyán?} — Quelles langues sait-il parler ?
\item \zh{他喜欢什么？} \emph{Tā xǐhuan shénme?} — Qu'aime-t-il ?
\end{enumerate}

\begin{center}
\begin{tabularx}{\linewidth}{|X|X|X|X|}
\hline
\rowcolor{popblueL} Caractères & Pinyin & Nature & Traduction \\
\hline
介绍 & jièshào & verbe & présenter, introduire \\
\hline
教育 & jiàoyù & nom & formation, éducation \\
\hline
技能 & jìnéng & nom & compétence \\
\hline
会 & huì & verbe modal & savoir (faire), être capable de \\
\hline
用 & yòng & verbe & utiliser \\
\hline
电脑 & diànnǎo & nom & ordinateur \\
\hline
爱好 & àihào & nom & loisir, passe-temps \\
\hline
首先 & shǒuxiān & connecteur & d'abord, en premier \\
\hline
然后 & ránhòu & connecteur & ensuite, puis \\
\hline
\end{tabularx}
\end{center}

\courssub{L'en-tête du CV chinois : les rubriques d'identité}

Un CV chinois commence toujours par un \textbf{en-tête} qui donne les informations d'identité. En chinois, on présente ces informations sous forme d'étiquettes suivies de deux-points \zh{：}.

\begin{center}
\begin{tabularx}{\linewidth}{|X|X|X|X|}
\hline
\rowcolor{popblueL} Caractères & Pinyin & Nature & Traduction \\
\hline
姓名 & xìngmíng & nom & nom et prénom \\
\hline
性别 & xìngbié & nom & sexe (\zh{男} nán « masculin » / \zh{女} nǚ « féminin ») \\
\hline
出生日期 & chūshēng rìqī & expression & date de naissance \\
\hline
地址 & dìzhǐ & nom & adresse \\
\hline
电话 & diànhuà & nom & téléphone \\
\hline
邮箱 & yóuxiāng & nom & courriel \\
\hline
\end{tabularx}
\end{center}

\begin{exemplebox}[Un en-tête complet]
\zh{姓名：李明。性别：男。出生日期：2006年5月12日。地址：雅温得。电话：690 12 34 56。邮箱：liming@email.com。}\\
\emph{Xìngmíng: Lǐ Míng. Xìngbié: nán. Chūshēng rìqī: èr-líng-líng-liù nián wǔ yuè shí'èr rì. Dìzhǐ: Yǎwēndé. Diànhuà: … Yóuxiāng: …}\\
« Nom : Li Ming. Sexe : masculin. Date de naissance : 12 mai 2006. Adresse : Yaoundé. Téléphone : … Courriel : … »
\end{exemplebox}

\begin{attention}[L'ordre des dates et des adresses]
En chinois, on va toujours \textbf{du plus grand au plus petit} : année → mois → jour (\zh{2006年5月12日}), pays → ville → quartier. C'est l'inverse de l'habitude française « 12 mai 2006 ». De même, le nom de famille précède le prénom : \zh{李明}, Lǐ est le nom de famille, Míng le prénom.
\end{attention}

\begin{savaistu}[CV en Chine et au Cameroun : quelques différences]
En Chine, le CV (简历 \zh{jiǎnlì}) inclut souvent une photo d'identité en haut à droite, l'état civil détaillé (âge, origine ethnique \zh{民族} \emph{mínzú}, situation familiale \zh{婚姻状况} \emph{hūnyīn zhuàngkuàng}) et parfois le \zh{户口} \emph{hùkǒu} (registre de résidence). Au Cameroun, comme en France, la photo n'est pas obligatoire et les informations personnelles sensibles (âge, situation familiale, origine) sont déconseillées, voire interdites par les règles de non-discrimination. Pour un stage ou un emploi dans un contexte international (entreprise chinoise au Cameroun, ambassade), adaptez le modèle : gardez la structure chinoise (rubriques, ordre chronologique) mais allégez l'en-tête aux standards internationaux (nom, contact, ville).
\end{savaistu}

\courssub{Les rubriques du corps du CV}

\courssub{Rubrique formation : 教育}

\begin{exemplebox}[La phrase de formation]
\zh{2020年到2024年，我在雅温得中学学习。}\\
\emph{Èr-líng-èr-líng nián dào èr-líng-èr-sì nián, wǒ zài Yǎwēndé zhōngxué xuéxí.}\\
« De 2020 à 2024, j'ai étudié au lycée de Yaoundé. »
\end{exemplebox}

Observons cette phrase. Elle suit un ordre très régulier, typique du chinois : \textbf{le temps d'abord, puis le sujet, puis le lieu, puis l'action}. La préposition \zh{在} (zài, « à, en ») introduit le lieu, placé \emph{avant} le verbe, contrairement au français (« j'étudie \emph{au lycée} » : le lieu vient après).

\begin{propriete}[Structure type d'une phrase de CV]
\textbf{[Temps] + Sujet + 在 + [Lieu] + Verbe (+ complément)}\\
Le complément de temps (avec \zh{从} cóng « depuis » ou \zh{到} dào « jusqu'à ») se place au tout début de la phrase ou juste après le sujet, jamais à la fin. Le complément de lieu introduit par \zh{在} se place \textbf{avant} le verbe.
\end{propriete}

\begin{popfigure}
\begin{tikzpicture}[node distance=0.35cm,
  bloc/.style={rectangle, rounded corners=3pt, draw=popblue, fill=popblueL, minimum height=0.9cm, align=center, font=\small},
  plus/.style={font=\large\bfseries, text=popdark}]
\node[bloc, align=center] (t){时间\\ temps};
\node[plus, right=of t] (p1) {+};
\node[bloc, right=of p1, fill=poporangeL, draw=poporange, align=center] (s){我\\ sujet};
\node[plus, right=of s] (p2) {+};
\node[bloc, right=of p2, fill=popgreenL, draw=popgreen, align=center] (z){在 + 地点\\ lieu};
\node[plus, right=of z] (p3) {+};
\node[bloc, right=of p3, fill=poppurpleL, draw=poppurple, align=center] (v){动词\\ verbe};
\node[below=0.5cm of s, align=center, font=\small\itshape, text=popmuted] {2020年到2024年，我 在 雅温得中学 学习。};
\end{tikzpicture}
\legende{Schéma de la phrase type : [时间] + 我在 + [地点] + [动词]}
\end{popfigure}

\courssub{Rubrique expérience : 经验}

\begin{exemplebox}[La phrase d'expérience]
\zh{2023年到2024年，我在杜阿拉公司实习。}\\
\emph{Èr-líng-èr-sān nián dào èr-líng-èr-sì nián, wǒ zài Dùālā gōngsī shíxí.}\\
« De 2023 à 2024, j'ai fait un stage à l'entreprise de Douala. »
\end{exemplebox}

Même structure que la formation : \textbf{Temps + Sujet + 在 + Lieu + Verbe}. Le verbe \zh{实习} (shíxí, « faire un stage ») s'utilise sans complément d'objet direct ici.

\courssub{Rubrique compétences : 技能}

\begin{exemplebox}[La phrase de compétences]
\zh{我会说中文、法语和英语。}\\
\emph{Wǒ huì shuō Zhōngwén, Fǎyǔ hé Yīngyǔ.}\\
« Je sais parler chinois, français et anglais. »
\end{exemplebox}

Le verbe modal \zh{会} (huì) exprime la \textbf{compétence acquise} (« savoir faire ») : \zh{会说} « savoir parler », \zh{会用电脑} « savoir utiliser l'ordinateur », \zh{会开车} « savoir conduire ». Notez l'énumération : \zh{中文、法语和英语} — la virgule chinoise \zh{、} sépare les éléments, et \zh{和} (hé, « et ») ne se place qu'entre les \textbf{deux derniers} éléments.

\courssub{Rubrique loisirs : 爱好}

\begin{exemplebox}[La phrase de loisirs]
\zh{我喜欢足球和音乐。}\\
\emph{Wǒ xǐhuan zúqiú hé yīnyuè.}\\
« J'aime le football et la musique. »
\end{exemplebox}

\courssub{La ponctuation chinoise : la virgule d'énumération 、}

\begin{propriete}[Le dunhao 、]
Dans une énumération en chinois, on utilise la virgule \zh{、} (appelée \emph{dunhao}) et \textbf{non} la virgule ordinaire \zh{，}. Le \zh{、} ne sert qu'à séparer des mots ou groupes de mots de même nature dans une liste : \zh{中文、法语和英语}. La virgule \zh{，}, elle, sépare des propositions ou des phrases : \zh{首先，我介绍我的教育。}
\end{propriete}

\begin{attention}[Erreur fréquente des francophones]
Ne traduisez pas la virgule française par \zh{，} partout. « Je parle chinois, français et anglais » se traduit \zh{我会说中文、法语和英语。} et non \zh{我会说中文，法语，和英语。} Autre piège : ne mettez jamais \zh{和} entre tous les éléments (\zh{中文和法语和英语} est lourd et incorrect dans une liste longue) ; \zh{和} ne vient qu'avant le dernier élément.
\end{attention}

\courssub{Les connecteurs utiles}

Pour transformer une suite de phrases en un texte fluide, le chinois utilise des \textbf{connecteurs}. En voici les cinq indispensables du CV.

\begin{center}
\begin{tabularx}{\linewidth}{|X|X|X|}
\hline
\rowcolor{popblueL} Connecteur & Exemple (caractères + pinyin) & Traduction \\
\hline
首先 shǒuxiān « d'abord » & 首先，我介绍我的教育。Shǒuxiān, wǒ jièshào wǒ de jiàoyù. & D'abord, je présente ma formation. \\
\hline
然后 ránhòu « ensuite » & 然后，我说我的技能。Ránhòu, wǒ shuō wǒ de jìnéng. & Ensuite, je parle de mes compétences. \\
\hline
也 yě « aussi » & 我也喜欢足球。Wǒ yě xǐhuan zúqiú. & J'aime aussi le football. \\
\hline
而且 érqiě « de plus » & 我会说中文，而且我会用电脑。Wǒ huì shuō Zhōngwén, érqiě wǒ huì yòng diànnǎo. & Je sais parler chinois, et de plus je sais utiliser l'ordinateur. \\
\hline
因为…所以… yīnwèi… suǒyǐ… « parce que… donc… » & 因为我喜欢中文，所以我学习中文。Yīnwèi wǒ xǐhuan Zhōngwén, suǒyǐ wǒ xuéxí Zhōngwén. & Parce que j'aime le chinois, j'apprends le chinois. \\
\hline
\end{tabularx}
\end{center}

\begin{propriete}[Place et emploi des connecteurs]
\begin{itemize}
\item \zh{首先} et \zh{然后} se placent en \textbf{tête de phrase}, suivis d'une virgule \zh{，}. On peut enchaîner : \zh{首先…然后…} « d'abord… ensuite… ».
\item \zh{也} se place \textbf{après le sujet, avant le verbe} : \zh{我也喜欢…} (jamais en début de phrase comme le « aussi » français).
\item \zh{而且} relie deux propositions et se place après la virgule : \zh{…，而且…}.
\item \zh{因为…所以…} : en chinois, les deux mots se combinent dans la même phrase, contrairement au français où « parce que » et « donc » ne s'emploient pas ensemble. Dire \zh{因为…} seul est possible ; dire \zh{所以…} seul aussi ; mais la paire complète est très naturelle.
\end{itemize}
\end{propriete}

\begin{attention}[Ne pas conjuguer !]
En chinois, le verbe \textbf{ne se conjugue pas} : \zh{我在雅温得中学学习} reste identique quelle que soit la personne — \zh{你学习} (tu étudies), \zh{他学习} (il étudie), \zh{我们学习} (nous étudions). N'ajoutez rien au verbe, ne changez rien. C'est le sujet qui porte la personne. De même, le temps est indiqué par le complément de temps (\zh{2020年到2024年}), pas par une terminaison verbale.
\end{attention}

\courssub{Organisation générale d'un CV chinois}

\begin{popfigure}
\begin{tikzpicture}[node distance=0.45cm,
  etape/.style={rectangle, rounded corners=3pt, draw=popblue, fill=popblueL, minimum width=4.6cm, minimum height=0.85cm, align=center, font=\small},
  fleche/.style={-{Stealth[length=2.5mm]}, thick, draw=popdark}]
\node[etape, fill=popgoldL, draw=popgold, align=center] (tete){姓名、性别、出生日期、地址、电话、邮箱\\ en-tête};
\node[etape, below=of tete, align=center] (jiao){教育\\ formation};
\node[etape, below=of jiao, align=center] (jing){经验\\ expérience};
\node[etape, below=of jing, align=center] (ji){技能\\ compétences};
\node[etape, below=of ji, align=center] (ai){爱好\\ loisirs};
\draw[fleche] (tete) -- (jiao);
\draw[fleche] (jiao) -- (jing);
\draw[fleche] (jing) -- (ji);
\draw[fleche] (ji) -- (ai);
\end{tikzpicture}
\legende{Organigramme de la structure d'un CV : en-tête → 教育 → 经验 → 技能 → 爱好}
\end{popfigure}

\courssub{Commentaire ligne par ligne du modèle}

Reprenons le texte de 李明 phrase par phrase.

\begin{enumerate}
\item \zh{我叫李明。} — Formule d'identité : \zh{叫} (jiào) « s'appeler ». Point final chinois \zh{。}.
\item \zh{首先，我介绍我的教育：…} — Le connecteur \zh{首先} annonce le plan ; les deux-points \zh{：} introduisent le détail de la rubrique.
\item \zh{2020年到2024年，我在雅温得中学学习。} — Phrase type : temps (\zh{到} « jusqu'à ») + sujet + \zh{在} + lieu + verbe. La virgule \zh{，} sépare le complément de temps du reste.
\item \zh{然后，我说我的技能：…} — Deuxième étape du plan avec \zh{然后}.
\item \zh{我会说中文、法语和英语，而且我会用电脑。} — Énumération avec \zh{、} ; \zh{和} avant le dernier élément ; \zh{而且} ajoute une compétence supplémentaire.
\item \zh{我也喜欢足球和音乐。} — \zh{也} placé après le sujet : la rubrique loisirs conclut le texte.
\end{enumerate}

\begin{methode}[Rédiger son propre mini-CV]
\begin{enumerate}
\item \textbf{Recopier la structure du modèle} : en-tête, puis \zh{首先} + formation, \zh{然后} + expérience (si applicable), \zh{然后} + compétences, \zh{也} + loisirs.
\item \textbf{Remplacer les informations personnelles} : ton nom, tes dates, ton école (\zh{我在……中学学习}), ton stage (\zh{我在……公司实习}), tes langues, tes logiciels, tes loisirs.
\item \textbf{Vérifier traits et tons} : relis chaque caractère complexe (\zh{教育}, \zh{技能}, \zh{经验}, \zh{简历}) trait par trait, contrôle la ponctuation (\zh{、} pour les listes, \zh{，} entre les phrases, \zh{。} à la fin, \zh{：} après les rubriques) et le pinyin de chaque mot.
\end{enumerate}
\end{methode}

\begin{exoresolu}[Thème : Traduire en chinois]
Énoncé : Traduisez ce mini-CV en chinois.
« Je m'appelle Marie. D'abord, ma formation : de 2021 à 2024, j'ai étudié au lycée de Douala. Ensuite, mon expérience : en 2024, j'ai fait un stage à l'entreprise de Yaoundé. Mes compétences : je sais parler français et anglais, et je sais utiliser l'ordinateur. J'aime aussi la musique. »
\tcblower
Solution :\\
\zh{我叫玛丽。首先，我介绍我的教育：2021年到2024年，我在杜阿拉中学学习。然后，我介绍我的经验：2024年，我在雅温得公司实习。然后，我说我的技能：我会说法语和英语，而且我会用电脑。我也喜欢音乐。}\\
\emph{Wǒ jiào Mǎlì. Shǒuxiān, wǒ jièshào wǒ de jiàoyù: èr-líng-èr-yī nián dào èr-líng-èr-sì nián, wǒ zài Dùālā zhōngxué xuéxí. Ránhòu, wǒ jièshào wǒ de jīngyàn: èr-líng-èr-sì nián, wǒ zài Yǎwēndé gōngsī shíxí. Ránhòu, wǒ shuō wǒ de jìnéng: wǒ huì shuō Fǎyǔ hé Yīngyǔ, érqiě wǒ huì yòng diànnǎo. Wǒ yě xǐhuan yīnyuè.}\\
Points de méthode : (1) \zh{首先} ouvre la première rubrique, \zh{然后} les suivantes ; (2) les deux-points \zh{：} annoncent le contenu ; (3) structure \textbf{Temps + Sujet + 在 + Lieu + Verbe} respectée pour formation et expérience ; (4) \zh{和} pour deux langues, \zh{而且} pour ajouter une compétence distincte ; (5) \zh{也} avant le verbe pour le loisir.
\end{exoresolu}

\begin{exoresolu}[Version : Traduire en français]
Énoncé : Traduisez ce texte en français.
\zh{我叫保罗。首先，2019年到2025年，我在布埃亚中学学习。然后，我会说法语、中文和英语，而且我会用电脑。因为我喜欢中文，所以我努力学习。我也喜欢足球。}
\tcblower
Solution :\\
« Je m'appelle Paul. D'abord, de 2019 à 2025, j'ai étudié au lycée de Buea. Ensuite, je sais parler français, chinois et anglais, et de plus je sais utiliser l'ordinateur. Parce que j'aime le chinois, j'étudie avec ardeur. J'aime aussi le football. »\\
Notes : \zh{布埃亚} (Bù'āiyà) = Buea. \zh{努力学习} (nǔlì xuéxí) = étudier avec ardeur / sérieusement. Structure \zh{因为…所以…} rendue par « Parce que…, … » (le « donc » est absorbé par la structure française).
\end{exoresolu}

\begin{exercice}[À toi d'écrire]
Rédige ton propre mini-CV en 5 à 7 phrases, en suivant le modèle de 李明 : présentation (\zh{我叫…}), formation avec dates et lieu, expérience (stage ou job d'été) si tu en as une, compétences (langues, informatique) et loisirs. Utilise au moins quatre connecteurs parmi \zh{首先}, \zh{然后}, \zh{也}, \zh{而且}, \zh{因为…所以…}. Soigne la ponctuation : \zh{、} dans les énumérations, \zh{。} à la fin de chaque phrase, \zh{：} après les titres de rubriques.
\end{exercice}

\begin{corrige}[Exercice « À toi d'écrire »]
Exemple de production attendue (les informations personnelles varient) :\\
\zh{我叫阿马杜。首先，我介绍我的教育：2019年到2025年，我在雅温得中学学习。然后，我介绍我的经验：2024年，我在杜阿拉公司实习。然后，我说我的技能：我会说法语、英语和中文，而且我会用电脑。因为我很喜欢语言，所以我努力学习中文。我也喜欢足球和音乐。}\\
\emph{Wǒ jiào Āmǎdù. Shǒuxiān, wǒ jièshào wǒ de jiàoyù: èr-líng-yī-jiǔ nián dào èr-líng-èr-wǔ nián, wǒ zài Yǎwēndé zhōngxué xuéxí. Ránhòu, wǒ jièshào wǒ de jīngyàn: èr-líng-èr-sì nián, wǒ zài Dùālā gōngsī shíxí. Ránhòu, wǒ shuō wǒ de jìnéng: wǒ huì shuō Fǎyǔ, Yīngyǔ hé Zhōngwén, érqiě wǒ huì yòng diànnǎo. Yīnwèi wǒ hěn xǐhuan yǔyán, suǒyǐ wǒ nǔlì xuéxí Zhōngwén. Wǒ yě xǐhuan zúqiú hé yīnyuè.}\\
« Je m'appelle Amadou. D'abord, je présente ma formation : de 2019 à 2025, j'ai étudié au lycée de Yaoundé. Ensuite, je présente mon expérience : en 2024, j'ai fait un stage à l'entreprise de Douala. Ensuite, je parle de mes compétences : je sais parler français, anglais et chinois, et de plus je sais utiliser l'ordinateur. Parce que j'aime beaucoup les langues, j'étudie le chinois avec ardeur. J'aime aussi le football et la musique. »\\
Critères de réussite : structure en-tête → formation → expérience → compétences → loisirs respectée ; au moins quatre connecteurs correctement placés ; \zh{在} + lieu avant le verbe ; \zh{、} dans la liste des langues (3 éléments) ; verbes non conjugués ; \zh{因为…所以…} bien formé.
\end{corrige}

\begin{aretenir}[L'essentiel de la section]
\begin{itemize}
\item Un CV chinois suit l'ordre : en-tête (\zh{姓名、性别、出生日期、地址、电话、邮箱}) → \zh{教育} → \zh{经验} → \zh{技能} → \zh{爱好}.
\item Phrase type : \textbf{[Temps] + 我在 + [Lieu] + Verbe} ; le lieu avec \zh{在} précède le verbe.
\item Énumération : virgule \zh{、} entre les éléments, \zh{和} seulement avant le dernier.
\item Connecteurs : \zh{首先} « d'abord », \zh{然后} « ensuite » en tête de phrase ; \zh{也} « aussi » après le sujet ; \zh{而且} « de plus » ; \zh{因为…所以…} « parce que… donc… ».
\item Le verbe chinois ne se conjugue pas : c'est le sujet et le complément de temps qui portent l'information.
\item Ponctuation spécifique : \zh{：} après rubrique, \zh{、} liste, \zh{，} proposition, \zh{。} phrase.
\end{itemize}
\end{aretenir}

\coursec{Leçon 40 — Expression écrite : caractères complexes liés à la rédaction d'un CV}

\courssub{Découverte : le CV chinois et ses rubriques}

\begin{textebox}[Dialogue : Préparer son CV pour un stage à Douala]
\zh{— 小李，你的简历准备好了吗？} \\
\textit{— Xiǎo Lǐ, nǐ de jiǎnlì zhǔnbèi hǎo le ma?} \\
« Petit Li, ton CV est-il prêt ? » \\

\zh{— 还差一点。我在写“教育背景”和“工作经验”这两部分。} \\
\textit{— Hái chà yīdiǎn. Wǒ zài xiě “jiàoyù bèijǐng” hé “gōngzuò jīngyàn” zhè liǎng bùfen.} \\
« Il manque encore un peu. J'écris les parties "Formation" et "Expérience professionnelle". » \\

\zh{— 别忘了“技能特长”和“爱好”。中文简历通常要写“自我评价”。} \\
\textit{— Bié wàngle “jìnéng tècháng” hé “àihào”. Zhōngwén jiǎnlì tōngcháng yào xiě “zìwǒ píngjià”.} \\
« N'oublie pas "Compétences" et "Loisirs". Un CV chinois comporte généralement une "Auto-évaluation". » \\
\zh{— 好的，我马上加上。谢谢王经理指点！} \\
\textit{— Hǎo de, wǒ mǎshàng jiā shàng. Xièxie Wáng jīnglǐ zhǐdiǎn!} \\
« D'accord, j'ajoute tout de suite. Merci pour vos conseils, M. Wang ! » \\
\end{textebox}

\begin{center}
\begin{tabularx}{\linewidth}{|X|X|X|X|}
\hline
\rowcolor{popblueL} \textbf{Caractères} & \textbf{Pinyin} & \textbf{Nature} & \textbf{Traduction} \\
\hline
简历 & jiǎnlì & n. & CV (curriculum vitæ) \\
\hline
教育背景 & jiàoyù bèijǐng & n. & Formation, parcours scolaire \\
\hline
工作经验 & gōngzuò jīngyàn & n. & Expérience professionnelle \\
\hline
实习经历 & shíxí jīnglì & n. & Expérience de stage \\
\hline
技能特长 & jìnéng tècháng & n. & Compétences et aptitudes particulières \\
\hline
爱好 & àihào & n. & Loisirs, centres d'intérêt \\
\hline
自我评价 & zìwǒ píngjià & n. & Auto-évaluation \\
\hline
个人信息 & gèrén xìnxī & n. & Informations personnelles \\
\hline
求职意向 & qiúzhí yìxiàng & n. & Poste visé / Objectif professionnel \\
\hline
联系方式 & liánxì fāngshì & n. & Coordonnées \\
\hline
\end{tabularx}
\end{center}

\begin{aretenir}[Différences CV Cameroun / Chine]
\begin{itemize}
    \item \textbf{Photo :} Obligatoire en Chine (coin sup. droit), optionnelle au Cameroun (souvent déconseillée pour éviter discrimination).
    \item \textbf{Informations personnelles :} En Chine : âge, sexe, ethnie (\zh{汉族} Hànzú), état civil, \zh{户口} (hùkǒu, registre de résidence), parti politique. Au Cameroun : nom, prénom, contact, nationalité, parfois âge.
    \item \textbf{Auto-évaluation (\zh{自我评价}) :} Incontournable en Chine (qualités, défauts, motivation). Rare au Cameroun (remplacé par la lettre de motivation).
    \item \textbf{Longueur :} 1 page A4 dans les deux pays pour un jeune diplômé.
\end{itemize}
\end{aretenir}

\courssub{Écriture des caractères complexes (1) : 简历, 经验}

\begin{propriete}[Analyse du caractère \zh{简} (jiǎn, simple, simplifié)]
\textbf{Radical (clé) :} \zh{竹} (zhú, bambou) — en haut. \textbf{Structure :} Haut-bas. \textbf{Nombre de traits :} 13.
\textbf{Ordre des traits :} 1. Bambou (6 traits : horizontales, verticales, crochet) $\rightarrow$ 2. \zh{间} (jiān, 7 traits : porte \zh{门} + soleil \zh{日}).
\textbf{Étymologie :} Bâtonnets de bambou sélectionnés (simples) pour écrire $\rightarrow$ simple, choisir $\rightarrow$ résumé (\zh{简历} = résumé de l'historique).
\textbf{Caractères proches :} \zh{篮} (lán, panier, radical bambou), \zh{策} (cè, stratagème, radical bambou).
\end{propriete}

\begin{propriete}[Analyse du caractère \zh{历} (lì, calendrier, expérience, passer par)]
\textbf{Radical (clé) :} \zh{厂} (chǎng, falaise/abri) — en haut à gauche. \textbf{Structure :} Semi-encerclement haut-gauche. \textbf{Nombre de traits :} 8.
\textbf{Ordre des traits :} 1. \zh{厂} (2 traits : horizontale, verticale crochet) $\rightarrow$ 2. \zh{止} (zhǐ, 4 traits) $\rightarrow$ 3. Trait horizontal final de liaison.
\textbf{Sens originel :} Calendrier, almanach $\rightarrow$ passer par (le temps), expérience.
\textbf{Composés :} \zh{历史} (lìshǐ, histoire), \zh{经历} (jīnglì, expérience vécue), \zh{日历} (rìlì, calendrier).
\end{propriete}

\begin{propriete}[Analyse du caractère \zh{经} (jīng, trame, passer par, gérer)]
\textbf{Radical (clé) :} \zh{糸} (mì, fil de soie) — à gauche. \textbf{Structure :} Gauche-droite. \textbf{Nombre de traits :} 13.
\textbf{Ordre des traits :} 1. \zh{糸} (6 traits : points, crochets, horizontales) $\rightarrow$ 2. \zh{巠} (jīng, 7 traits : verticale, horizontales, verticale centrale).
\textbf{Sens :} Fils de chaîne (verticaux) sur un métier à tisser $\rightarrow$ constant, régulier $\rightarrow$ passer par, gérer, sutra bouddhique.
\textbf{Caractères proches :} \zh{径} (jìng, chemin, radical \zh{辶}), \zh{茎} (jīng, tige, radical \zh{艹}).
\end{propriete}

\begin{propriete}[Analyse du caractère \zh{验} (yàn, vérifier, expérience)]
\textbf{Radical (clé) :} \zh{马} (mǎ, cheval) — à droite. \textbf{Structure :} Gauche-droite. \textbf{Nombre de traits :} 10.
\textbf{Partie gauche :} \zh{ㄧ} + \zh{丨} + \zh{豕} (shǐ, porc) simplifié $\rightarrow$ idée d'inspection.
\textbf{Ordre des traits :} Partie gauche (5 traits) $\rightarrow$ \zh{马} (3 traits : horizontales, verticale, points) $\rightarrow$ trait final.
\textbf{Sens :} À l'origine, inspecter un cheval avant achat $\rightarrow$ vérifier, tester $\rightarrow$ expérience (résultat de la vérification).
\textbf{Composés :} \zh{实验} (shíyàn, expérience scientifique), \zh{检验} (jiǎnyàn, inspection), \zh{经验} (jīngyàn, expérience professionnelle/vie).
\end{propriete}

\begin{experience}[Atelier écriture : Caractères du CV]
\textbf{Consigne :} Sur papier quadrillé (grille 米字格), écrire 5 fois chaque caractère ci-dessous en respectant l'ordre des traits et les proportions. Identifier le radical et donner un mot composé.
\begin{enumerate}
    \item \zh{简} (Radical : \zh{竹}) $\rightarrow$ \zh{简单} (jiǎndān, simple)
    \item \zh{历} (Radical : \zh{厂}) $\rightarrow$ \zh{历史} (lìshǐ, histoire)
    \item \zh{经} (Radical : \zh{糸}) $\rightarrow$ \zh{经济} (jīngjì, économie)
    \item \zh{验} (Radical : \zh{马}) $\rightarrow$ \zh{实验} (shíyàn, expérience)
\end{enumerate}
\textbf{Auto-correction :} Vérifier la hauteur du radical \zh{竹} (occupe 1/3 haut), l'ouverture du \zh{厂}, l'alignement du \zh{糸} à gauche.
\end{experience}

\courssub{Écriture des caractères complexes (2) : 教育, 技能, 聘}

\begin{propriete}[Analyse du caractère \zh{教} (jiāo, enseigner)]
\textbf{Radical :} \zh{攵} (pū, frapper/taper) — à droite (forme compressée de \zh{攴}). \textbf{Structure :} Gauche-droite. \textbf{Traits :} 11.
\textbf{Ordre :} Gauche \zh{孝} (xiào, piété filiale : \zh{耂} + \zh{子}) $\rightarrow$ Droite \zh{攵} (4 traits : point, horizontale, point, trait oblique long).
\textbf{Logique :} \zh{孝} (éduquer avec piété) + action de guider (\zh{攵}) = enseigner.
\textbf{Composés :} \zh{教育} (jiàoyù), \zh{教授} (jiàoshòu, professeur), \zh{教室} (jiàoshì, salle de classe).
\end{propriete}

\begin{propriete}[Analyse du caractère \zh{育} (yù, élever, éduquer)]
\textbf{Radical :} \zh{月} (yuè, chair/corps) — en bas (forme de \zh{肉} ròu). \textbf{Structure :} Haut-bas. \textbf{Traits :} 8.
\textbf{Ordre :} Haut \zh{立} (lì, debout, 5 traits) $\rightarrow$ Bas \zh{月} (4 traits : deux verticales, deux horizontales intérieures).
\textbf{Logique :} Un enfant (\zh{立}) qui grandit dans le corps (\zh{月} $\leftarrow$ \zh{肉}) $\rightarrow$ élever, éduquer.
\textbf{Attention :} Ne pas confondre \zh{月} (lune) et \zh{肉} (viande/corps) : en bas de caractère, c'est presque toujours \zh{肉} (ex: \zh{肚} dù, \zh{脚} jiǎo, \zh{育} yù).
\end{propriete}

\begin{propriete}[Analyse du caractère \zh{技} (jì, technique, habileté)]
\textbf{Radical :} \zh{扌} (shǒu, main) — à gauche. \textbf{Structure :} Gauche-droite. \textbf{Traits :} 7.
\textbf{Ordre :} \zh{扌} (3 traits) $\rightarrow$ \zh{支} (zhī, 4 traits : horizontale, verticale, horizontale, crochet).
\textbf{Logique :} Main qui soutient/manipule (\zh{支}) avec adresse $\rightarrow$ technique, art, métier.
\textbf{Caractères proches :} \zh{支} (zhī, soutenir, branche), \zh{枝} (zhī, branche d'arbre, radical \zh{木}), \zh{技} (jì, radical \zh{扌}).
\end{propriete}

\begin{propriete}[Analyse du caractère \zh{能} (néng, capacité, pouvoir)]
\textbf{Radical :} \zh{月} (yuè $\leftarrow$ \zh{肉}) — en bas. \textbf{Structure :} Haut-bas. \textbf{Traits :} 10.
\textbf{Ordre :} Haut (partie complexe : \zh{匕} + \zh{厶} + traits) $\rightarrow$ Bas \zh{月} (4 traits).
\textbf{Mnémotechnique :} Un homme fort (\zh{匕} assis, \zh{厶} privé) avec de la viande (\zh{月}) a de l'énergie.
\textbf{Attention :} \zh{会} (huì, savoir/appris) vs \zh{能} (néng, capacité innée/physique/autorisation). CV : \zh{技能} (compétences techniques), \zh{能力} (nénglì, capacité générale).
\end{propriete}

\begin{propriete}[Analyse du caractère \zh{聘} (pìn, engager, inviter)]
\textbf{Radical :} \zh{耳} (ěr, oreille) — à gauche. \textbf{Structure :} Gauche-droite. \textbf{Traits :} 13.
\textbf{Partie droite :} \zh{并} (bìng, ensemble) + \zh{寸} (cùn, pouce/mesure/loi).
\textbf{Ordre :} \zh{耳} (6 traits) $\rightarrow$ \zh{并} (6 traits) $\rightarrow$ \zh{寸} (3 traits).
\textbf{Logique :} Écouter (\zh{耳}) les termes réunis (\zh{并}) selon la règle (\zh{寸}) $\rightarrow$ contrat d'embauche.
\textbf{Composés :} \zh{聘用} (pìnyòng, embaucher), \zh{聘书} (pìnshū, lettre d'embauche), \zh{特聘} (tèpìn, recrutement spécial).
\end{propriete}

\begin{exemplebox}[Exercices d'écriture guidée : Radicaux et sens]
Complétez le tableau en reliant le caractère à son radical et au sens du radical.
\begin{center}
\begin{tabularx}{\linewidth}{|X|X|X|X|}
\hline
\rowcolor{popblueL} \textbf{Caractère} & \textbf{Radical (Clé)} & \textbf{Sens du radical} & \textbf{Indice mnémotechnique} \\
\hline
\zh{教} & \zh{攵} (pū) & Action de taper/guider & Éduquer en guidant (comme un maître) \\
\hline
\zh{育} & \zh{月} (\zh{肉} ròu) & Corps, chair & L'enfant grandit dans le corps \\
\hline
\zh{技} & \zh{扌} (shǒu) & Main, action manuelle & Habileté de la main \\
\hline
\zh{能} & \zh{月} (\zh{肉} ròu) & Corps, viande & Force physique, énergie \\
\hline
\zh{聘} & \zh{耳} (ěr) & Oreille, écouter & Écouter les conditions du contrat \\
\hline
\end{tabularx}
\end{center}
\end{exemplebox}

\courssub{Modèle de CV commenté, structure et connecteurs}

\begin{textebox}[Modèle de CV : Amina Bello (Terminale A, Yaoundé)]
\zh{\textbf{个人简历} (Gèrén Jiǎnlì)} \\

\zh{\textbf{个人信息} (Gèrén Xìnxī)} \\
\zh{姓名：阿明娜 \quad 性别：女 \quad 民族：汉族/非洲裔} \\
\zh{出生年月：2006年5月 \quad 户籍：喀麦隆 雅温得} \\
\zh{电话：+237 6XX XX XX XX \quad 邮箱：amina.bello@email.cm} \\
\zh{求职意向：实习生 / 行政助理} \\

\zh{\textbf{教育背景} (Jiàoyù Bèijǐng)} \\
\zh{2021年9月 -- 2024年6月 \quad 雅温得第一高中 (Yaoundé Lycée N°1)} \\
\zh{专业：文学系列 (Série A) \quad 核心课程：法语、中文、历史、地理、哲学} \\
\zh{主修成绩：中文 HSK 3 级 (240/300)、法语 B2、英语 B1} \\

\zh{\textbf{实习经历} (Shíxí Jīnglì)} \\
\zh{2023年7月 -- 2023年8月 \quad 雅温得中国文化中心 (Centre Culturel Chinois de Yaoundé)} \\
\zh{职位：行政实习生} \\
\zh{主要职责：} \\
\zh{1. 协助组织汉语角活动，接待访客 50+ 人次。} \\
\zh{2. 翻译中法文活动宣传资料，管理微信公众号排版。} \\
\zh{3. 整理图书档案，协助 HSK 考试监考工作。} \\

\zh{\textbf{技能特长} (Jìnéng Tècháng)} \\
\zh{语言：中文 (HSK 3, 熟练口语)、法语 (母语级)、英语 (流利)、富尔贝语 (母语)} \\
\zh{办公软件：熟练掌握 Word、Excel、PPT、Canva} \\
\zh{其他：持有驾照 (B 类)、擅长摄影与短视频剪辑 (CapCut)} \\

\zh{\textbf{爱好与特长} (Àihào yǔ Tècháng)} \\
\zh{足球 (前锋)、阅读非洲文学、非洲鼓演奏、志愿者服务 (红十字会青年团)} \\

\zh{\textbf{自我评价} (Zìwǒ Píngjià)} \\
\zh{性格开朗，责任心强，具备良好的跨文化沟通能力。} \\
\zh{作为喀麦隆学生，熟悉中非文化差异，渴望在中资企业实习，为中喀合作贡献力量。} \\
\zh{抗压能力强，适应快节奏工作环境，服从安排，执行力强。} \\
\end{textebox}

\begin{propriete}[Structure type d'un CV chinois (标准简历结构)]
\begin{enumerate}
    \item \zh{个人信息} (Gèrén Xìnxī) : Nom, sexe, naissance, contact, \textbf{photo}.
    \item \zh{求职意向} (Qiúzhí Yìxiàng) : Poste visé, ville, salaire attendu (optionnel).
    \item \zh{教育背景} (Jiàoyù Bèijǐng) : \textbf{Ordre antéchronologique} (récent $\rightarrow$ ancien). École, dates, spécialité, mentions, \textbf{HSK}.
    \item \zh{工作/实习经历} (Gōngzuò/Shíxí Jīnglì) : Entreprise, dates, poste, \textbf{réalisations chiffrées} (verbes d'action).
    \item \zh{技能特长} (Jìnéng Tècháng) : Langues (niveau certifié), Informatique (logiciels), Permis, Certificats.
    \item \zh{爱好特长} (Àihào Tècháng) : Sports, arts, bénévolat (montrer soft skills).
    \item \zh{自我评价} (Zìwǒ Píngjià) : 3-5 phrases : personnalité + atouts + motivation + disponibilité.
\end{enumerate}
\end{propriete}

\begin{propriete}[Connecteurs logiques pour CV et rédaction formelle]
\begin{center}
\begin{tabularx}{\linewidth}{|X|X|X|}
\hline
\rowcolor{popblueL} \textbf{Chinois} & \textbf{Pinyin} & \textbf{Français / Emploi} \\
\hline
首先 & shǒuxiān & Premièrement / Tout d'abord (début énumération) \\
\hline
其次 & qícì & Deuxièmement / En second lieu \\
\hline
此外 & cǐwài & De plus / En outre / Par ailleurs \\
\hline
再者 & zàishuō & Qui plus est / En plus (argument supplémentaire) \\
\hline
最后 & zuìhòu & Enfin / En dernier lieu \\
\hline
因此 & yīncǐ & Par conséquent / C'est pourquoi (conséquence) \\
\hline
不过 & bùguò & Cependant / Toutefois (concession, oral/écrit courant) \\
\hline
但是 & dànshì & Mais / Cependant (opposition forte) \\
\hline
例如 & lìrú & Par exemple / Notamment (illustration) \\
\hline
特别是 & tèbié shì & Surtout / Plus particulièrement (focus) \\
\hline
\end{tabularx}
\end{center}
\end{propriete}

\begin{exemplebox}[Utilisation des connecteurs dans \zh{自我评价}]
\zh{我认为我适合这个岗位。\textbf{首先}，我通过了HSK 3级考试，中文沟通无障碍。\textbf{其次}，我在中国文化中心实习过，熟悉中方办公流程。\textbf{此外}，我性格开朗，抗压能力强。\textbf{最后}，我非常渴望为中喀合作贡献力量。}
\textit{Wǒ rènwéi wǒ shìhé zhège gǎngwèi. \textbf{Shǒuxiān}, wǒ tōngguò le HSK 3 jí kǎoshì, Zhōngwén gōutōng wú zhàng'ài. \textbf{Qícì}, wǒ zài Zhōngguó Wénhuà Zhōngxīn shíxí guò, shúxī Zhōngfāng bàngōng liúchéng. \textbf{Cǐwài}, wǒ xìnggé kāilǎng, kàngyā nénglì qiáng. \textbf{Zuìhòu}, wǒ fēicháng kěwàng wèi Zhōng-Kā hézuò gòngxiàn lìliàng.}
« Je pense être adaptée à ce poste. \textbf{Premièrement}, j'ai réussi le HSK 3, communication chinoise fluide. \textbf{Deuxièmement}, j'ai fait un stage au Centre Culturel Chinois, je connais les procédures administratives chinoises. \textbf{De plus}, je suis extravertie et résistante au stress. \textbf{Enfin}, je souhaite ardemment contribuer à la coopération sino-camerounaise. »
\end{exemplebox}

\courssub{Traduction : thème et version}

\begin{methode}[Thème : 5 étapes pour passer du français au chinois]
\textbf{Étape 1 : Identifier le noyau de la phrase.} Repérez le \cle{sujet}, le \cle{verbe} principal et les \cle{compléments} (temps, lieu, manière, objet).\\
\textbf{Étape 2 : Souligner les indicateurs temps/lieu.} En français, ils sont souvent en fin de phrase (ex: \emph{« De 2021 à 2024 »}, \emph{« au lycée de Douala »}). En chinois, ils \textbf{doivent} monter devant le verbe.\\
\textbf{Étape 3 : Appliquer l'ordre canonique chinois :} \cle{Sujet + Temps + Lieu + Verbe (+ Objet)}.\\
\textbf{Étape 4 : Choisir les bons mots-outils et classificateurs.} Vérifiez : \zh{了} (action achevée), \zh{过} (expérience vécue), \zh{会} (compétence acquise), et le classificateur correct (ex: \zh{种} pour les langues, \zh{家} pour entreprises, \zh{个} générique).\\
\textbf{Étape 5 : Relire à voix haute en pinyin.} Contrôlez les tons, l'ordre des mots et la ponctuation chinoise (， 。).
\end{methode}

\begin{propriete}[Règle d'or : L'ordre des compléments Temps / Lieu / Manière]
En français, l'ordre est souvent \textbf{Sujet -- Verbe -- Complément de Lieu -- Complément de Temps} (ex: \emph{J'ai étudié à Douala en 2021}).
En chinois, l'ordre est immuable : \textbf{Sujet -- Temps -- Lieu -- Manière -- Verbe}.
\zh{我 2021年 在杜阿拉 用功 学习。} (Wǒ 2021 nián zài Dùyālā yònggōng xuéxí.)
\emph{Point de vigilance :} Le mot \zh{在} (zài) introduit le lieu et se place \textbf{avant} le verbe, après le temps. Le complément de manière (\zh{用功}) se place après le lieu.
\end{propriete}

\begin{popfigure}
\legende{Schéma comparatif de l'ordre des mots : Français vs Chinois}
\begin{tikzpicture}[
    node distance=0.7cm and 1.0cm,
    box/.style={rectangle, draw=popblue, thick, fill=popblueL, rounded corners, minimum height=1.1cm, text width=2.0cm, align=center, font=\small},
    frbox/.style={rectangle, draw=poporange, thick, fill=poporangeL, rounded corners, minimum height=1.1cm, text width=2.0cm, align=center, font=\small},
    arrow/.style={-Stealth, thick, popdark}
]
% French row
\node[frbox, align=center] (fr_s){Sujet\\\emph{Je}};
\node[frbox, right=of fr_s, align=center] (fr_v){Verbe\\\emph{ai étudié}};
\node[frbox, right=of fr_v, align=center] (fr_m){Manière\\\emph{assidûment}};
\node[frbox, right=of fr_m, align=center] (fr_l){Lieu\\\emph{à Douala}};
\node[frbox, right=of fr_l, align=center] (fr_t){Temps\\\emph{en 2021}};

% Chinese row
\node[box, below=1.3cm of fr_s, align=center] (zh_s){Sujet\\\zh{我} (Wǒ)};
\node[box, right=of zh_s, align=center] (zh_t){Temps\\\zh{2021年} (2021 nián)};
\node[box, right=of zh_t, align=center] (zh_l){Lieu\\\zh{在杜阿拉} (zài Dùyālā)};
\node[box, right=of zh_l, align=center] (zh_m){Manière\\\zh{用功} (yònggōng)};
\node[box, right=of zh_m, align=center] (zh_v){Verbe\\\zh{学习} (xuéxí)};

% Arrows
\draw[arrow] (fr_s) -- (fr_v);
\draw[arrow] (fr_v) -- (fr_m);
\draw[arrow] (fr_m) -- (fr_l);
\draw[arrow] (fr_l) -- (fr_t);

\draw[arrow] (zh_s) -- (zh_t);
\draw[arrow] (zh_t) -- (zh_l);
\draw[arrow] (zh_l) -- (zh_m);
\draw[arrow] (zh_m) -- (zh_v);

% Cross mapping
\draw[dashed, popmuted, thick] (fr_s.south) -- (zh_s.north);
\draw[dashed, popmuted, thick] (fr_t.south) .. controls +(down:1cm) and +(up:1cm) .. (zh_t.north) node[midway, left, font=\tiny, popmuted] {Déplacement};
\draw[dashed, popmuted, thick] (fr_l.south) .. controls +(down:1cm) and +(up:1cm) .. (zh_l.north);
\draw[dashed, popmuted, thick] (fr_m.south) .. controls +(down:1cm) and +(up:1cm) .. (zh_m.north);
\draw[dashed, popmuted, thick] (fr_v.south) -- (zh_v.north);

% Labels
\node[above=0.2cm of fr_s, font=\bfseries\small, poporange] {Français (SVO + Compléments finaux)};
\node[below=0.2cm of zh_s, font=\bfseries\small, popblue] {Chinois (S -- T -- L -- M -- V)};
\end{tikzpicture}
\end{popfigure}

\courssub{Thèmes guidés : mise en application}

\begin{exemplebox}[Thème guidé 1]
\textbf{Phrase source :} « Je m'appelle Amina et je parle trois langues. »

\textbf{Analyse pas à pas :}
\begin{enumerate}
    \item \textbf{Segment 1 :} « Je m'appelle Amina » $\rightarrow$ Sujet \zh{我} (Wǒ) + Verbe \zh{叫} (jiào) + Objet \zh{阿明娜} (Āmíngnà). Structure SVO commune.
    \item \textbf{Segment 2 :} « et je parle trois langues » $\rightarrow$ Sujet \zh{我} (Wǒ) + Verbe modal \zh{会} (huì, \emph{savoir}) + Verbe \zh{说} (shuō, \emph{parler}) + Objet \zh{三种语言} (sān zhǒng yǔyán).
    \item \textbf{Point clé :} Le classificateur de \zh{语言} (langue) est \cle{种} (zhǒng), pas \zh{个}. \zh{会} exprime la compétence acquise.
\end{enumerate}

\textbf{Production finale :}
\zh{我叫阿明娜，我会说三种语言。}
\textit{Wǒ jiào Āmíngnà, wǒ huì shuō sān zhǒng yǔyán.}
\end{exemplebox}

\begin{exemplebox}[Thème guidé 2]
\textbf{Phrase source :} « De 2021 à 2024, j'ai étudié au lycée de Douala. »

\textbf{Analyse pas à pas :}
\begin{enumerate}
    \item \textbf{Repérage :} Sujet = « je » (\zh{我}). Verbe = « ai étudié » (\zh{学习} xuéxí, action achevée). Temps = « De 2021 à 2024 » (\zh{2021年到2024年}). Lieu = « au lycée de Douala » (\zh{在杜阿拉高中}).
    \item \textbf{Réordonnancement (Étape 3) :} Sujet \zh{我} + Temps \zh{2021年到2024年} + Lieu \zh{在杜阿拉高中} + Verbe \zh{学习}.
    \item \textbf{Aspect :} L'action est terminée dans le passé $\rightarrow$ ajouter \zh{了} (le) après le verbe pour marquer l'accompli.
\end{enumerate}

\textbf{Production finale :}
\zh{2021年到2024年，我在杜阿拉高中学习了。}
\textit{2021 nián dào 2024 nián, wǒ zài Dùyālā gāozhōng xuéxí le.}
\emph{Note :} La virgule chinoise sépare souvent l'adverbial de temps initial du reste de la phrase.
\end{exemplebox}

\begin{center}
\begin{tabularx}{\linewidth}{|X|X|X|X|}
\hline
\rowcolor{popblueL} \textbf{Français (Source)} & \textbf{Analyse (Temps/Lieu/Verbe/Aspect)} & \textbf{Ordre Chinois Cible} & \textbf{Chinois (Cible)} \\
\hline
Je m'appelle Amina & S: Je / V: m'appelle / O: Amina & S + V + O & \zh{我叫阿明娜} \\
\hline
Je parle trois langues & S: Je / V: sais parler / O: 3 langues & S + 会 + V + Num + 种 + N & \zh{我会说三种语言} \\
\hline
De 2021 à 2024, j'ai étudié à Douala & T: 2021-2024 / L: à Douala / V: étudié (+passé) & T + S + 在 + L + V + 了 & \zh{2021年到2024年，我在杜阿拉学习了} \\
\hline
J'ai travaillé dans une entreprise & T: (passé implicite) / L: dans une entreprise / V: travaillé & S + T + 在 + L + V + 了 & \zh{我在公司工作了} \\
\hline
Je sais utiliser l'ordinateur & S: Je / V: sais utiliser / O: ordinateur & S + 会 + V + O & \zh{我会用电脑} \\
\hline
J'ai fait un stage à Yaoundé & T: (passé) / L: à Yaoundé / V: stage & T + S + 在 + L + 实习 + 了 & \zh{2023年我在雅温得实习了} \\
\hline
\end{tabularx}
\end{center}

\courssub{Méthode de la version : du chinois vers le français}

\begin{methode}[Version : 5 étapes pour passer du chinois au français]
\textbf{Étape 1 : Segmenter la phrase chinoise.} Repérez les blocs : [Sujet] [Temps] [在 Lieu] [Manière] [Verbe 了/过/会] [Objet].\\
\textbf{Étape 2 : Identifier la rubrique du CV.} Est-ce \zh{教育背景} (Formation) ? \zh{工作经验} (Expérience) ? \zh{技能特长} (Compétences) ? Cela guide le vocabulaire français (ex: \zh{学习} $\rightarrow$ \emph{étudier / suivre des études} ; \zh{工作} $\rightarrow$ \emph{travailler / occuper un poste} ; \zh{实习} $\rightarrow$ \emph{faire un stage}).\\
\textbf{Étape 3 : Restituer l'ordre français (SVO).} Remettez le Temps et le Lieu \textbf{après} le verbe (ou en début de phrase avec préposition). Traduisez \zh{在} par \emph{à / au / en / dans}.\\
\textbf{Étape 4 : Gérer l'aspect et les classificateurs.} \zh{了} $\rightarrow$ passé composé / imparfait selon contexte. \zh{过} $\rightarrow$ « a déjà fait » (expérience de vie). \zh{会} $\rightarrow$ « sait / connaît / maîtrise » (jamais « will »). Supprimez les classificateurs chinois inutiles en français (pas de « trois \emph{sortes de} langues », mais « trois langues »).\\
\textbf{Étape 5 : Soigner le style professionnel.} Utilisez le vocabulaire du CV : \emph{maîtriser, assurer, coordonner, développer, gérer} au lieu de \emph{faire, savoir, mettre}.
\end{methode}

\begin{exemplebox}[Version guidée]
\textbf{Phrase source :} \zh{我的爱好是篮球和阅读。}
\textit{Wǒ de àihào shì lánqiú hé yuèdú.}

\textbf{Analyse :}
\begin{itemize}
    \item \zh{我的爱好} (Wǒ de àihào) : Sujet = « Mes loisirs / Mes centres d'intérêt » (\zh{的} marque l'appartenance).
    \item \zh{是} (shì) : Verbe « être » (identification).
    \item \zh{篮球和阅读} (lánqiú hé yuèdú) : Complément = « le basket-ball et la lecture » (\zh{和} = et).
\end{itemize}

\textbf{Traduction française naturelle :}
\emph{« Mes loisirs sont le basket-ball et la lecture. »} \\
\emph{Variante CV :} « Centres d'intérêt : basket-ball, lecture. »
\end{exemplebox}

\begin{exemplebox}[Version guidée : Phrase complexe]
\textbf{Phrase source :} \zh{2022年到2023年，我在杜阿拉的一家科技公司实习，学会了Python编程。}
\textit{2022 nián dào 2023 nián, wǒ zài Dùyālā de yī jiā kējì gōngsī shíxí, xuéhuì le Python biānchéng.}

\textbf{Analyse :}
\begin{itemize}
    \item \zh{2022年到2023年} : Adverbial de temps en tête.
    \item \zh{我在杜阿拉的一家科技公司实习} : Proposition 1. S=\zh{我}, L=\zh{在杜阿拉的一家科技公司}, V=\zh{实习} (pas de \zh{led} mais dates = passé).
    \item \zh{学会了Python编程} : Proposition 2 (conséquence/résultat). V=\zh{学会了} (verbe résultatif : apprendre + réussir), O=\zh{Python编程}.
\end{itemize}

\textbf{Traduction :}
\emph{« De 2022 à 2023, j'ai effectué un stage dans une entreprise technologique à Douala, où j'ai appris la programmation Python. »} \\
\emph{Variante CV :} « 2022-2023 : Stage en entreprise tech (Douala) — Acquisition de Python. »
\end{exemplebox}

\courssub{Pièges de traduction et points de grammaire critiques}

\begin{attention}[Piège n°1 : Le classificateur des langues et des catégories]
\emph{Erreur fréquente :} \zh{三个语言} (sān gè yǔyán) $\leftarrow$ Calque sur « trois langues » / « three languages ».
\emph{Correction :} \zh{三种语言} (sān \textbf{zhǒng} yǔyán).
\emph{Règle :} \cle{种} (zhǒng) est le classificateur pour les types, catégories, sortes (langues, fruits \zh{水果}, fleurs \zh{花}, styles \zh{风格}). \zh{个} (gè) est générique mais \textbf{incorrect} pour les langues.
\end{attention}

\begin{propriete}[Règle : \zh{会} (huì) = Compétence acquise, pas futur]
Dans un CV, \zh{会} ne traduit \textbf{jamais} le futur français (« je ferai », « je pourrai »).
\begin{itemize}
    \item \zh{我会说中文。} $\rightarrow$ « Je \textbf{sais parler} chinois / Je \textbf{maîtrise} le chinois. » (Compétence actuelle).
    \item \zh{我会用Excel。} $\rightarrow$ « Je \textbf{sais utiliser} Excel / Je \textbf{maîtrise} Excel. »
\end{itemize}
Pour exprimer le futur dans un autre contexte, on utilise \zh{将} (jiāng), \zh{会} (huì) + adverbe de temps futur (\zh{明年}...), ou \zh{打算} (dǎsuàn, « avoir l'intention de »). Mais dans un CV, on ne parle que de ce qu'on \emph{sait déjà faire}.
\end{propriete}

\begin{attention}[Piège n°2 : L'ordre des compléments (Temps / Lieu / Manière)]
\emph{Erreur fréquente :} \zh{我在杜阿拉2021年学习。} (Ordre français : Lieu avant Temps).
\emph{Correction :} \zh{我2021年在杜阿拉学习。} (Temps avant Lieu).
\emph{Astuce mnémotechnique :} « \textbf{Quand} ? $\rightarrow$ \textbf{Où} ? $\rightarrow$ \textbf{Comment ?} $\rightarrow$ \textbf{Qui fait quoi ?} » (Temps $\rightarrow$ Lieu $\rightarrow$ Manière $\rightarrow$ Verbe).
\end{attention}

\begin{attention}[Piège n°3 : L'article indéfini / défini et les métiers]
\emph{Erreur :} Traduire \zh{我是学生} par « Je suis \textbf{un} étudiant » $\rightarrow$ Correct (statut). Mais \zh{我是老师} $\rightarrow$ « Je suis professeur » (métier = pas d'article en français après \emph{être}).
\emph{Vigilance :} \zh{我有一个哥哥} $\rightarrow$ « J'ai \textbf{un} grand frère » (article nécessaire pour dénombrer). \zh{我哥哥是医生} $\rightarrow$ « Mon frère est médecin » (pas d'article pour la profession après \emph{être}).
\emph{Au CV :} \zh{求职意向：工程师} $\rightarrow$ « Poste visé : Ingénieur » (pas d'article).
\end{attention}

\begin{attention}[Piège n°4 : \zh{学会} vs \zh{学} + \zh{了}]
\zh{学会} (xuéhuì) = Verbe résultatif : « apprendre ET réussir à / maîtriser ». \emph{Focus sur le résultat.}
\zh{学了} (xué le) = « avoir étudié » (processus achevé, pas de garantie de maîtrise).
\textbf{CV :} Préférer \zh{学会了} / \zh{掌握了} (zhǎngwò le, maîtrisé) pour les compétences techniques.
\end{attention}

\begin{exemplebox}[Exemples traduits : Thème / Version croisés]
\begin{enumerate}
    \item \zh{我会用电脑。} \textit{Wǒ huì yòng diànnǎo.} $\rightarrow$ « Je sais utiliser l'ordinateur. / Je maîtrise l'outil informatique. »
    \item \zh{2020年到2023年，我在雅温得大学学习法律。} \textit{2020 nián dào 2023 nián, wǒ zài Yǎwēndé dàxué xuéxí fǎlǜ.} $\rightarrow$ « De 2020 à 2023, j'ai étudié le droit à l'université de Yaoundé. »
    \item \zh{我的技能包括英语和法语。} \textit{Wǒ de jìnéng bāokuò Yīngyǔ hé Fǎyǔ.} $\rightarrow$ « Mes compétences incluent l'anglais et le français. »
    \item \zh{我有两年工作经验。} \textit{Wǒ yǒu liǎng nián gōngzuò jīngyàn.} $\rightarrow$ « J'ai deux ans d'expérience professionnelle. » (\zh{有} = avoir, pas \zh{是}).
    \item \zh{我获得过优秀学生奖。} \textit{Wǒ huòdé guò yōuxiù xuéshēng jiǎng.} $\rightarrow$ « J'ai \textbf{déjà} obtenu le prix d'étudiant excellent. » (\zh{过} = expérience vécue au moins une fois).
    \item \zh{我能熟练操作办公软件。} \textit{Wǒ néng shúliàn cāozuò bàngōng ruǎnjiàn.} $\rightarrow$ « Je suis capable d'utiliser couramment les logiciels bureautiques. » (\zh{能} + \zh{熟练} = capacité avérée).
\end{enumerate}
\end{exemplebox}

\begin{popfigure}
\legende{Tableau récapitulatif : Les 5 étapes de la méthode Thème / Version}
\begin{tikzpicture}[
    node distance=0.3cm and 0.5cm,
    stepS/.style={rectangle, draw=popblue, thick, fill=popblueL, rounded corners, minimum height=1.1cm, text width=3.5cm, align=center, font=\small},
    stepv/.style={rectangle, draw=popgreen, thick, fill=popgreenL, rounded corners, minimum height=1.1cm, text width=3.5cm, align=center, font=\small},
    arrow/.style={-Stealth, thick, popdark}
]
% Theme column
\node[stepS, align=center] (t1){\textbf{THÈME (Fr $\rightarrow$ Zh)}\\\textbf{1. Identifier} : Sujet, Verbe, Compléments};
\node[stepS, below=of t1] (t2) {\textbf{2. Souligner} : Temps \& Lieu (souvent fin phrase Fr)};
\node[stepS, below=of t2] (t3) {\textbf{3. Réordonner} : S + \textcolor{poporange}{Temps} + \textcolor{popgreen}{Lieu} + \textcolor{poppurple}{Manière} + V (+ O)};
\node[stepS, below=of t3] (t4) {\textbf{4. Choisir} : Classif. (\zh{种}, \zh{家}, \zh{年}), Aspect (\zh{了}, \zh{过}), Modaux (\zh{会}, \zh{能})};
\node[stepS, below=of t4] (t5) {\textbf{5. Vérifier} : Tons, Ponctuation (，。), Caractères};

% Version column
\node[stepv, right=2cm of t1, align=center] (v1){\textbf{VERSION (Zh $\rightarrow$ Fr)}\\\textbf{1. Segmenter} : [S] [T] [在 L] [M] [V 了/过] [O]};
\node[stepv, below=of v1] (v2) {\textbf{2. Contextualiser} : Rubrique CV (Formation, Exp., Compétences)};
\node[stepv, below=of v2, align=center] (v3){\textbf{3. Restituer ordre Fr} : S + V + (Lieu) + (Temps) \\ \zh{在} $\rightarrow$ \emph{à/au/en/dans}};
\node[stepv, below=of v3] (v4) {\textbf{4. Adapter} : \zh{会}=sait/maîtrise ; \zh{过}=a déjà fait ; \zh{学会}=maîtriser ; Suppr. classif.};
\node[stepv, below=of v4] (v5) {\textbf{5. Style pro} : Vocab. action (gérer, développer), pas d'anglicismes};

% Arrows
\draw[arrow] (t1) -- (t2);
\draw[arrow] (t2) -- (t3);
\draw[arrow] (t3) -- (t4);
\draw[arrow] (t4) -- (t5);

\draw[arrow] (v1) -- (v2);
\draw[arrow] (v2) -- (v3);
\draw[arrow] (v3) -- (v4);
\draw[arrow] (v4) -- (v5);

% Horizontal links
\draw[dashed, popmuted] (t1.east) -- (v1.west);
\draw[dashed, popmuted] (t3.east) -- (v3.west) node[midway, above, font=\tiny, popmuted] {Inversion T/L/M};
\draw[dashed, popmuted] (t4.east) -- (v4.west) node[midway, above, font=\tiny, popmuted] {Choix lexicaux / Aspect};
\end{tikzpicture}
\end{popfigure}

\courssub{Exercice résolu : Thème et Version complets}

\begin{exoresolu}[Entraînement Bac : Paragraphe « Expérience \& Compétences »]
\textbf{Consigne :} Traduire le paragraphe suivant en chinois (Thème), puis retraduire la version chinoise proposée en français (Version).

\textbf{A. THÈME (Français $\rightarrow$ Chinois)}
« De 2022 à 2023, j'ai fait un stage dans une entreprise informatique à Douala. J'ai appris à programmer en Python. Je sais utiliser Excel et Word. Mes loisirs sont le football et la lecture. »

\textbf{B. VERSION (Chinois $\rightarrow$ Français)}
\zh{2022年到2023年，我在杜阿拉的一家科技公司实习。我学会了Python编程。我会用Excel和Word。我的爱好是足球和阅读。}

\tcblower

\textbf{Solution détaillée A (Thème)}

\begin{enumerate}
    \item \emph{« De 2022 à 2023, j'ai fait un stage dans une entreprise informatique à Douala. »}
    \begin{itemize}
        \item Temps : \zh{2022年到2023年}
        \item Lieu : \zh{在杜阿拉的一家科技公司} (dans une entreprise technologique de Douala). \zh{一家} (yī jiā) = classificateur pour entreprises/shops.
        \item Verbe : \zh{实习} (shíxí, faire un stage). Action achevée $\rightarrow$ \zh{了}.
        \item Structure : Temps, Sujet, Lieu, Verbe+了.
    \end{itemize}
    \zh{2022年到2023年，我在杜阿拉的一家科技公司实习了。}

    \item \emph{« J'ai appris à programmer en Python. »}
    \begin{itemize}
        \item « Appris à » = acquisition compétence $\rightarrow$ \zh{学会} (xuéhuì, verbe résultatif : apprendre + réussir à).
        \item Objet : \zh{Python编程} (Python biānchéng, programmation Python).
    \end{itemize}
    \zh{我学会了Python编程。}

    \item \emph{« Je sais utiliser Excel et Word. »}
    \begin{itemize}
        \item Compétence $\rightarrow$ \zh{会用} (huì yòng).
        \item Objets : \zh{Excel和Word}.
    \end{itemize}
    \zh{我会用Excel和Word。}

    \item \emph{« Mes loisirs sont le football et la lecture. »}
    \begin{itemize}
        \item Sujet : \zh{我的爱好}.
        \item Verbe : \zh{是}.
        \item Objets : \zh{足球和阅读}.
    \end{itemize}
    \zh{我的爱好是足球和阅读。}
\end{enumerate}

\textbf{Paragraphe chinois final :}
\zh{2022年到2023年，我在杜阿拉的一家科技公司实习了。我学会了Python编程。我会用Excel和Word。我的爱好是足球和阅读。}

\textbf{Solution détaillée B (Version)}

\begin{enumerate}
    \item \zh{2022年到2023年，我在杜阿拉的一家科技公司实习。}
    $\rightarrow$ « De 2022 à 2023, j'ai fait un stage dans une entreprise technologique à Douala. » (\zh{实习} sans \zh{了} ici mais contexte passé clair par les dates ; \zh{科技公司} = entreprise high-tech / informatique).
    \item \zh{我学会了Python编程。}
    $\rightarrow$ « J'ai appris la programmation Python / J'ai appris à programmer en Python. » (\zh{学会} = apprendre et maîtriser).
    \item \zh{我会用Excel和Word。}
    $\rightarrow$ « Je sais utiliser Excel et Word / Je maîtrise Excel et Word. » (\zh{会} = compétence).
    \item \zh{我的爱好是足球和阅读。}
    $\rightarrow$ « Mes loisirs sont le football et la lecture. » (CV : « Centres d'intérêt : football, lecture. »).
\end{enumerate}

\textbf{Note de correction :} Pour le thème phrase 1, l'ajout de \zh{了} après \zh{实习} est recommandé pour marquer l'accompli bien que les dates le suggèrent. En version, l'absence de \zh{了} ne change pas le passé car les dates ancrent l'action.
\end{exoresolu}

\courssub{Grille d'évaluation et production autonome}

\begin{propriete}[Grille d'évaluation du CV chinois (Bac LV2 — 20 points)]
\begin{center}
\begin{tabularx}{\linewidth}{|X|X|X|}
\hline
\rowcolor{popblueL} \textbf{Critère} & \textbf{Indicateurs de réussite} & \textbf{Points} \\
\hline
\textbf{Structure \& Rubriques} (4 pts) & 6 rubriques présentes, ordre logique, titres corrects (\zh{个人信息}...\zh{自我评价}) & 4 \\
\hline
\textbf{Vocabulaire spécialisé} (5 pts) & Mots justes : \zh{实习}, \zh{掌握}, \zh{熟练}, \zh{协助}, \zh{组织}, \zh{翻译}, HSK niveau, classificateurs (\zh{种}, \zh{家}, \zh{年}) & 5 \\
\hline
\textbf{Grammaire \& Ordre des mots} (5 pts) & S-T-L-M-V respecté, \zh{了}/\zh{过}/\zh{会} corrects, \zh{在} bien placé, connecteurs (\zh{首先}...) utilisés & 5 \\
\hline
\textbf{Caractères \& Écriture} (3 pts) & Caractères complexes (\zh{简历}, \zh{经验}, \zh{教育}, \zh{技能}, \zh{聘}) bien formés, traits dans l'ordre, pas de confusion proches & 3 \\
\hline
\textbf{Cohérence \& Style pro} (3 pts) & Ton formel, phrases complètes, auto-évaluation personnalisée, pas de fautes de soin & 3 \\
\hline
\textbf{TOTAL} & & \textbf{20} \\
\hline
\end{tabularx}
\end{center}
\end{propriete}

\begin{methode}[Production autonome : Rédiger son propre CV en chinois (Devoir maison / Bac blanc)]
\textbf{Étape 1 : Remplir la fiche d'identité.} Écrire en chinois : Nom chinois (transcription phonétique), Date naissance (AAAA年MM月), Ville (Yaoundé, Douala, Buea...), Téléphone, Email, Poste visé (\zh{实习生} / \zh{行政助理} / \zh{市场专员}...).
\textbf{Étape 2 : Construire \zh{教育背景}.} Lycée actuel + Lycée précédent. Spécialité (Série A/D/C). Mentionner HSK si passé. Notes moyennes si bonnes.
\textbf{Étape 3 : Détailler \zh{实习经历} / \zh{工作经验}.} Même si c'est un job d'été (baby-sitting, vente marché, aide bureau), formuler : \zh{时间} + \zh{地点} + \zh{职位} + \zh{职责} (verbes d'action : \zh{负责}, \zh{协助}, \zh{完成}, \zh{管理}). Chiffrer si possible (\zh{接待客户20人次}).
\textbf{Étape 4 : Lister \zh{技能特长}.} Langues (niveau CECRL/HSK), Informatique (Word, Excel, PPT, Réseaux sociaux, Montage vidéo), Permis, Autres (photo, chant, sport).
\textbf{Étape 5 : Rédiger \zh{爱好} et \zh{自我评价}.} 3 loisirs variés (sport, culture, social). Auto-évaluation : 4 phrases types (Personnalité + Atout majeur + Expérience clé + Motivation Chine/Cameroun).
\textbf{Étape 6 : Relecture croisée.} Échanger avec un camarade. Vérifier : Ordre T-L-V ? Classificateurs ? Tons ? Caractères complexes ? Ponctuation ?
\end{methode}

\begin{exercice}[Production écrite : Mon CV en chinois]
Rédigez votre CV complet en chinois simplifié sur une feuille A4 (ou dans le cahier) en suivant le modèle de la leçon.
\textbf{Contraintes :}
\begin{enumerate}
    \item Utiliser \textbf{au minimum 5} des caractères complexes étudiés : \zh{简历, 经验, 教育, 技能, 聘, 实习, 掌握, 熟练, 协助, 组织}.
    \item Intégrer \textbf{au moins 3 connecteurs} : \zh{首先}, \zh{其次}, \zh{此外}, \zh{最后}.
    \item Respecter l'ordre \textbf{S - Temps - Lieu - Verbe} dans toutes les phrases de la rubrique \zh{实习经历} et \zh{教育背景}.
    \item Écrire \zh{自我评价} en 4 phrases minimum.
\end{enumerate}
\textbf{Remise :} Version manuscrite (soin calligraphique) + Version tapée (Pinyin + Caractères) pour correction professeur.
\end{exercice}

\begin{corrige}[Exercice : Mon CV en chinois]
\textit{Correction individuelle. Le professeur vérifiera particulièrement :}
\begin{itemize}
    \item L'écriture correcte de \zh{简历} (radical \zh{竹}), \zh{验} (radical \zh{马}), \zh{育} (radical \zh{月}/\zh{肉}), \zh{技} (radical \zh{扌}), \zh{聘} (radical \zh{耳}).
    \item L'usage de \zh{种} pour les langues (\zh{三种语言}) et \zh{家} pour l'entreprise (\zh{一家公司}).
    \item La place de \zh{了} (action passée achevée) vs \zh{过} (expérience de vie) vs \zh{会} (compétence).
    \item L'ordre strict : \zh{2023年} (Temps) \zh{在杜阿拉} (Lieu) \zh{实习了} (Verbe).
    \item La cohérence du \zh{自我评价} : pas de futur (\zh{会} = sais), vocabulaire professionnel.
\end{itemize}
\textbf{Exemple de phrase type pour \zh{自我评价} :}
\zh{我性格严谨，责任心强。} (Personnalité) \\
\zh{我掌握中文HSK 3级和法语英语双语。} (Atout majeur) \\
\zh{我有在中资机构实习的经验。} (Expérience clé) \\
\zh{我希望能加入贵公司，为中喀合作贡献力量。} (Motivation)
\end{corrige}

\begin{savaistu}[Le Thème/Version au Baccalauréat : un entraînement CV]
Au Cameroun, l'épreuve de chinois LV2 au Baccalauréat (Terminale A) comporte systématiquement un exercice de \textbf{thème} (3 à 5 phrases) et un exercice de \textbf{version} (3 à 5 phrases). Les sujets portent quasi exclusivement sur :
\begin{itemize}
    \item La présentation personnelle (nom, âge, nationalité, famille).
    \item Le parcours scolaire (dates, établissements, spécialités) $\rightarrow$ Rubrique \zh{教育背景}.
    \item Les stages, jobs d'été, expériences bénévoles $\rightarrow$ Rubrique \zh{工作/实习经历}.
    \item Les langues parlées, logiciels maîtrisés, permis de conduire $\rightarrow$ Rubrique \zh{技能特长}.
    \item Les centres d'intérêt $\rightarrow$ Rubrique \zh{爱好}.
\end{itemize}
\textbf{Stratégie gagnante :} Traitez votre vrai CV (ou celui d'un camarade) comme un « texte ressource » tout au long de l'année. Traduisez chaque ligne en chinois, faites-le corriger, apprenez-le par cœur. Le jour de l'examen, vous n'aurez qu'à « piocher » dans vos phrases déjà validées. C'est la méthode la plus efficace pour garantir 18-20/20 à cette partie prévisible de l'épreuve.

\end{savaistu}

\begin{savaistu}[Le radical \zh{攵} (pū) : l'action d'enseigner]
Le radical \zh{攵} (forme compressée de \zh{攴} pū) vient d'une main tenant un bâton pour taper ou guider. Il figure dans de nombreux caractères liés à l'action, l'éducation, la gouvernance :
\begin{itemize}
    \item \zh{教} (jiāo) : enseigner (main guide l'enfant \zh{孝}).
    \item \zh{政} (zhèng) : politique, gouverner (main corrige \zh{正}).
    \item \zh{敲} (qiāo) : frapper, taper (main + haut \zh{高} phonétique).
    \item \zh{敬} (jìng) : respecter (main + \zh{苟} phonétique $\rightarrow$ attitude soignée).
    \item \zh{敦} (dūn) : sincère, épais (main + \zh{享} phonétique $\rightarrow$ action solide).
\end{itemize}
\emph{Astuce :} Quand vous voyez \zh{攵} à droite, pensez « action de la main / autorité / éducation ».
\end{savaistu}


\coursec{Grille d'évaluation et production autonome}

Après avoir travaillé la traduction thème/version dans la section précédente, nous passons maintenant à l'étape décisive : \textbf{l'évaluation de votre production écrite} et la \textbf{rédaction finale de votre propre mini-CV}. Cette section vous donne les outils pour juger la qualité d'un CV en chinois, corriger vos erreurs les plus fréquentes et produire un texte autonome, noté sur 20.

\courssub{Critères d'évaluation et grille sur 20}

L'évaluation d'un CV en chinois au niveau Terminale (HSK 3-4) repose sur quatre piliers. Chaque critère est noté sur un barème précis, pour un total de 20 points. Observez le tableau ci-dessous, qui sert de référence pour l'enseignant comme pour l'élève.

\begin{popfigure}
\begin{tikzpicture}[
    cell/.style={rectangle, draw=popline, minimum height=1.1cm, align=center, font=\small},
    header/.style={cell, fill=popblue, font=\bfseries, text=white},
    crit/.style={cell, fill=popblueL, font=\bfseries, text=popdark, minimum width=3.2cm},
    score/.style={cell, fill=popcream, minimum width=1.8cm},
    desc/.style={cell, fill=white, minimum width=6.5cm},
]
% En-têtes
\node[header, minimum width=3.2cm] (h1) at (0,0) {Critère};
\node[header, minimum width=1.8cm] (h2) at (3.2,0) {Points};
\node[header, minimum width=6.5cm] (h3) at (5.0,0) {Descripteurs (niveau « Maîtrise satisfaisante »)};
% Ligne 1 : Caractères
\node[crit] (c1) at (0,-1.1) {Caractères};
\node[score] (s1) at (3.2,-1.1) {/6};
\node[desc] (d1) at (5.0,-1.1) {Écriture exacte des 30-40 caractères clés (traits, ordre, radicaux). Pas de confusion 简/艰, 历/厉, 技/枝. Pinyin correct si fourni.};
% Ligne 2 : Grammaire
\node[crit] (c2) at (0,-2.2) {Grammaire \& ordre};
\node[score] (s2) at (3.2,-2.2) {/6};
\node[desc] (d2) at (5.0,-2.2) {Structure Sujet + Temps + Lieu + Verbe respectée. Bons classificateurs (个, 年, 门, 种). 了/过/着 bien employés. Négation 不/没 juste.};
% Ligne 3 : Contenu
\node[crit] (c3) at (0,-3.3) {Contenu \& rubriques};
\node[score] (s3) at (3.2,-3.3) {/4};
\node[desc] (d3) at (5.0,-3.3) {6 rubriques présentes : 个人信息, 求职意向, 教育背景, 工作经验, 技能特长, 自我评价. Infos vraies, cohérentes, 6-8 phrases.};
% Ligne 4 : Connecteurs
\node[crit] (c4) at (0,-4.4) {Connecteurs \& présentation};
\node[score] (s4) at (3.2,-4.4) {/4};
\node[desc] (d4) at (5.0,-4.4) {Utilisation de 3+ connecteurs (首先, 其次, 此外, 最后, 因为…所以…). Mise en page aérée, puces ou retraits, police lisible. Pas de virgules chinoises en fin de ligne.};
% Total
\node[header, minimum width=3.2cm] (ht1) at (0,-5.5) {TOTAL};
\node[header, minimum width=1.8cm] (ht2) at (3.2,-5.5) {/20};
\node[header, minimum width=6.5cm] (ht3) at (5.0,-5.5) {≥ 16 = Excellent | 12-15 = Bon | 8-11 = Passable | < 8 = À retravailler};
% Lignes de séparation
\draw[popline] (0,0) -- (11.5,0);
\draw[popline] (0,-1.1) -- (11.5,-1.1);
\draw[popline] (0,-2.2) -- (11.5,-2.2);
\draw[popline] (0,-3.3) -- (11.5,-3.3);
\draw[popline] (0,-4.4) -- (11.5,-4.4);
\draw[popline] (0,-5.5) -- (11.5,-5.5);
\draw[popline] (0,-6.6) -- (11.5,-6.6);
\draw[popline] (0,0) -- (0,-6.6);
\draw[popline] (3.2,0) -- (3.2,-6.6);
\draw[popline] (5.0,0) -- (5.0,-6.6);
\draw[popline] (11.5,0) -- (11.5,-6.6);
\end{tikzpicture}
\legende{Grille d'évaluation du mini-CV chinois (Terminale A) — 4 critères, 20 points}
\end{popfigure}

\begin{propriete}[Règle — Un bon CV chinois est bref]
Au niveau Terminale (HSK 3-4), \textbf{un CV efficace tient en 6 à 8 phrases complètes}. On ne rédige pas de paragraphes longs : chaque rubrique = une phrase (ou deux au maximum pour \zh{工作经验}). La concision prouve la maîtrise de la synthèse, qualité très prisée dans le monde professionnel chinois. \textbf{Exemple de structure visée :}
\begin{itemize}
    \item Phrase 1 : \zh{我叫李明，今年十八岁，来自喀麦隆。} (Infos persos)
    \item Phrase 2 : \zh{我希望申请贵公司的市场助理岗位。} (Objectif)
    \item Phrase 3 : \zh{二〇二一年至二〇二四年，我在雅温得第二中学就读高中。} (Éducation)
    \item Phrase 4 : \zh{暑假我在杜阿拉的一家中餐馆做过服务员。} (Expérience)
    \item Phrase 5 : \zh{我会说中文、法文和英语，会用办公软件。} (Compétences)
    \item Phrase 6 : \zh{我性格开朗，吃苦耐劳，有团队合作精神。} (Auto-éval.)
\end{itemize}
\end{propriete}

\courssub{Auto-évaluation et évaluation croisée entre pairs}

Avant de remettre votre copie, vous devez la relire \textbf{trois fois}, chacune avec un focus différent. C'est la méthode des « trois lectures » utilisée par les correcteurs du baccalauréat.

\begin{methode}[S'auto-corriger : la méthode des trois lectures]
\textbf{Lecture 1 — Caractères (yeux sur le trait)} \\
Parcourez chaque caractère lentement. Vérifiez : nombre de traits, ordre des traits, radicaux. Comparez avec votre fiche de révision des 40 caractères du module. Entourez au crayon tout caractère douteux. \\
\textbf{Lecture 2 — Grammaire et ordre des mots (yeux sur la structure)} \\
Isolez chaque phrase. Identifiez : Sujet — Temps — Lieu — Verbe — Complément. Vérifiez la place des compléments (le lieu \textbf{précède} le verbe en chinois !), les classificateurs, les particules aspectuelles (\zh{了, 过, 着}). \\
\textbf{Lecture 3 — Contenu, connecteurs, présentation (yeux sur le sens)} \\
Le CV a-t-il les 6 rubriques ? Les infos sont-elles logiques ? Y a-t-il au moins 3 connecteurs (\zh{首先, 其次, 此外, 最后, 因为…所以…}) ? La mise en page est-elle propre (retraits, puces, pas de ratures) ?
\end{methode}

\begin{popfigure}
\begin{tikzpicture}[
    check/.style={rectangle, draw=popgreen, minimum size=0.7cm, thick},
    label/.style={font=\small, text=popdark, anchor=west, align=left, text width=10cm},
    title/.style={font=\bfseries\large, text=popblue, anchor=west}
]
\node[title] at (0,4.5) {✅  Check-list d'auto-évaluation avant remise};
% Ligne 1
\node[check] (c1) at (0,3.5) {};
\node[label] at (1,3.5) {Traits corrects : ordre, nombre, radicaux (简历, 经验, 教育, 技能, 聘…)};
% Ligne 2
\node[check] (c2) at (0,2.5) {};
\node[label] at (1,2.5) {Ordre des mots : Sujet + Temps + Lieu + Verbe + Complément};
% Ligne 3
\node[check] (c3) at (0,1.5) {};
\node[label] at (1,1.5) {6 rubriques présentes : 个人信息, 求职意向, 教育背景, 工作经验, 技能特长, 自我评价};
% Ligne 4
\node[check] (c4) at (0,0.5) {};
\node[label] at (1,0.5) {Connecteurs utilisés ( ≥ 3 ) : 首先, 其次, 此外, 最后, 因为…所以… };
% Ligne 5
\node[check] (c5) at (0,-0.5) {};
\node[label] at (1,-0.5) {Ponctuation chinoise ： ， 。 ； pas de virgule française, pas de point final anglais .};
% Ligne 6
\node[check] (c6) at (0,-1.5) {};
\node[label] at (1,-1.5) {Pinyin exact si demandé (tons sur voyelles : ā á ǎ à, ē é ě è…)};
% Décoration : flèches
\draw[popgreen, thick, ->] (0.35,3.5) -- (0.35,3.2);
\draw[popgreen, thick, ->] (0.35,2.5) -- (0.35,2.2);
\draw[popgreen, thick, ->] (0.35,1.5) -- (0.35,1.2);
\draw[popgreen, thick, ->] (0.35,0.5) -- (0.35,0.2);
\draw[popgreen, thick, ->] (0.35,-0.5) -- (0.35,-0.8);
\draw[popgreen, thick, ->] (0.35,-1.5) -- (0.35,-1.8);
\end{tikzpicture}
\legende{Check-list visuelle à cocher avant de rendre son CV}
\end{popfigure}

Une fois votre auto-évaluation faite, échangez votre brouillon avec un camarade. L'évaluateur pair utilise la \textbf{même grille sur 20} et note en marge : \zh{✓} (bon), \zh{✗} (erreur), \zh{?} (incompréhensible). Il écrit un commentaire constructif en bas : \zh{字写得很工整，但“在中学学习”语序不对。} (« Caractères soignés, mais ordre des mots incorrect pour "étudier au lycée" »).

\courssub{Correction collective : 5 erreurs types sur un CV élève}

Voici un CV rédigé par un élève de Terminale. Il contient \textbf{5 erreurs volontaires} représentatives des pièges les plus fréquents. Nous allons l'analyser ensemble.

\begin{attention}[Une erreur de trait peut changer le sens — la précision graphique est notée]
\begin{itemize}
    \item \zh{历} (lì, histoire, expérience) vs \zh{厉} (lì, sévère, féroce) : \zh{工作经历} (expérience pro) ≠ \zh{工作经厉} (sans sens).
    \item \zh{技} (jì, compétence, technique) vs \zh{枝} (zhī, branche, classif. pour crayons) : \zh{技能} (compétence) ≠ \zh{枝能}.
    \item \zh{简} (jiǎn, simple) vs \zh{艰} (jiān, difficile) : \zh{简历} (CV) ≠ \zh{艰历}.
    \item \zh{聘} (pìn, embaucher) : radical \zh{讠} (parole, à gauche) + phonétique \zh{甫} (fǔ). Ne pas confondre avec \zh{职} (zhí, poste, radical \zh{耳}).
    \item \zh{职} (zhí, poste) vs \zh{植} (zhí, planter) : \zh{职位} (poste) ≠ \zh{植位}.
\end{itemize}
\textbf{Conseil :} Pour chaque caractère complexe, identifiez son \textbf{radical-clé} (clé de dictionnaire) et son \textbf{phonétique}. Ex. \zh{技} : radical \zh{扌} (main) + phonétique \zh{枝} (zhī) → action manuelle habile.
\end{attention}

\begin{textebox}[CV élève à corriger — 5 erreurs à repérer]
个人艰历 \\
姓名：王伟 \quad 性别：男 \quad 年龄：18 岁 \\
电话：6 90 12 34 56 \quad 邮箱：wangwei@email.cn \\
求职意向：市场助理 \\
教育背景：二〇二一年九月 到 二〇二四年六月，我 学习 在 雅温得第一中学。 \\
工作经厉：二〇二三年暑假，我 在 一家 中国 餐馆 做 服务员，学 到 很多。 \\
技能特长：我 会 说 中文、法文、英文，会 用 办公 软件，还 会 做 中国 菜。 \\
自我评价：我 性格 开朗，不 怕 吃苦，因为 我 有 团队 精神，所以 适合 这 个 工作。
\end{textebox}

\begin{exoresolu}[Correction collective guidée]
\textbf{Énoncé :} Repérez et corrigez les 5 erreurs dans le CV ci-dessus. \\
\textbf{Solution détaillée :} \\
\noindent\textbf{Erreur 1 — Caractère faux (radical) dans le titre :} \zh{个人艰历}. \zh{简} (simple, radical \zh{竹}) est remplacé par \zh{艰} (difficile, radical \zh{艮}). \textbf{Correction :} \zh{个人简历} (gèrén jiǎnlì). \\
\noindent\textbf{Erreur 2 — Caractère faux (phonétique) dans la rubrique Expérience :} \zh{工作经厉}. \zh{历} (lì, expérience, radical \zh{止}) est remplacé par \zh{厉} (lì, sévère, radical \zh{厂}). \textbf{Correction :} \zh{工作经验} (gōngzuò jīngyàn). \\
\noindent\textbf{Erreur 3 — Ordre des mots (Lieu / Verbe) :} \zh{我 学习 在 雅温得第一中学} (rubrique Éducation). En chinois, le complément de lieu introduit par \zh{在} se place \textbf{avant} le verbe. \textbf{Correction :} \zh{我 在 雅温得第一中学 就读} (jiùdú, « étudier dans ») ou \zh{我 在 雅温得第一中学 读书}. Structure : \textbf{Sujet + 在 + Lieu + Verbe}. \\
\noindent\textbf{Erreur 4 — Complément de résultat incomplet / Particule aspectuelle manquante :} \zh{学 到 很多} (rubrique Expérience). \zh{学到} (verbe + complément de résultat) appelle un objet. De plus, action passée → \zh{了}. \textbf{Correction :} \zh{学到了很多经验} (xué dào le hěn duō jīngyàn) ou \zh{学到了很多东西}. \\
\noindent\textbf{Erreur 5 — Sujet manquant après 所以 :} \zh{因为 我 有 团队 精神，所以 适合 这 个 工作} (rubrique Auto-éval). \zh{所以} introduit une conséquence ; le sujet \zh{我} doit être répété (ou pronominalisé). \textbf{Correction :} \zh{因为 我 有 团队 精神，所以 \underline{我} 适合 这 个 工作}. Plus fluide : \zh{我 有 团队 精神，\underline{因此} 适合 这 个 工作}. \\
\noindent\textbf{Bilan :} 2 erreurs de caractères (简/艰, 历/厉), 1 erreur d'ordre des mots (Lieu/Verbe), 1 erreur de structure verbale (学到/了), 1 erreur de syntaxe (sujet manquant après 所以).
\end{exoresolu}

\begin{exemplebox}[Phrase fautive → correcte : l'ordre Lieu + Verbe]
\begin{center}
\begin{tabularx}{\linewidth}{|X|X|X|}\hline
\rowcolor{popblueL} \textbf{Fautive (calque français)} & \textbf{Correcte (structure chinoise)} & \textbf{Traduction} \\ \hline
\zh{我 学习 在 中学} & \zh{我 在 中学 学习} / \zh{我 在 中学 读书} & J'étudie au lycée \\ \hline
\zh{我 工作 在 杜阿拉} & \zh{我 在 杜阿拉 工作} / \zh{我 在 杜阿拉 上班} & Je travaille à Douala \\ \hline
\zh{他 住 在 喀麦隆} & \zh{他 住 在 喀麦隆} (correct : verbe d'état) & Il habite au Cameroun \\ \hline
\end{tabularx}
\end{center}
\emph{Règle d'or :} Avec les verbes d'action (travailler, étudier, manger…), le complément de lieu (\zh{在} + lieu) se place \textbf{avant} le verbe. C'est l'erreur n°1 des francophones. (Exception : verbes de posture/existence \zh{住, 放, 站, 坐, 躺} → \zh{V + 在 + Lieu} possible).
\end{exemplebox}

\courssub{Rédaction finale : votre mini-CV individuel (6-8 phrases)}

C'est à vous maintenant. Vous allez rédiger \textbf{votre propre mini-CV} en chinois simplifié, sur une feuille A4, en respectant la grille d'évaluation. Vous avez 30 minutes.

\begin{savaistu}[Les entreprises chinoises au Cameroun apprécient le CV bilingue]
Depuis une vingtaine d'années, de grandes entreprises chinoises sont implantées au Cameroun : \textbf{China Harbour Engineering Company (CHEC)} au port en eau profonde de \textbf{Kribi}, \textbf{Sinohydro} pour les barrages (Lom Pangar, Nachtigal), \textbf{Huawei} et \textbf{ZTE} pour les télécoms à \textbf{Yaoundé} et \textbf{Douala}, \textbf{China Road and Bridge Corporation} pour les routes. Ces entreprises recrutent localement et \textbf{valorisent les candidats capables de présenter un CV bilingue français-chinois \textbf{même sommaire}}. Cela prouve : (1) une ouverture culturelle, (2) une rigueur graphique (caractères), (3) une capacité de synthèse. Au lycée, ce mini-CV est votre premier « passeport » pour un stage ou un job d'été dans ces structures.
\end{savaistu}

\begin{exercice}[Rédaction finale — Mon mini-CV en chinois (30 min)]
\textbf{Consignes :}
\begin{enumerate}
    \item Rédigez votre CV sur une feuille A4, à l'encre bleue ou noire. Écrivez \textbf{uniquement en caractères chinois simplifiés} (pinyin optionnel en petit sous chaque ligne).
    \item Respectez les \textbf{6 rubriques} dans l'ordre : \zh{个人信息} — \zh{求职意向} — \zh{教育背景} — \zh{工作经验} — \zh{技能特长} — \zh{自我评价}.
    \item Écrivez \textbf{6 à 8 phrases complètes} (une par rubrique, deux max pour Expérience).
    \item Utilisez \textbf{au moins 3 connecteurs} parmi : \zh{首先, 其次, 此外, 最后, 因为…所以…, 不但…而且…}.
    \item Soignez l'écriture : traits dans l'ordre, radicaux visibles, pas de « pattes de mouche ».
    \item Relisez avec la \textbf{check-list} (3 lectures) avant de rendre.
\end{enumerate}
\textbf{Vocabulaire d'appui (révisé du module) :}
\begin{center}
\begin{tabularx}{\linewidth}{|X|X|X|X|}\hline
\rowcolor{popblueL} Caractères & Pinyin & Nature & Traduction \\ \hline
姓名 & xìngmíng & n & nom complet \\ \hline
求职意向 & qiúzhí yìxiàng & n & poste visé \\ \hline
市场助理 & shìchǎng zhùlǐ & n & assistant marketing \\ \hline
教育背景 & jiàoyù bèijǐng & n & formation \\ \hline
就读 & jiùdú & v & étudier (dans une école) \\ \hline
工作经验 & gōngzuò jīngyàn & n & expérience pro \\ \hline
实习 & shíxí & v/n & stage / faire un stage \\ \hline
技能特长 & jìnéng tècháng & n & compétences / atouts \\ \hline
办公软件 & bàngōng ruǎnjiàn & n & logiciels bureautiques \\ \hline
性格开朗 & xìnggé kāilǎng & adj & extraverti, gai \\ \hline
吃苦耐劳 & chīkǔ nàiláo & idiome & courageux, endurant \\ \hline
团队合作精神 & tuánduì hézuò jīngshén & n & esprit d'équipe \\ \hline
\end{tabularx}
\end{center}
\end{exercice}

\begin{corrige}[Exercice — Rédaction finale : modèle de mini-CV élève (exemple)]
\textbf{Voici un modèle de production attendue (niveau « Bon », ~14/20). L'élève adapte les infos vraies.}

\begin{textebox}[Mini-CV modèle — Wang Wei (élève Terminale A)]
\textbf{个人简历} \\
\textbf{个人信息}：我叫王伟，男，二〇〇六年生，来自喀麦隆雅温得。\\
\textbf{求职意向}：我希望申请贵公司的市场助理岗位。\\
\textbf{教育背景}：首先，二〇二一年至二〇二四年，我在雅温得第一中学就读高中。\\
\textbf{工作经验}：其次，二〇二三年暑假，我在杜阿拉一家中餐馆实习过，当过服务员，学到了服务经验。\\
\textbf{技能特长}：此外，我会说中文、法文和英文，熟练使用办公软件，还会做简单的中国菜。\\
\textbf{自我评价}：最后，因为我性格开朗，不但吃苦耐劳，而且有团队合作精神，所以我相信我适合这个工作。
\end{textebox}

\noindent\textbf{Analyse selon la grille :}
\begin{itemize}
    \item \textbf{Caractères /6} : 5,5 — écriture soignée, \zh{简历/经验/技能/聘} corrects, un trait hésitant sur \zh{耐}.
    \item \textbf{Grammaire /6} : 5 — ordre S+T+L+V respecté, \zh{在…实习过/当过} bien construits, \zh{不但…而且…} correct.
    \item \textbf{Contenu /4} : 4 — 6 rubriques, infos cohérentes, 7 phrases.
    \item \textbf{Connecteurs /4} : 4 — \zh{首先, 其次, 此外, 最后, 因为…所以…, 不但…而且…} (6 connecteurs !).
    \item \textbf{Total : 18,5/20} → Excellent. Présentation aérée, puces claires.
\end{itemize}

\end{corrige}

\begin{aretenir}[Bilan de la leçon 40 — Ce qu'il faut retenir pour le Bac]
\begin{itemize}
    \item Un CV chinois = \textbf{6 rubriques standardisées} + \textbf{6-8 phrases} + \textbf{connecteurs logiques}.
    \item \textbf{Ordre des mots immuable} : Sujet + Temps + \zh{在} Lieu + Verbe (+ \zh{了/过}).
    \item \textbf{Caractères complexes du module} : \zh{简历, 经验, 教育, 技能, 聘, 职, 耐, 劳, 实, 习} — connaître radical, ordre traits, phonétique.
    \item \textbf{Classificateurs} : \zh{一家} (entreprise/resto), \zh{一门} (langue/matière), \zh{一种} (compétence), \zh{个} (période générique).
    \item \textbf{Auto-correction} : 3 lectures (caractères → grammaire → sens) + check-list visuelle.
    \item \textbf{Contexte réel} : CV bilingue = atout majeur pour stages dans entreprises chinoises au Cameroun (Kribi, Douala, Yaoundé, Buea).
\end{itemize}
\end{aretenir}

\vspace{1em}
\noindent\textbf{Félicitations !} Vous avez terminé le Module 5 « Mass médias et télécommunication ». Vous savez maintenant rédiger, évaluer et corriger un CV en chinois, compétence concrète pour votre avenir professionnel. \zh{加油 !}

\newpage

\coursec{Activité d'intégration}

\coursec{Leçon 40 — Expression écrite : caractères complexes liés à la rédaction d'un CV}

\courssub{Découverte : Le CV chinois et ses rubriques}
\begin{textebox}{Appel à candidatures — Bourse d'excellence Gouvernement Chinois (CSC) / Ambassade de Chine au Cameroun}
\zh{中国驻喀麦隆大使馆 教育处} \\
\zh{关于 2025 年 中国政府奖学金 选拔 通知} \\
\zh{各位 同学：} \\
\zh{为 促进 中喀 两国 教育 交流，中国政府 奖学金 现 面向 喀麦隆 优秀 高中 毕业生 开放 申请。} \\
\zh{申请 材料 包括：} \\
\zh{1. 个人 简历 (中文版，手写 或 打印 均可)；} \\
\zh{2. 最高 学历 证明 及 成绩单；} \\
\zh{3. 汉语 水平 证明 (HSK 3 级 以上 优先) 或 法语/英语 成绩；} \\
\zh{4. 推荐信 两 封。} \\
\zh{请 于 3 月 31 日 前 将 材料 发送 至 邮箱：education@chinaembassy.cm。} \\
\zh{初选 通过 者 将 受邀 参加 面试 (中文 简短 自我介绍 及 答辩)。} \\
\zh{祝 同学们 取得 优异 成绩！} \\
\zh{中国驻喀麦隆大使馆 教育处} \\
\zh{2025 年 1 月 15 日}
\tcblower
\textbf{Traduction :} Ambassade de Chine au Cameroun — Section Éducation. Avis de sélection pour la Bourse du Gouvernement Chinois 2025. Chers élèves, pour promouvoir les échanges éducatifs sino-camerounais, la bourse du gouvernement chinois est ouverte aux excellents bacheliers camerounais. Dossier : 1. CV personnel (version chinoise, manuscrit ou tapé) ; 2. Diplôme le plus élevé et relevés de notes ; 3. Preuve de niveau chinois (HSK 3+ prioritaire) ou résultats français/anglais ; 4. Deux lettres de recommandation. Envoi avant le 31 mars à education@chinaembassy.cm. Les présélectionnés seront convoqués pour un entretien (brève présentation orale en chinois et questions). Bonne chance à tous ! Section Éducation, Ambassade de Chine au Cameroun, 15 janvier 2025.
\end{textebox}

\begin{definition}[Vocabulaire clé : Les rubriques standard du CV chinois (\zh{简历标准栏目})]
\begin{center}
\begin{tabularx}{\linewidth}{|X|X|X|X|}
\hline
\rowcolor{popblueL} \textbf{Caractères} & \textbf{Pinyin} & \textbf{Nature} & \textbf{Traduction} \\
\hline
\zh{个人信息} & gèrén xìnxī & Nom & Informations personnelles \\
\hline
\zh{教育背景} & jiàoyù bèijǐng & Nom & Formation / Parcours scolaire \\
\hline
\zh{工作经验} / \zh{实习经验} & gōngzuò jīngyàn / shíxí jīngyàn & Nom & Expérience professionnelle / Stage \\
\hline
\zh{技能特长} & jìnéng tècháng & Nom & Compétences et talents \\
\hline
\zh{语言能力} & yǔyán nénglì & Nom & Compétences linguistiques \\
\hline
\zh{社会实践} & shèhuì shíjiàn & Nom & Pratique sociale / Bénévolat \\
\hline
\zh{爱好} & àihào & Nom & Centres d'intérêt / Loisirs \\
\hline
\zh{自我评价} & zìwǒ píngjià & Nom & Auto-évaluation \\
\hline
\end{tabularx}
\end{center}
\end{definition}

\courssub{Écriture des caractères complexes (1) : \zh{简历}、\zh{经验}}
\begin{definition}[Analyse graphique : \zh{简} (jiǎn), \zh{历} (lì), \zh{验} (yàn)]
\begin{center}
\begin{tabularx}{\linewidth}{|X|X|X|X|}
\hline
\rowcolor{popblueL} \textbf{Caractère} & \textbf{Pinyin} & \textbf{Radical (Clé) \& Nb traits} & \textbf{Aide-mémoire / Piège \& Ordre des traits} \\
\hline
\zh{简} & jiǎn & \zh{竹} (zhú, bambou, 6 tr.) + 12 = 18 tr. & \textbf{Structure :} Haut \zh{竹} (simplifié de \zh{⺮}), bas \zh{间} (jiān). \textbf{Ordre :} \zh{竹} (haut→bas, gauche→droite), puis \zh{间} (porte \zh{门} + soleil \zh{日}). \textbf{Piège :} Ne pas confondre avec \zh{检} (jiǎn, inspecter, radical \zh{木}). \\
\hline
\zh{历} & lì & \zh{止} (zhǐ, arrêter, 4 tr.) + 5 = 9 tr. & \textbf{Structure :} Simplifié de \zh{曆} (soleil \zh{日} + \zh{历}). \textbf{Ordre :} \zh{止} (horizontale, verticale, horizontale, crochet), puis \zh{厂} (haut→bas), puis \zh{口}. Sens : passer par, calendrier, expérience. \\
\hline
\zh{验} & yàn & \zh{马} (mǎ, cheval, 10 tr.) + 8 = 18 tr. & \textbf{Structure :} Gauche \zh{马} (sens/phonétique : monter à cheval pour tester), droite \zh{乛} + \zh{工} + \zh{十}. \textbf{Ordre :} \zh{马} d'abord, puis partie droite (haut→bas). \textbf{Piège :} Ne pas écrire le traditionnel \zh{騐} ni confondre avec \zh{骗} (piàn, tromper, radical \zh{马} mais structure diff.). \\
\hline
\end{tabularx}
\end{center}
\end{definition}

\begin{exemplebox}[Mots composés et phrases modèles]
\begin{itemize}
    \item \zh{简历} (jiǎnlì) — CV, résumé. \emph{Ex:} \zh{请 投递 你 的 简历。} (Qǐng tóudì nǐ de jiǎnlì.) — Veuillez envoyer votre CV.
    \item \zh{经验} (jīngyàn) — Expérience. \emph{Ex:} \zh{我 有 三 年 教学 经验。} (Wǒ yǒu sān nián jiàoxué jīngyàn.) — J'ai trois ans d'expérience dans l'enseignement.
    \item \zh{经历} (jīnglì) — Parcours, expérience vécue. \emph{Ex:} \zh{这 是 一段 宝贵 的 经历。} (Zhè shì yī duàn bǎoguì de jīnglì.) — C'est une expérience précieuse.
    \item \zh{检验} (jiǎnyàn) — Inspecter, tester. \emph{Ex:} \zh{质量 检验 合格。} (Zhìliàng jiǎnyàn hégé.) — Le contrôle qualité est conforme.
\end{itemize}
\end{exemplebox}

\courssub{Écriture des caractères complexes (2) : \zh{教育}、\zh{技能}、\zh{聘}}
\begin{definition}[Analyse graphique : \zh{教} (jiào), \zh{育} (yù), \zh{技} (jì), \zh{能} (néng), \zh{聘} (pìn)]
\begin{center}
\begin{tabularx}{\linewidth}{|X|X|X|X|}
\hline
\rowcolor{popblueL} \textbf{Caractère} & \textbf{Pinyin} & \textbf{Radical (Clé) \& Nb traits} & \textbf{Aide-mémoire / Piège \& Ordre des traits} \\
\hline
\zh{教} & jiào & \zh{攵} (pū, action main/bâton, 4 tr.) + 7 = 11 tr. & \textbf{Structure :} Haut \zh{孝} (xiào, piété filiale) simplifié, bas \zh{攵} (frapper/enseigner). \textbf{Ordre :} Haut→bas ; \zh{攵} = point, horizontale, crochet vertical, point. \\
\hline
\zh{育} & yù & \zh{月} (yuè, chair/lune, 4 tr.) + 9 = 13 tr. & \textbf{Structure :} Haut \zh{立} (lì, debout) + \zh{月} (corps). Élever un enfant. \textbf{Ordre :} \zh{立} puis \zh{月}. \textbf{Piège :} Ne pas ajouter un trait au \zh{月} (≠ \zh{毓} yù traditionnel). \\
\hline
\zh{技} & jì & \zh{扌} (shǒu, main, 3 tr.) + 4 = 7 tr. & \textbf{Structure :} Phonétique \zh{枝} (zhī) sans bois \zh{木}. Main qui travaille l'artisanat. \textbf{Ordre :} \zh{扌} (horizontale, crochet vertical, relevé), puis \zh{枝} simplifié (horizontal, vertical, horizontal). \\
\hline
\zh{能} & néng & \zh{月} (yuè, chair/lune, 4 tr.) + 6 = 10 tr. & \textbf{Structure :} Anciennement \zh{能} (ours). Haut \zh{匕} + \zh{厶} + \zh{月}. Capacité, talent. \textbf{Ordre :} Haut→bas, gauche→droite. \\
\hline
\zh{聘} & pìn & \zh{耳} (ěr, oreille, 6 tr.) + 7 = 13 tr. & \textbf{Structure :} Gauche \zh{耳} (écouter), droite \zh{并} (bìng, ensemble). Inviter officiellement (écouter la proposition). \textbf{Ordre :} \zh{耳} puis \zh{并}. \textbf{Piège :} \zh{并} (6 tr., deux \zh{立} côte à côte) ≠ \zh{关} (guān, 6 tr., \zh{门} + \zh{关}). \\
\hline
\end{tabularx}
\end{center}
\end{definition}

\begin{exemplebox}[Mots composés et phrases modèles]
\begin{itemize}
    \item \zh{教育} (jiàoyù) — Éducation. \emph{Ex:} \zh{他 从事 教育 工作 十 年 了。} (Tā cóngshì jiàoyù gōngzuò shí nián le.) — Il travaille dans l'éducation depuis 10 ans.
    \item \zh{技能} (jìnéng) — Compétence, habileté. \emph{Ex:} \zh{熟练 掌握 计算机 技能。} (Shúliàn zhǎngwò jìsuànjī jìnéng.) — Maîtriser les compétences informatiques.
    \item \zh{能力} (nénglì) — Capacité, aptitude. \emph{Ex:} \zh{提高 沟通 能力。} (Tígāo gōutōng nénglì.) — Améliorer la capacité de communication.
    \item \zh{聘用} (pìnyòng) / \zh{招聘} (zhāopìn) — Embaucher / Recruter. \emph{Ex:} \zh{公司 正在 招聘 实习生。} (Gōngsī zhèngzài zhāopìn shíxīshēng.) — L'entreprise recrute des stagiaires.
\end{itemize}
\end{exemplebox}

\courssub{Modèle de CV commenté, structure et connecteurs}
\begin{propriete}[Structures grammaticales et connecteurs utiles pour le CV]
\begin{center}
\begin{tabularx}{\linewidth}{|X|X|X|}
\hline
\rowcolor{popblueL} \textbf{Fonction} & \textbf{Structure / Modèle} & \textbf{Exemple CV} \\
\hline
Organiser le discours (Connecteurs) & \zh{首先 / 其次 / 此外 / 最后，} + Sujet + Verbe / Adjectif & \zh{首先，我 主修 文科。其次，我 通过 了 HSK 3 级。} \\
\hline
Éducation (Antéchronologique) & \zh{时间 + 在 + Lieu/École + 学习 / 就读 + 专业} & \zh{2021 年 至 2024 年 在 雅温得 勒克莱尔 高中 就读，主修 文科。} \\
\hline
Expérience / Stage & \zh{时间 + 在 + Structure + 担任 / 实习 + Poste，负责 + Tâches} & \zh{2023 年 暑假 在 「中非友谊」 文化中心 实习，负责 接待 与 资料 整理。} \\
\hline
Compétences (Maîtrise) & \zh{熟练 掌握 / 擅长 / 精通 + Nom (logiciel, langue, sport)} & \zh{熟练 掌握 Microsoft Office；擅长 足球 与 摄影。} \\
\hline
Niveau de langue & \zh{语言 + 程度} : \zh{母语 / 流利 / 良好 / 基础 / HSK X 级} & \zh{法语 母语，英语 流利，汉语 HSK 3 级 (备考 HSK 4)。} \\
\hline
Motivation (Oral / Auto-évaluation) & \zh{我想 去 中国 学习，因为... / 为了...} & \zh{我想 去 中国 攻读 国际 关系，因为 致力 于 中喀 合作。} \\
\hline
\end{tabularx}
\end{center}
\end{propriete}

\begin{popfigure}
\begin{tikzpicture}[
    mindmap,
    grow cyclic,
    every node/.style={concept, execute at begin node=\hskip0pt},
    concept color=popblue!80,
    level 1/.append style={level distance=4.5cm, sibling angle=90, concept color=popgreen!70},
    level 2/.append style={level distance=3cm, sibling angle=45, concept color=poporange!70},
    text=white,
    font=\small\bfseries
]
\node[align=center]{CV Chinois\\ \zh{简历}}
    child { node[align=center]{Infos Perso\\ \zh{个人信息}}
        child { node {Nom \zh{姓名}} }
        child { node {Contact \zh{联系方式}} }
        child { node {Lycée \zh{就读学校}} }
    }
    child { node[align=center]{Formation\\ \zh{教育背景}}
        child { node {2021--2024 \zh{时间}} }
        child { node {Lycée Leclerc \zh{学校}} }
        child { node {Série A \zh{文科班}} }
    }
    child { node[align=center]{Compétences\\ \zh{技能特长}}
        child { node {Informatique \zh{办公软件}} }
        child { node {Langues \zh{语言能力}} }
        child { node {Sport/Art \zh{体育/艺术}} }
    }
    child { node[align=center]{Expérience\\ \zh{经验/实践}}
        child { node {Bénévolat \zh{志愿者}} }
        child { node {Job été \zh{兼职}} }
        child { node {Projets \zh{项目经历}} }
    };
\end{tikzpicture}
\legende{Structure mentale d'un CV chinois standard (antéchronologique).}
\end{popfigure}

\begin{exoresolu}{CV Modèle — Élève fictif : Jean-Pierre M. (Terminale A, Lycée Leclerc, Yaoundé)}
\zh{个人 简历} \\
\zh{个人信息} \\
\zh{姓名：马·让·皮埃尔 (Mǎ Ràng Pí'ěr)} \\
\zh{性别：男 \quad 出生年月：2006 年 5 月 12 日} \\
\zh{籍贯：喀麦隆 杜阿拉 \quad 现居住：雅温得} \\
\zh{联系电话：+237 6XX XX XX XX \quad 邮箱：jp.mbarga@email.cm} \\
\zh{就读学校：雅温得 勒克莱尔 高中 \quad 年级：高三 (A) 班} \\
\tcblower
\textbf{Commentaire détaillé du corrigé (pour l'enseignant / auto-correction élève)} \\

\textbf{1. Rubrique \zh{个人信息} (Informations personnelles)} \\
\emph{Choix linguistiques :} Le nom est transcrit phonétiquement (\zh{马} Mǎ pour Mbarga, \zh{让·皮埃尔} pour Jean-Pierre, point médian `·` standard). \zh{籍贯} (jíguàn) = lieu d'origine ancestrale (Douala), \zh{现居住} (xiàn jūzhù) = résidence actuelle (Yaoundé). Structure nominale standard, sans verbe. \zh{高三 (A) 班} (Gāosān A bān) rend « Terminale A » (année terminale, série A).

\textbf{2. Rubrique \zh{教育背景} (Formation)} \\
\zh{教育背景} \\
\zh{2021 年 9 月 -- 2024 年 6 月 \quad 雅温得 勒克莱尔 高中 \quad 文科班} \\
\zh{主修课程：法语、文学、历史、地理、数学、哲学} \\
\zh{班级 排名：前 10\% \quad 获奖：2023 年 校级 「优秀 学生」 称号} \\
\emph{Analyse :} Format antéchronologique. \zh{主修课程} (cours principaux). \zh{班级排名} (classement). \zh{校级} (niveau établissement). \zh{称号} (titre/récompense). \textbf{Caractères cibles :} \zh{教} (dans \zh{教育}, \zh{教学}), \zh{育} (dans \zh{教育}, \zh{培育}).

\textbf{3. Rubrique \zh{语言能力} (Compétences linguistiques)} \\
\zh{语言能力} \\
\zh{法语：母语} \\
\zh{英语：流利 (已通过 剑桥 FCE / 校内 高级)} \\
\zh{汉语：HSK 3 级 (2024 年 3 月 获得)，正在 备考 HSK 4 级} \\
\emph{Analyse :} Vocabulaire précis : \zh{母语}, \zh{流利}, \zh{通过} (réussir un examen), \zh{获得} (obtenir), \zh{备考} (préparer un examen). \textbf{Caractère cible :} \zh{能} (dans \zh{能力}, \zh{能够}).

\textbf{4. Rubrique \zh{技能特长} (Compétences techniques \& autres)} \\
\zh{技能特长} \\
\zh{办公软件：熟练 掌握 Microsoft Word、Excel、PowerPoint} \\
\zh{打字：中文 拼音 输入 (60 字/分)、法文 盲打} \\
\zh{体育：擅长 足球 (校队 前锋)，长跑} \\
\zh{其他：持有 驾照 (B 类)，会做 喀麦隆 传统 菜 (恩多莱、椰子鸡)} \\
\emph{Analyse :} \zh{熟练掌握} (maîtriser couramment), \zh{擅长} (exceller dans). \zh{盲打} (mángdǎ = frappe au clavier sans regarder). \zh{前锋} (qiánfēng = avant-centre). \textbf{Caractères cibles :} \zh{技} (dans \zh{技能}, \zh{技术}), \zh{能} (dans \zh{技能}, \zh{能力}).

\textbf{5. Rubrique \zh{社会实践 与 实习经验} (Pratique sociale \& Expérience)} \\
\zh{社会实践 与 实习经验} \\
\zh{2023 年 7 月 -- 8 月 \quad 雅温得 「中非友谊」 文化中心 \quad 志愿者} \\
\zh{协助 组织 「中国 文化 周」 活动，负责 接待 来宾、资料 整理 及 中法 口语 交流} \\
\zh{2022 年 12 月 \quad 杜阿拉 家庭 餐厅 \quad 服务员 (兼职)} \\
\zh{锻炼 了 服务 意识 与 团队 合作 能力} \\
\emph{Analyse :} \zh{社会实践} (pratique sociale/bénévolat) très valorisé en Chine. \zh{志愿者} (bénévole). \zh{协助组织} (assister à l'organisation). \zh{接待} (accueillir), \zh{资料整理} (classement de documents), \zh{口语交流} (échange oral). \zh{锻炼了...能力} (a exercé/développé la capacité de...). \textbf{Caractères cibles :} \zh{验} (dans \zh{经验}, \zh{检验}), \zh{历} (dans \zh{经历}, \zh{履历}).

\textbf{6. Rubrique \zh{爱好 与 自我评价} (Loisirs \& Auto-évaluation)} \\
\zh{爱好} \\
\zh{阅读 中国 近代 史、喀麦隆 民间 故事；足球 (关注 上海 海港 / 山东 泰山)；摄影} \\
\zh{自我评价} \\
\zh{性格 开朗，吃苦耐劳，有 强烈 的 跨文化 交际 愿望。希望 获得 中国 政府 奖学金，赴华 攻读 国际 关系 或 汉语 国际 教育 专业，未来 致力 于 中喀 教育 合作。} \\
\emph{Analyse :} \zh{吃苦耐劳} (chīkǔ nàiláo = endurer la peine/travailler dur, chengyu très apprécié). \zh{跨文化交际} (communication interculturelle). \zh{攻读} (gōngdú = poursuivre des études supérieures). \zh{致力于} (zhìlì yú = s'engager à/dé vouer à). \textbf{Caractère cible :} \zh{聘} (contexte \zh{聘用} / \zh{招聘} dans l'appel, essentiel à reconnaître).
\end{exoresolu}

\courssub{Traduction : Thème et Version}
\begin{methode}[Méthode de traduction Thème / Version (Bac / HSK)]
\begin{enumerate}
    \item \textbf{Analyser la phrase source :} Identifier le sujet, le verbe, les compléments (temps, lieu, manière), la structure (SVO, sujet-commentaire, \zh{把}/\zh{被}).
    \item \textbf{Ordonner selon la syntaxe cible :} \textbf{Français $\rightarrow$ Chinois :} Temps -- Lieu -- Sujet -- Verbe -- Objet. \textbf{Chinois $\rightarrow$ Français :} Restituer les articles, les prépositions, les temps (passé/composé/imparfait selon \zh{了}/\zh{过}/contexte).
    \item \textbf{Traiter les mots-outils :} \zh{的} (N), \zh{得} (V/degré), \zh{地} (V/manière), \zh{了} (accompli/état), \zh{过} (expérience).
    \item \textbf{Vérifier les classificateurs :} \zh{个/位} (personnes), \zh{项/段} (expérience/projet), \zh{门} (langue/matière), \zh{种} (type).
    \item \textbf{Relire :} Ponctuation chinoise pleine chasse, pinyin mental pour les tons.
\end{enumerate}
\end{methode}

\begin{exercice}[Thème : Français $\rightarrow$ Chinois (Écrire en caractères + pinyin)]
\begin{enumerate}
    \item Je m'appelle Marie Koum, j'ai 18 ans, je suis élève au Lycée de Buea.
    \item De 2021 à 2024, j'ai étudié en série A (lettres) au Lycée Leclerc à Yaoundé.
    \item Je parle couramment le français (langue maternelle) et l'anglais, et j'ai le niveau HSK 3 en chinois.
    \item Je maîtrise bien Word et Excel ; je sais faire la cuisine camerounaise (ndolé).
    \item L'été dernier, j'ai fait un stage au Centre Culturel « Amitié Sino-Camerounaise » à Douala.
    \item Je souhaite obtenir la bourse CSC pour étudier les relations internationales en Chine.
\end{enumerate}
\end{exercice}

\begin{exercice}[Version : Chinois $\rightarrow$ Français (Traduire en français naturel)]
\begin{enumerate}
    \item \zh{我 叫 保罗·阿约，今年 十九 岁，来自 喀麦隆 杜阿拉。}
    \item \zh{2022 年 至 2025 年 在 雅温得 政府 高中 就读，主修 理科。}
    \item \zh{法语 母语，英语 良好，汉语 HSK 4 级 (二零二四 年 六月 通过)。}
    \item \zh{熟练 掌握 Python 编程，擅长 篮球，曾 任 班长 一 年。}
    \item \zh{二零二三 年 暑假 在 北京 餐厅 兼职，锻炼 了 沟通 能力。}
    \item \zh{我 希望 赴华 攻读 计算机 科学，为 中喀 科技 合作 做 贡献。}
\end{enumerate}
\end{exercice}

\begin{corrige}[Exercices Thème \& Version]
\textbf{Thème (Propositions de rédaction) :}
\begin{enumerate}
    \item \zh{我 叫 马丽·库姆 (Mǎlì Kùmǔ)，今年 十八 岁，在 布埃亚 高中 就读。} (Wǒ jiào Mǎlì Kùmǔ, jīnnián shíbā suì, zài Bù'āiyà Gāozhōng jiùdú.)
    \item \zh{2021 年 至 2024 年 在 雅温得 勒克莱尔 高中 就读，主修 文科 (A 班)。} (2021 nián zhì 2024 nián zài Yǎwēndé Lèikèlái Gāozhōng jiùdú, zhǔxiū wénkē A bān.)
    \item \zh{法语 是 母语，英语 流利，汉语 HSK 3 级。} (Fǎyǔ shì mǔyǔ, Yīngyǔ liúlì, Hànyǔ HSK 3 jí.)
    \item \zh{熟练 掌握 Word 和 Excel；擅长 做 喀麦隆 菜 (如 恩多莱)。} (Shúliàn zhǎngwò Word hé Excel; shàncháng zuò Kāmàilóng cài (rú Ēnduólái).)
    \item \zh{去年 暑假 在 杜阿拉 「中喀 友谊」 文化中心 实习。} (Qùnián shǔjià zài Dūālā "Zhōng-Kā Yǒuyì" Wénhuà Zhōngxīn shíxí.)
    \item \zh{我 想 申请 CSC 奖学金，去 中国 攻读 国际 关系 专业。} (Wǒ xiǎng shēnqǐng CSC jiǎngxuéjīn, qù Zhōngguó gōngdú guójì guānxì zhuānyè.)
\end{enumerate}

\textbf{Version :}
\begin{enumerate}
    \item Je m'appelle Paul Ayot, j'ai 19 ans, je viens de Douala (Cameroun).
    \item De 2022 à 2025, j'ai étudié au Lycée Gouvernemental de Yaoundé, en série scientifique (S).
    \item Français langue maternelle, anglais bon niveau, chinois HSK 4 (réussi en juin 2024).
    \item Maîtrise courante de la programmation Python, bon au basket, a été délégué de classe pendant un an.
    \item L'été 2023, job d'été dans un restaurant pékinois, a exercé ma capacité de communication.
    \item Je souhaite aller en Chine pour faire un master en informatique, afin de contribuer à la coopération technologique sino-camerounaise.
\end{enumerate}
\end{corrige}

\courssub{Grille d'évaluation et production autonome}
\courssub{Objectifs de l'activité}
\begin{itemize}
    \item Mobiliser le vocabulaire spécialisé du CV (rubriques, diplômes, compétences, expériences) et les \cle{caractères complexes} étudiés (\zh{简历}, \zh{经验}, \zh{教育}, \zh{技能}, \zh{聘}) dans une production écrite autonome.
    \item Structurer un texte cohérent en respectant les conventions du CV chinois (ordre antéchronologique, rubriques standardisées, phrases nominales ou verbales concises).
    \item Utiliser les \cle{connecteurs logiques} (\zh{首先}, \zh{其次}, \zh{此外}, \zh{最后}) pour organiser l'information.
    \item Développer la compétence \textbf{expression orale} : présenter son parcours face à un jury simulé (Ambassade de Chine à Yaoundé) et répondre à une question simple.
    \item Acquérir une démarche d'\textbf{auto-évaluation et d'évaluation par les pairs} à l'aide d'une grille critériée sur 20 points.
\end{itemize}

\courssub{Consignes}
\textbf{Étape 1 — Rédaction individuelle (30 min)}
\begin{enumerate}
    \item \textbf{Analyse du modèle :} Relisez le CV commenté (rubriques \zh{个人信息}, \zh{教育背景}, \zh{工作/实习经验}, \zh{技能特长}, \zh{语言能力}, \zh{爱好}). Identifiez les \cle{caractères cibles} et les \cle{connecteurs}.
    \item \textbf{Remplissage de la fiche « Mon profil » (brouillon) :} Complétez en chinois (caractères + pinyin) vos informations réelles ou réalistes (élève de Terminale A, Lycée de Yaoundé/Douala/Buea).
    \item \textbf{Rédaction du CV final (6 à 8 phrases / items) :} Sur la feuille de copie fournie (format A4, modèle vierge avec rubriques pré-imprimées), écrivez votre \zh{简历} en \textbf{caractères chinois} (soin calligraphique exigé pour les 8 caractères complexes). Respectez la structure :
    \begin{itemize}
        \item \zh{个人信息} : Nom, âge, lycée, ville (Yaoundé/Douala...), contact.
        \item \zh{教育背景} : \zh{2021--2024 就读于 喀麦隆 雅温得 某 高中，主修 文科} (exemple).
        \item \zh{语言能力} : \zh{法语 母语，英语 流利，汉语 HSK 3 级 (或 正在备考)}.
        \item \zh{技能特长} : Informatique (Word, Excel), sport, musique, etc. Utilisez \zh{擅长} / \zh{熟练掌握}.
        \item \zh{爱好} : \zh{喜欢 足球，阅读 中国 历史，做 喀麦隆 菜 (如 恩多莱、椰子鸡)}.
    \end{itemize}
    \item \textbf{Relecture :} Vérifiez l'ordre des traits des 8 caractères complexes, l'accord des classificateurs (\zh{个}, \zh{门}, \zh{项}, \zh{段}), la place des compléments (temps/lieu avant le verbe).
\end{enumerate}

\textbf{Étape 2 — Évaluation croisée par les pairs (15 min)}
\begin{enumerate}
    \setcounter{enumi}{4}
    \item Échangez votre CV avec votre binôme. Utilisez la \textbf{Grille d'évaluation sur 20} (distribuée / affichée) pour noter le CV de votre camarade. Écrivez un commentaire constructif en français ou en chinois simple (\zh{字写得很工整}, \zh{词汇 丰富}, \zh{注意 「经验」 的 写法}).
    \item Rendez le CV annoté à son auteur. Prenez 2 minutes pour lire la note et le commentaire.
\end{enumerate}

\textbf{Étape 3 — Simulation orale : « Jury de l'Ambassade » (25 min)}
\begin{enumerate}
    \setcounter{enumi}{6}
    \item Trois volontaires viennent au tableau. La classe devient le \textbf{Jury de sélection CSC} (présidé par l'enseignant / un élève « Ambassadeur »).
    \item Chaque candidat \textbf{présente son CV à l'oral} (1 min 30 max) : \zh{各位 评委 好，我叫... 今年... 岁，来自 喀麦隆 雅温得。我的 优势 是...} (utiliser \zh{首先}...\zh{其次}...\zh{最后}).
    \item Le Jury pose \textbf{une question simple} par candidat (ex. \zh{你会说什么语言？}, \zh{你为什么 想 去 中国 学习？}, \zh{你的 爱好 是 什么？}).
    \item Le Jury délibère (vote à main levée) et désigne le \textbf{« Meilleur CV 2025 »}. Le texte gagnant est affiché au « Mur de la Réussite » de la classe.
\end{enumerate}

\courssub{Ressources à mobiliser}
\begin{definition}[Rappel synthétique : Les 8 caractères complexes cibles (Ordre des traits \& Radicaux)]
\begin{center}
\begin{tabularx}{\linewidth}{|X|X|X|X|}
\hline
\rowcolor{popblueL} \textbf{Caractère} & \textbf{Pinyin} & \textbf{Radical (Clé) \& Nb traits} & \textbf{Piège principal} \\
\hline
\zh{简} & jiǎn & \zh{竹} (6) + 12 = 18 & $\neq$ \zh{检} (bois \zh{木}) \\
\hline
\zh{历} & lì & \zh{止} (4) + 5 = 9 & Simplifié de \zh{曆} \\
\hline
\zh{验} & yàn & \zh{马} (10) + 8 = 18 & $\neq$ \zh{骗} (tromper) \\
\hline
\zh{教} & jiào & \zh{攵} (4) + 7 = 11 & Bas \zh{攵} : point, horiz., crochet, point \\
\hline
\zh{育} & yù & \zh{月} (4) + 9 = 13 & $\neq$ \zh{毓} (traditionnel) \\
\hline
\zh{技} & jì & \zh{扌} (3) + 4 = 7 & Phonétique \zh{枝} sans bois \\
\hline
\zh{能} & néng & \zh{月} (4) + 6 = 10 & Anciennement « ours » \\
\hline
\zh{聘} & pìn & \zh{耳} (6) + 7 = 13 & Droite \zh{并} $\neq$ \zh{关} \\
\hline
\end{tabularx}
\end{center}
\end{definition}

\begin{attention}[Erreurs fréquentes des francophones — Check-list avant remise]
\begin{itemize}
    \item \textbf{Ordre des mots (Temps/Lieu) :} \emph{Faux :} \zh{我 在 2021 年 就读 高中。} $\rightarrow$ \textbf{Correct :} \zh{2021 年 我 在 高中 就读。} (Temps en tout début ou après sujet).
    \item \textbf{Classificateurs :} \zh{一 \textbf{个} 经验} (\textbf{Faux}) $\rightarrow$ \zh{一 \textbf{项} 经验} / \zh{一 \textbf{段} 经历}. \zh{一 \textbf{门} 语言} (pas \zh{个}). \zh{一 \textbf{位} 老师} (respect).
    \item \textbf{\zh{了} vs \zh{过} :} CV = faits accomplis, expérience vécue $\rightarrow$ \zh{去 \textbf{过} 中国} (expérience de vie), \zh{学 \textbf{了} 三 年 汉语} (action terminée avec durée). Pas de \zh{了} en fin de phrase isolée sur une ligne de CV.
    \item \textbf{\zh{的} / \zh{得} / \zh{地} :} \zh{中文 写 \textbf{得} 好} (complément de degré), \zh{认真 \textbf{地} 学习} (adverbial), \zh{我的 \textbf{的} 简历} (possession).
    \item \textbf{Caractères proches :} \zh{经验} (\zh{验} cheval) $\neq$ \zh{骗} (tromper). \zh{教育} (\zh{教}+\zh{育}) $\neq$ \zh{教养} (éducation familiale). \zh{并} (ensemble) $\neq$ \zh{关} (fermer).
\end{itemize}
\end{attention}

\begin{aretenir}[Bilan personnel — À coller dans le cahier de chinois]
\begin{itemize}
    \item \textbf{Ce que je sais faire maintenant :} Rédiger un \zh{简历} structuré en chinois (rubriques, vocabulaire spécifique, connecteurs \zh{首先/其次/此外/最后}). Écrire correctement les 8 caractères complexes \zh{简历经验教育技能聘} (ordre des traits, radicaux). Me présenter oralement en 1 min 30 face à un jury. Traduire des phrases Thème/Version sur le thème.
    \item \textbf{Points de vigilance pour le Bac / HSK :} Respecter l'ordre « Temps -- Lieu -- Verbe ». Choisir le bon classificateur (\zh{项/段/门/位}). Distinguer \zh{的/得/地}. Ne pas confondre \zh{验/骗}, \zh{教/孝}, \zh{并/关}, \zh{简/检}.
    \item \textbf{Mon score auto-évalué (sur 20) :} \underline{\hspace{3cm}} \quad \textbf{Commentaire du binôme :} \underline{\hspace{5cm}}
    \item \textbf{Objectif pour la prochaine fois :} \underline{\hspace{8cm}}
\end{itemize}
\end{aretenir}

\begin{savaistu}[Le saviez-vous ? Le CV en Chine : particularités culturelles]
En Chine, le CV (\zh{简历}) standard pour un étudiant ou jeune diplômé comporte souvent une \textbf{photo d'identité} (formelle, fond bleu/rouge) en haut à droite, l'\textbf{âge} (ou date de naissance), le \textbf{sexe}, l'\textbf{ethnie} (\zh{民族}, ex. \zh{汉族} Hànzú) et le \textbf{statut politique} (\zh{政治面貌} : \zh{共青团员} / \zh{中共党员} / \zh{群众}). Ces informations, interdites ou déconseillées en France/Cameroun (loi anti-discrimination), sont courantes en Chine pour les administrations et grandes entreprises. Pour une bourse CSC, on exige le \zh{政治面貌}. En revanche, la rubrique \zh{爱好} est prise très au sérieux : elle révèle la personnalité et l'équilibre du candidat. Mentionner un sport d'équipe (football \zh{足球}, basket \zh{篮球}) ou un art traditionnel (calligraphie \zh{书法}, peinture \zh{国画}, thé \zh{茶艺}) est un atout majeur.
\end{savaistu}

\newpage

\coursec{Méthodes et exercices résolus}

\coursec{Leçon 40 — Expression écrite : caractères complexes liés à la rédaction d'un CV}

\courssub{Découverte : le CV chinois et ses rubriques}

\begin{textebox}[Modèle de CV chinois standard — Étudiant Terminale]
\textbf{个人简历} (Gèrén Jiǎnlì) \\ 
\noindent\rule{\linewidth}{0.4pt}\par
\textbf{个人信息} (Gèrén Xìnxī) \\
姓名：李华 \quad 性别：男 \quad 出生年月：2006年9月 \\
电话：+237 6XX XX XX XX \quad 邮箱：lihua@email.cm \\
地址：喀麦隆，布埃亚 \\ 
\noindent\rule{\linewidth}{0.4pt}\par
\textbf{教育经历} (Jiàoyù Jīnglì) \\
2023--2026 \quad \textbf{布埃亚第一高中} (Bù'āiyà Dì-yī Gāozhōng) — Terminale A \\
\quad \textit{Spécialité : Littérature, mention Bien au Probatoire 2024} \\
2020--2023 \quad \textbf{布埃亚双语中学} (Bù'āiyà Shuāngyǔ Zhōngxué) \\
\noindent\rule{\linewidth}{0.4pt}\par
\textbf{工作经验} (Gōngzuò Jīngyàn) \\
2024.07--2024.08 \quad \textbf{喀麦隆电信公司} (Kāmàilóng Diànxìn Gōngsī) — Stagiaire \\
\quad \textit{Missions : accueil clientèle, saisie de données, initiation au chinois des affaires} \\
\noindent\rule{\linewidth}{0.4pt}\par
\textbf{技能} (Jìnéng) \\
语言：法语 (母语), 英语 (流利), 汉语 (HSK 3) \\
Informatique : Pack Office, Python (bases), WeChat Work \\
Permis : Permis B (voiture) \\
\noindent\rule{\linewidth}{0.4pt}\par
\textbf{兴趣爱好} (Xìngqù Àihào) \\
足球 (队长), 书法, 非洲鼓, 讲故事比赛 \\
\noindent\rule{\linewidth}{0.4pt}\par
\end{textebox}

\begin{definition}[Les cinq rubriques obligatoires d'un CV chinois]
Un CV (简历 jiǎnlì) se structure en blocs fixes, dans cet ordre :
\begin{enumerate}
\item \textbf{个人信息} (gèrén xìnxī) — Identité, contact, nationalité, photo (souvent exigée en Chine).
\item \textbf{教育经历} (jiàoyù jīnglì) — Formation, du plus récent au plus ancien (ordre antéchronologique).
\item \textbf{工作经验} (gōngzuò jīngyàn) — Expériences pro / stages, idem ordre antéchronologique.
\item \textbf{技能} (jìnéng) — Langues (niveau HSK/CECRL), informatique, permis, certifications.
\item \textbf{兴趣爱好} (xìngqù àihào) — Centres d'intérêt (sport, arts, vie associative).
\end{enumerate}
\end{definition}

\begin{center}
\begin{tabularx}{\linewidth}{|X|X|X|X|}
\hline \rowcolor{popblueL} Caractères & Pinyin & Nature & Traduction \\ \hline
简历 & jiǎnlì & nom & CV, curriculum vitæ \\ \hline
个人信息 & gèrén xìnxī & nom & informations personnelles \\ \hline
教育经历 & jiàoyù jīnglì & nom & parcours éducatif, formation \\ \hline
工作经验 & gōngzuò jīngyàn & nom & expérience professionnelle \\ \hline
技能 & jìnéng & nom & compétences, aptitudes \\ \hline
兴趣爱好 & xìngqù àihào & nom & centres d'intérêt, loisirs \\ \hline
招聘 & zhāopìn & verbe & recruter, embaucher \\ \hline
应聘 & yìngpìn & verbe & postuler, répondre à une offre \\ \hline
面试 & miànshì & verbe/nom & passer un entretien / entretien d'embauche \\ \hline
录用 & lùyòng & verbe & admettre, engager (après sélection) \\ \hline
\end{tabularx}
\end{center}

\courssub{Écriture des caractères complexes (1) : 简历, 经验}

\begin{methode}[Décomposer, écrire, mémoriser les caractères du CV]
\begin{enumerate}
\item \cle{Analyser la structure} : identifier la \cle{clé (radical)} — indice de sens — et la \cle{partie phonétique} — indice de prononciation.
\item \cle{Respecter l'ordre des traits} (règles universelles) : 
\begin{itemize}
\item Haut $\rightarrow$ bas (上 $\rightarrow$ 下) ; Gauche $\rightarrow$ droite (左 $\rightarrow$ 右).
\item Horizontal avant vertical (横 $\rightarrow$ 竖) ; Extérieur avant intérieur ; Fermer la boîte en dernier (ex. 国, 园).
\item Point final en haut à gauche s'il est isolé (ex. 门, 斗).
\end{itemize}
\item \cle{S'écrire le mot complet} : en chinois, le sens stable est au niveau du mot (dissyllabique le plus souvent). Apprenez \zh{简历}, \zh{经验}, \zh{教育}, \zh{技能}, \zh{招聘} comme des blocs inséparables.
\end{enumerate}
\end{methode}

\begin{propriete}[Analyse graphique : 简历 jiǎnlì « CV »]
\begin{itemize}
\item \textbf{简 (jiǎn, simple, abrégé)} : Clé \textbf{⺮ (zhú, bambou, 6 traits)} + phonétique \textbf{间 (jiān)}. 
Ordre : traits du bambou (points, horizontales, crochet) $\rightarrow$ 间 (porte 日 + 日). 
\textit{Sens} : bambou utilisé pour faire des tablettes écrites simplifiées.
\item \textbf{历 (lì, traverser, calendrier, expérience)} : Clé \textbf{厂 (hǎn, falaise, 2 traits)} + \textbf{力 (lì, force)}. 
Ordre : 厂 (horizontal, pliure verticale) $\rightarrow$ 力 (horizontal, crochet descendant). 
\textit{Attention} : ne pas confondre avec \zh{厉 (lì, sévère)} — même son, clé 厂 mais 冎 (jiǒng) à l'intérieur.
\end{itemize}
\end{propriete}

\begin{propriete}[Analyse graphique : 经验 jīngyàn « expérience »]
\begin{itemize}
\item \textbf{经 (jīng, trame, passer par, gérer)} : Clé \textbf{纟 (sī, soie, 3 traits : diagonale, diagonale, horizontale)} + phonétique \textbf{圣 (shèng)}. 
\textit{Piège} : la clé 纟 s'écrit en \textbf{trois traits distincts}, jamais en un seul geste.
\item \textbf{验 (yàn, examiner, vérifier)} : Clé \textbf{马 (mǎ, cheval, 3 traits : horizontale, pliure, point)} + phonétique \textbf{佥 (qiān)} à droite. 
Forme simplifiée : 马 (gauche) + 佥 (droite). Forme traditionnelle : 驗 (cheval en bas).
\end{itemize}
\end{propriete}

\begin{exoresolu}[Correction d'erreurs fréquentes sur 简历 et 经验]
Énoncé. Corrigez les mots soulignés, nommez la clé fautive et donnez le pinyin exact.
\begin{enumerate}
\item L'élève écrit : « Mon \underline{简厉} est prêt. »
\item L'élève écrit : « J'ai beaucoup d'\underline{经马}. »
\end{enumerate}
\tcblower
Solution détaillée.
\begin{enumerate}
\item \textbf{简厉 $\rightarrow$ 简历 (jiǎnlì)}. Erreur : confusion homophone \zh{厉 (lì, sévère)} / \zh{历 (lì, parcourir)}. La clé de \zh{历} est \textbf{厂} (falaise) avec \textbf{力} dedans ; \zh{厉} a \textbf{厂} + \textbf{冎}. Dans « CV », c'est bien \zh{历} (parcours).
\item \textbf{经马 $\rightarrow$ 经验 (jīngyàn)}. Erreur : caractère incomplet. \zh{验} = clé \textbf{马} (cheval, à gauche) + \textbf{佥} (à droite). L'élève n'a écrit que la clé. Rappel : 纟 (soie, 3 traits) pour 经 ; 马 (cheval, 3 traits) pour 验.
\end{enumerate}
\end{exoresolu}

\courssub{Écriture des caractères complexes (2) : 教育, 技能, 聘}

\begin{propriete}[Analyse graphique : 教育 jiàoyù « éducation / formation »]
\begin{itemize}
\item \textbf{教 (jiāo, enseigner)} : Clé \textbf{攵 (pū, frapper/main tenant bâton, 4 traits)} + phonétique \textbf{孝 (xiào, piété filiale)}. 
Ordre : 孝 (haut) $\rightarrow$ 攵 (bas : point, horizontale, diagonale, point final). 
\textit{Distinction} : \zh{教 (jiāo)} $\neq$ \zh{数 (shù, nombre)} — clé de 数 est \textbf{攵} aussi mais phonétique \textbf{娄 (lóu)}.
\item \textbf{育 (yù, élever, éduquer)} : Clé \textbf{月 (yuè, chair/corps)} + \textbf{立 (lì, debout)} + \textbf{月} (bas). 
Structure : haut-milieu-bas. Ordre : 月 (haut) $\rightarrow$ 立 $\rightarrow$ 月 (bas).
\end{itemize}
\end{propriete}

\begin{propriete}[Analyse graphique : 技能 jìnéng « compétence, habileté »]
\begin{itemize}
\item \textbf{技 (jì, technique, art, habileté)} : Clé \textbf{扌 (shǒu, main, 3 traits)} + phonétique \textbf{支 (zhī) + 又 (yòu)}. 
Décomposition pratique : 扌 (gauche) $\rightarrow$ 十 $\rightarrow$ 一 $\rightarrow$ 一 (partie haute droite) $\rightarrow$ 又 (bas droite). 
\textit{Sens} : main qui manie une technique.
\item \textbf{能 (néng, pouvoir, capacité)} : Clé \textbf{月 (yuè, chair)} + \textbf{厶 (sī, privé)} + \textbf{匕 (bǐ, cuillère)} $\times$ 2. 
Ordre : 月 $\rightarrow$ 厶 $\rightarrow$ 匕 (gauche) $\rightarrow$ 匕 (droite). 
\textit{Note} : 月 clé « chair » (viande) s'écrit comme le caractère lune mais signifie « partie du corps » ici.
\end{itemize}
\end{propriete}

\begin{propriete}[Analyse graphique : 招聘 zhāopìn « recruter » et 聘 pìn « engager »]
\begin{itemize}
\item \textbf{招 (zhāo, appeler, recruter)} : Clé \textbf{扌 (main)} + phonétique \textbf{召 (zhào, convoquer)}. 
Ordre : 扌 $\rightarrow$ 召 (刀 + 口).
\item \textbf{聘 (pìn, engager, embaucher)} : Clé \textbf{耳 (ěr, oreille, 6 traits)} + phonétique \textbf{甹 (pìn)}. 
Ordre : 耳 (gauche, contour puis traits intérieurs) $\rightarrow$ 甹 (droite). 
\textit{Piège majeur} : \zh{聘 (pìn, oreille + 甹)} $\neq$ \zh{骋 (chěng, cheval + 甹, galoper)}. La clé change tout : oreille $\rightarrow$ écouter/embaucher ; cheval $\rightarrow$ courir.
\end{itemize}
\end{propriete}

\begin{center}
\begin{tabularx}{\linewidth}{|X|X|X|X|}
\hline \rowcolor{popgreenL} Caractères & Pinyin & Nature & Traduction \\ \hline
教育 & jiàoyù & nom/verbe & éducation, former \\ \hline
教 & jiāo & verbe & enseigner \\ \hline
育 & yù & verbe & élever, éduquer \\ \hline
技能 & jìnéng & nom & compétence, qualification \\ \hline
技 & jì & nom & technique, art, habileté \\ \hline
能 & néng & verbe/modal & pouvoir, être capable de \\ \hline
招聘 & zhāopìn & verbe & recruter \\ \hline
聘 & pìn & verbe & engager, embaucher \\ \hline
应聘 & yìngpìn & verbe & postuler \\ \hline
面试 & miànshì & verbe/nom & entretien (d'embauche) \\ \hline
\end{tabularx}
\end{center}

\begin{exoresolu}[Discrimination de caractères proches — Examen blanc]
Énoncé. Choisissez le bon caractère entre parenthèses et justifiez par la clé.
\begin{enumerate}
\item Cette entreprise veut (招聘 / 招骋) des ingénieurs.
\item Mon (教育 / 数育) a été fait à Yaoundé.
\item Il a beaucoup d'expérience : (经验 / 经马).
\end{enumerate}
\tcblower
Solution détaillée.
\begin{enumerate}
\item \textbf{招聘} (zhāopìn). Clé \textbf{耳} (oreille) pour \zh{聘} = embaucher (écouter les candidatures). \zh{骋} a la clé \textbf{马} (cheval) = galoper.
\item \textbf{教育} (jiàoyù). \zh{教} a la clé \textbf{攵} (action de la main) + \textbf{孝}. \zh{数} a la même clé mais phonétique \textbf{娄} = nombre.
\item \textbf{经验} (jīngyàn). \zh{验} = clé \textbf{马} + \textbf{佥} (examiner un cheval $\rightarrow$ vérifier). \zh{马} seul = cheval.
\end{enumerate}
\end{exoresolu}

\courssub{Modèle de CV commenté, structure et connecteurs}

\begin{textebox}[CV commenté — Li Hua (extrait rubriques clés)]
\textbf{教育经历} \\
2023--2026 \quad 布埃亚第一高中 \quad Terminale A \\
\quad 主修文学，2024年会考获良好成绩。 \\
\textbf{工作经验} \\
2024年7月--8月 \quad 喀麦隆电信公司 \quad 实习生 \\
\quad 负责客户接待、数据录入；协助中文商务培训。 \\
\textbf{技能} \\
语言：法语 (母语)、英语 (流利)、汉语 (HSK 3) \\
办公软件：Word, Excel, PowerPoint ; Python (基础) \\
\textbf{兴趣爱好} \\
足球 (校队队长)、中国书法、非洲鼓
\end{textebox}

\begin{propriete}[Structure de phrase nominale du CV (sans sujet répété)]
En chinois, les rubriques d'un CV s'écrivent en \cle{phrases nominales elliptiques} (sujet omis = « je ») :
\[
\textbf{Temps} + \textbf{Lieu (在 + lieu)} + \textbf{Verbe} + \textbf{Complément/Objet} (+ \textbf{了/过} si accompli)
\]
\begin{center}
\begin{tabularx}{\linewidth}{|X|X|X|}
\hline \rowcolor{poporangeL} Français & Chinois (caractères + pinyin) & Structure \\ \hline
J'ai étudié à Buea en 2023. & 2023年在布埃亚学习。 (2023 nián zài Bù'āiyà xuéxí.) & Temps + 在+Lieu + V \\ \hline
J'ai fait un stage chez Camtel. & 在喀麦隆电信公司实习过。 (Zài Kāmàilóng Diànxìn Gōngsī shíxí guo.) & 在+Lieu + V + 过 \\ \hline
Je sais parler chinois. & 会说汉语。 (Huì shuō Hànyǔ.) & Verbe modal + V + Objet \\ \hline
Diplômé de l'Univ. Yaoundé II. & 毕业于雅温得第二大学。 (Bìyè yú Yǎwēndé Dì-èr Dàxué.) & Verbe + 于 + Lieu \\ \hline
\end{tabularx}
\end{center}
\end{propriete}

\begin{definition}[Connecteurs utiles pour la lettre de motivation (求职信 qiúzhíxìn)]
\begin{center}
\begin{tabularx}{\linewidth}{|X|X|X|}
\hline \rowcolor{poppurpleL} Chinois & Pinyin & Français \\ \hline
首先 & shǒuxiān & d'abord, en premier lieu \\ \hline
然后 & ránhòu & ensuite, puis \\ \hline
此外 & cǐwài & de plus, par ailleurs \\ \hline
最后 & zuìhòu & enfin, pour terminer \\ \hline
因此 & yīncǐ & c'est pourquoi, par conséquent \\ \hline
希望 & xīwàng & espérer (que) \\ \hline
\end{tabularx}
\end{center}
\end{definition}

\begin{experience}[Atelier oral : « Mon CV en 2 minutes »]
\textbf{Consigne} : En binôme, présentez votre CV oralement en chinois (1 min 30 chacun). Utilisez les connecteurs \zh{首先、然后、此外、最后}. L'écouteur note les infos clés sur une grille (Nom, Formation, Stage, Compétences, Loisir) puis restitue en français.
\textbf{Modèle de départ} : \zh{首先，我叫……，我在……读书。然后，2024年我在……实习过。此外，我会……最后，我的爱好是……。}
\end{experience}

\courssub{Traduction : thème et version}

\begin{methode}[Réussir un thème (Français $\rightarrow$ Chinois) sur le CV]
\begin{enumerate}
\item \cle{Segmenter} la phrase française : [Sujet] [Verbe] [Compléments temps/lieu/objet].
\item \cle{Réordonner} selon le schéma chinois : \textbf{Sujet + Temps + 在+Lieu + Verbe + Objet + Durée/Particules (了/过)}.
\item \cle{Choisir le vocabulaire précis} : « postuler » = \zh{应聘 / 申请} ; « embaucher » = \zh{招聘 / 录用} ; « diplômé de » = \zh{毕业于} ; « stage » = \zh{实习}.
\item \cle{Insérer les classificateurs} : \zh{一所学校} (école), \zh{一家公司} (entreprise), \zh{一个职位} (poste).
\item \cle{Vérifier} : tons, ordre des mots, particules, ponctuation chinoise.
\end{enumerate}
\end{methode}

\begin{exoresolu}[Thème : trois phrases types Bac]
Énoncé. Traduisez en chinois (caractères + pinyin).
\begin{enumerate}
\item « Marie cherche un emploi à Yaoundé. »
\item « L'année dernière, j'ai travaillé six mois dans une entreprise de Douala. »
\item « Cette société recrute des jeunes qui parlent chinois. »
\end{enumerate}
\tcblower
Solution détaillée.
\begin{enumerate}
\item \textbf{玛丽在雅温得找工作。} (Mǎlì zài Yǎwēndé zhǎo gōngzuò.) 
\textit{Analyse} : Sujet 玛丽 + Lieu \zh{在雅温得} (avant verbe) + Verbe \zh{找工作}. Pas de « un » (emploi indéfini).
\item \textbf{去年我在杜阿拉的一家公司工作了六个月。} (Qùnián wǒ zài Dù'ālā de yì jiā gōngsī gōngzuò le liù gè yuè.) 
\textit{Analyse} : Temps \zh{去年} (tête) + Sujet \zh{我} + Lieu \zh{在杜阿拉的一家公司} (classif. \zh{家}) + Verbe \zh{工作} + \zh{了} (accompli) + Durée \zh{六个月} (classif. \zh{个} pour mois).
\item \textbf{这家公司招聘会说汉语的年轻人。} (Zhè jiā gōngsī zhāopìn huì shuō Hànyǔ de niánqīngrén.) 
\textit{Analyse} : Sujet \zh{这家公司} + Verbe \zh{招聘} + Objet = Relative \zh{会说汉语的} + Nom \zh{年轻人}. 
\textit{Règle d'or} : La relative en \zh{的} \textbf{précède} le nom (inverse du français).
\end{enumerate}
\end{exoresolu}

\begin{methode}[Réussir une version (Chinois $\rightarrow$ Français) sur extrait de CV]
\begin{enumerate}
\item \cle{Segmenter} en blocs : Sujet / Temps / Lieu / Verbe (+ 了/过) / Complément / Relative (\zh{的}+Nom).
\item \cle{Convertir les marqueurs d'aspect} : \zh{了} $\rightarrow$ passé composé / accompli ; \zh{过} $\rightarrow$ expérience vécue (passé composé ou imparfait selon contexte) ; Date explicite $\rightarrow$ passé composé.
\item \cle{Traduire la relative} \zh{...的 + Nom} par « qui/que ... » placé \textbf{après} le nom en français.
\item \cle{Soigner le français} : articles (le/la/les/un/une), accords, conjugaisons correctes, vocabulaire professionnel (stage, CDI, compétences).
\end{enumerate}
\end{methode}

\begin{exoresolu}[Version : extrait de CV complet]
Énoncé. Traduisez en français naturel :
\zh{我叫王伟，2022年毕业于北京大学。我在一家电信公司工作过两年，会说法语和英语。我想在喀麦隆找工作。}
\tcblower
Solution détaillée.
\begin{enumerate}
\item \textbf{Segmentation} : 
\zh{我叫王伟} / \zh{2022年毕业于北京大学} / \zh{我在一家电信公司工作过两年} / \zh{会说法语和英语} / \zh{我想在喀麦隆找工作}。
\item \textbf{Analyse grammaticale} :
\begin{itemize}
\item \zh{毕业于} + date $\rightarrow$ passé composé « j'ai obtenu mon diplôme de / j'ai été diplômé de ».
\item \zh{工作过} (\zh{过} = expérience) + durée \zh{两年} $\rightarrow$ « j'ai travaillé deux ans / pendant deux ans ».
\item \zh{会说} = « savoir parler » (compétence acquise).
\item \zh{想} + verbe = « vouloir / avoir l'intention de ».
\end{itemize}
\item \textbf{Traduction rédigée} :
« Je m'appelle Wang Wei. En 2022, j'ai obtenu mon diplôme de l'Université de Pékin. J'ai travaillé deux ans dans une entreprise de télécommunications, et je sais parler français et anglais. Je souhaite trouver un emploi au Cameroun. »
\item \textbf{Points de vigilance} : \zh{北京大学} = « Université de Pékin » (nom conventionnel) ; \zh{喀麦隆} = « le Cameroun » (article défini pays masculin) ; \zh{一家公司} = « une entreprise » (classificateur \zh{家} rendu par article indéfini).
\end{enumerate}
\end{exoresolu}

\courssub{Grille d'évaluation et production autonome}

\begin{aretenir}[Grille d'évaluation — Expression écrite CV (sur 10 pts)]
\begin{center}
\begin{tabularx}{\linewidth}{|X|X|X|}
\hline \rowcolor{popgoldL} Critère & Indicateurs & Points \\ \hline
\textbf{Respect du format} & 5 rubriques présentes, ordre logique, mise en page claire & 2 \\ \hline
\textbf{Vocabulaire spécialisé} & Mots-clés exacts (简历, 教育经历, 工作经验, 技能, 招聘, 应聘, 实习, 毕业于) & 2 \\ \hline
\textbf{Caractères complexes} & Écriture correcte des 15 caractères cibles (clés, traits, pas de confusion 历/厉, 聘/骋, 教/数) & 2 \\ \hline
\textbf{Grammaire / Syntaxe} & Ordre S + Temps + 在+Lieu + V + Compl ; Relative en \zh{的} avant nom ; Particules \zh{了/过/的} justes ; Classificateurs \zh{所/家/个} & 2 \\ \hline
\textbf{Cohérence / Ponctuation} & Phrases nominales elliptiques, tons notés en pinyin, ponctuation chinoise ，。 ； pas de calque français & 2 \\ \hline
\end{tabularx}
\end{center}
\end{aretenir}

\begin{exercice}[Production autonome — Épreuve type Bac (20 min)]
\textbf{Consigne} : Vous êtes \zh{张明} (Zhāng Míng), élève de Terminale A à Douala. Rédigez votre CV complet en chinois (caractères + pinyin) pour postuler à un stage d'été dans une entreprise sino-camerounaise de BTP. 
\textbf{Contraintes} :
\begin{itemize}
\item 5 rubriques obligatoires remplies (inventez des détails réalistes : lycée, dates, stage 2023, langues, sport).
\item Minimum 8 phrases nominales (segments) au total.
\item Utilisez au moins 3 connecteurs (首 先, 然后, 此外, 最后) dans une courte lettre d'accompagnement (求职信, 50--80 caractères) placée en haut.
\item Soignez les caractères complexes : \zh{简历, 经验, 教育, 技能, 招聘, 实习, 毕业于, 面试}.
\end{itemize}
\end{exercice}

\begin{corrige}[Exercice — Proposition de correction (exemple)]
\textbf{求职信} \\
首先，我叫张明，是杜阿拉第一高中终末年级学生。然后，我对贵公司的建筑项目很感兴趣。此外，我会说法语、英语和汉语。最后，希望获得面试机会。 \\
\noindent\rule{\linewidth}{0.4pt}\par
\textbf{个人简历} \\
\textbf{个人信息} \\
姓名：张明 \quad 性别：男 \quad 出生：2006年3月 \\
电话：+237 6XX XX XX XX \quad 邮箱：zhangming@email.cm \\
\textbf{教育经历} \\
2023--2026 \quad 杜阿拉第一高中 \quad Terminale A \\
\quad 2024年会考通过，主修科学数学。 \\
\textbf{工作经验} \\
2023年7月--8月 \quad 中喀建筑公司 \quad 实习生 \\
\quad 协助工程师整理图纸，现场翻译法语/汉语。 \\
\textbf{技能} \\
语言：法语 (母语), 英语 (B2), 汉语 (HSK 4) \\
Informatique : AutoCAD (基础), Office \\
Permis : Permis B \\
\textbf{兴趣爱好} \\
篮球 (校队), 中国书法, 阅读建筑杂志 \\
\noindent\rule{\linewidth}{0.4pt}\par
\textbf{Note} : Cette correction sert de modèle ; les détails personnels varient. Vérifiez chaque caractère cible avec la grille ci-dessus.
\end{corrige}

\begin{attention}[Les 5 erreurs qui coûtent des points au Bac — Récapitulatif]
\begin{enumerate}
\item \textbf{Ordre des mots calqué sur le français} : \zh{*我工作在杜阿拉} $\rightarrow$ \textbf{\zh{我在杜阿拉工作}} (Lieu \zh{在...} AVANT le verbe).
\item \textbf{Confusion caractères proches / homophones} : \zh{简历} $\neq$ \zh{简厉} (厉=严厉) ; \zh{招聘} $\neq$ \zh{招骋} (骋=驰骋) ; \zh{教育} $\neq$ \zh{数育} (数=数字) ; \zh{经验} $\neq$ \zh{经马} (马=cheval seul). \textit{Une clé fautive = mot nul}.
\item \textbf{Relative mal placée} : \zh{*年轻人会说汉语的} $\rightarrow$ \textbf{\zh{会说汉语的年轻人}} (\zh{的} + Nom). En chinois, le modificateur précède toujours le nom.
\item \textbf{Mots-outils mal employés} : \zh{过} seulement pour l'expérience vécue (\zh{去过、做过}) ; \zh{了} pour l'accompli ponctuel ou changement d'état ; \zh{的} pour relative/possession. Ne pas mettre \zh{了} après chaque verbe passé.
\item \textbf{Classificateurs et tons oubliés} : \zh{一所学校、一家公司、一个职位、六个月} (pas \zh{个} pour école/entreprise) ; Tons exacts sur mots-clés : \zh{jiǎnlì, jīngyàn, jiàoyù, jìnéng, zhāopìn, shíxí, bìyè yú, miànshì}. Un ton faux sur mot-clé = pénalité.
\end{enumerate}
\end{attention}

\begin{savaistu}[CV en Chine vs Cameroun : trois différences culturelles]
\begin{itemize}
\item \textbf{Photo d'identité} : En Chine, la photo (formelle, fond bleu/rouge) est \textbf{quasi obligatoire} sur le CV ; au Cameroun, elle est facultative (sauf hôtellerie, mannequinat).
\item \textbf{Informations personnelles} : En Chine, on indique souvent l'âge (出生年月), le lieu de naissance (籍贯), l'état civil, le parti politique (党员) ; au Cameroun (RGPD / droit du travail), âge, état civil, religion, ethnie sont \textbf{interdits} ou déconseillés pour éviter la discrimination.
\item \textbf{Recommandations / Guanxi} : En Chine, les \zh{关系} (guānxi, réseaux) et lettres de recommandation (推荐信) pèsent lourd ; au Cameroun, le diplôme et l'expérience priment officiellement, mais le « piston » (réseau informel) existe aussi.
\item \textbf{Format} : Chine : souvent une page A4 très structurée, parfois modèle officiel ; Cameroun : format libre, 1--2 pages, lettre de motivation séparée.
\end{itemize}
\end{savaistu}

\begin{popfigure}
\begin{tikzpicture}[
  node distance=1.2cm and 1.5cm,
  box/.style={rectangle, rounded corners, draw=popblue, thick, fill=popblueL, text width=2.8cm, align=center, font=\small},
  arrow/.style={-Stealth, thick, popdark}
]
\node[box, align=center] (sujet){Sujet \\ (我 / 姓名)};
\node[box, right=of sujet, align=center] (temps){Temps \\ (2024年 / 去年)};
\node[box, right=of temps, align=center] (lieu){Lieu \\ (在 + 杜阿拉 / 公司)};
\node[box, right=of lieu, align=center] (verbe){Verbe \\ (工作 / 实习 / 找工作)};
\node[box, right=of verbe, align=center] (comp){Complément \\ (Objet / Durée / 了 / 过)};
\draw[arrow] (sujet) -- (temps);
\draw[arrow] (temps) -- (lieu);
\draw[arrow] (lieu) -- (verbe);
\draw[arrow] (verbe) -- (comp);
\node[below=0.8cm of verbe, font=\bfseries, text=poporange] {Ordre canonique du CV chinois};
\end{tikzpicture}
\legende{Schéma de structure de phrase pour les rubriques « Expérience » et « Formation » : S + Temps + Lieu + V + Compl.}
\end{popfigure}

\newpage

\coursec{Exercices}

\courssub{Je vérifie mes connaissances}

\begin{exercice}[QCM : Vocabulaire et structure du CV]
Chaque question ne comporte qu'une seule réponse correcte. Choisis la bonne proposition.
\begin{enumerate}
    \item Quel terme chinois correspond à « CV / Curriculum Vitae » ?
    \begin{enumerate}
        \item 简历 (jiǎnlì)
        \item 简历 (jiǎnlì) \textbf{Note : les deux écritures sont identiques, c'est le terme standard.}
        \item 履历 (lǚlì)
        \item 简介 (jiǎnjiè)
    \end{enumerate}
    \item Dans un CV chinois, la rubrique « Expérience professionnelle » s'écrit généralement :
    \begin{enumerate}
        \item 工作经验 (gōngzuò jīngyàn)
        \item 学习经历 (xuéxí jīnglì)
        \item 爱好特长 (àihào tècháng)
        \item 个人信息 (gèrén xìnxī)
    \end{enumerate}
    \item Pour indiquer la période « de 2020 à 2024 », on utilise la structure :
    \begin{enumerate}
        \item 2020年 — 2024年
        \item 2020年到2024年
        \item 2020年至2024年
        \item Les trois propositions sont correctes selon le registre (formel/informel).
    \end{enumerate}
    \item Le caractère \zh{聘} (pìn) dans \zh{应聘} (yìngpìn) signifie :
    \begin{enumerate}
        \item Embaucher / Engager (verbe)
        \item Postuler (verbe)
        \item Entretien (nom)
        \item Salaire (nom)
    \end{enumerate}
\end{enumerate}
\end{exercice}

\begin{exercice}[Vrai ou Faux : Connaissance du CV chinois]
Indique si l'affirmation est vraie (V) ou fausse (F) et justifie ta réponse en une phrase en français.
\begin{enumerate}
    \item La photo d'identité est obligatoire et se place en bas à droite du CV chinois.
    \item La rubrique \zh{教育背景} (jiàoyù bèijǐng) ne concerne que les diplômes universitaires, pas le lycée.
    \item Dans la partie \zh{技能} (jìnéng), on mentionne obligatoirement le niveau HSK pour la langue chinoise.
    \item L'ordre chronologique inverse (du plus récent au plus ancien) est la norme pour lister les expériences.
    \item Le connecteur \zh{熟练掌握} (shúliàn zhǎngwò) s'utilise uniquement pour les langues étrangères.
\end{enumerate}
\end{exercice}

\begin{exercice}[Vocabulaire : Associer caractères, pinyin et sens]
Relie chaque caractère complexe du thème à son pinyin (avec tons) et à sa traduction. Écris les correspondances sous la forme : \textbf{Caractère — Pinyin — Traduction}.
\begin{center}
\begin{tabularx}{\linewidth}{|X|X|X|X|}
\hline
\rowcolor{popblueL} \textbf{Caractère} & \textbf{Pinyin} & \textbf{Nature} & \textbf{Traduction} \\ \hline
简 & jiǎn & adj. / verbe & simple ; simplifier ; choisir \\ \hline
历 & lì & nom / verbe & calendrier ; expérience ; passer par \\ \hline
经 & jīng & nom / verbe & trame ; passer par ; expérience \\ \hline
验 & yàn & nom / verbe & vérifier ; expérience ; test \\ \hline
教 & jiāo & verbe & enseigner ; éduquer \\ \hline
育 & yù & verbe & élever ; éduquer ; former \\ \hline
技 & jì & nom & technique ; habileté ; compétence \\ \hline
能 & néng & verbe / nom & pouvoir ; capacité ; compétence \\ \hline
聘 & pìn & verbe & engager ; embaucher ; inviter \\ \hline
应 & yìng & verbe & répondre ; postuler ; faire face \\ \hline
\end{tabularx}
\end{center}
\end{exercice}

\begin{exercice}[Pinyin vers Caractères : Reconnaissance auditive/visuelle]
Écris les caractères chinois simplifiés correspondant aux pinyin suivants (attention aux tons et aux homophones fréquents). Puis entoure le radical (clé) de chaque caractère composé.
\begin{enumerate}
    \item jiǎnlì \hspace{3cm} (CV)
    \item jīngyàn \hspace{3cm} (expérience)
    \item jiàoyù \hspace{3cm} (éducation)
    \item jìnéng \hspace{3cm} (compétence)
    \item yìngpìn \hspace{3cm} (postuler)
    \item gōngzuò \hspace{3cm} (travail)
    \item zìwo \hspace{3cm} (soi-même / auto-)
    \item pinggu \hspace{3cm} (évaluation)
    \item xiangmu \hspace{3cm} (projet)
    \item jiangli \hspace{3cm} (récompense)
\end{enumerate}
\end{exercice}

\courssub{Je m'entraîne}

\begin{exercice}[Traduction : Thème et Version (phrases isolées)]
Traduis les phrases suivantes. Respecte l'ordre des mots chinois (Sujet + Temps + Lieu + Verbe + Objet) et l'usage des mots-outils.
\begin{enumerate}
    \item \textbf{Thème :} Je m'appelle Paul. Je suis camerounais.
    \item \textbf{Thème :} De 2020 à 2024, j'ai étudié au Lycée de Bafoussam.
    \item \textbf{Thème :} J'ai de l'expérience de travail. J'ai travaillé dans une entreprise chinoise.
    \item \textbf{Version :} \zh{我的技能包括中文、法文和英文。}
    \item \textbf{Version :} \zh{我熟练掌握办公软件和编程基础。}
    \item \textbf{Thème :} Mes loisirs sont le football et la musique. Je suis en bonne santé.
    \item \textbf{Version :} \zh{应聘职位：中文助教。联系电话：6XXXXXXXX。}
    \item \textbf{Thème :} J'espère obtenir un entretien. Merci de votre attention.
\end{enumerate}
\end{exercice}

\begin{exercice}[Discrimination de caractères : Choix du bon graphème]
Complète les phrases en choisissant le bon caractère parmi la paire proposée. Écris le caractère correct dans l'espace vide.
\begin{enumerate}
    \item Mon CV est très \_\_\_\_\_ (简 / 间) 单, facile à lire.
    \item Il a une riche \_\_\_\_\_ (历 / 厉) 史 de travail en Chine.
    \item Cette expérience est très \_\_\_\_\_ (历 / 厉) 害, il a beaucoup appris.
    \item Ma \_\_\_\_\_ (技 / 枝) 能 principale est la programmation Python.
    \item Il a cassé une \_\_\_\_\_ (技 / 枝) de l'arbre devant l'école.
    \item Nous devons \_\_\_\_\_ (验 / 检) 查 产品质量.
    \item Le médecin va \_\_\_\_\_ (验 / 检) 查 我的身体.
\end{enumerate}
\end{exercice}

\begin{exercice}[Réorganisation de mots : Structure de phrase CV]
Remets les mots dans l'ordre correct pour former des phrases grammaticales adaptées à un CV. Écris la phrase complète en caractères, puis donne le pinyin et la traduction.
\begin{enumerate}
    \item 我 / 在 / 学习 / 雅温得中学 / 2020年到2024年
    \item 会说 / 我 / 中文 / 法文 / 和 / 英文
    \item 工作经验 / 有 / 两年 / 我 / 的
    \item 熟练 / 办公软件 / 使用 / 我 / 微软
    \item 足球 / 爱好 / 是 / 和 / 音乐 / 我的
\end{enumerate}
\end{exercice}

\begin{exercice}[Correction d'erreurs : CV d'un élève fictif]
Le texte ci-dessous est un extrait de CV rédigé par un élève de Terminale. Il contient \textbf{5 erreurs} (vocabulaire, grammaire, ordre des mots, caractères, registre). Identifie-les, souligne-les et propose la correction.
\begin{textebox}[CV de M. Étienne (avec erreurs)]
姓名： Étienne
国籍： 喀麦隆
教育经历：
2020-2024 雅温得中学 高中毕业
工作经历：
2022-2023 在 一家 餐厅 服务员 (错词序)
技能：
中文 HSK3， 法文 母语， 英文 好
爱好： 看书， 打篮球， 听音乐
自我评价： 我 很 努力， 希望 给 我 一个 机会。
\end{textebox}
\end{exercice}

\courssub{Je me prépare au Bac}

\begin{exercice}[Compréhension de l'écrit : Lettre de motivation + CV (Type Bac)]
Lis le texte suivant (lettre de motivation simplifiée accompagnant un CV), réponds aux questions en français et traite les exercices de langue.
\begin{textebox}[Lettre de motivation pour un stage d'été]
尊敬的 经理 先生/女士：

您好！

我叫 马可 (Mǎkě)， 今年 十九 岁， 是 喀麦隆 雅温得 第一 大学 的 大二 学生， 专业 是 国际 贸易。我 很 希望 能 到 贵公司 实习 两 个月 (七月 到 八月)。

在 学校， 我 系统 学习 了 国际 贸易 理论 和 商务 中文。我 通过 了 HSK 4 级 考试， 能 用 中文 进行 日常 交流 和 简单 的 商务 沟通。此外， 我 熟练 使用 Excel 和 PowerPoint， 有 组织 校园 活动 的 经验， 做事 认真 负责， 学习 能力 强。

随信 附上 我的 简历。期待 您 的 回复。祝 工作 顺利！

此致
敬礼

马可
2025年6月1日
\end{textebox}

\textbf{Questions de compréhension (en français) :}
\begin{enumerate}
    \item Qui est l'expéditeur ? Quel est son âge, sa nationalité, son lieu d'études et sa spécialité ?
    \item Quel est l'objet de cette lettre ? Quelle est la durée et la période souhaitées ?
    \item Quels sont les trois atouts principaux que Marco met en avant pour justifier sa candidature (langues, informatique, expérience) ?
    \item Quelle formule de politesse de fin de lettre est utilisée ici ? Traduisez-la.
\end{enumerate}

\textbf{Exercices de langue :}
\begin{enumerate}
    \setcounter{enumi}{4}
    \item Trouve dans le texte l'équivalent chinois de : « étudiant en deuxième année » ; « théorie du commerce international » ; « joindre son CV à la lettre ».
    \item Transforme la phrase « 我通过了HSK 4级考试 » en utilisant la structure \zh{考过} (kǎoguò) + \zh{了} (le) pour marquer l'expérience passée.
    \item Remplace « 做事认真负责 » par un synonyme formel à 4 caractères (chengyu) souvent utilisé dans les CV.
\end{enumerate}
\end{exercice}

\begin{exercice}[Thème et Version : Textes courts (Type Bac)]
\textbf{A. Thème (Français $\rightarrow$ Chinois) — Environ 80 caractères}
Traduis le paragraphe suivant en chinois. Soigne l'ordre des mots, les mots-outils (了, 过, 的, 地, 得) et le vocabulaire du CV.
\begin{quote}
« Je m'appelle Marie. J'ai 20 ans. De 2021 à 2024, j'ai étudié la langue chinoise à l'Université de Yaoundé II. J'ai obtenu le HSK 4. Je parle couramment chinois et français. J'ai fait un stage de deux mois dans une entreprise de commerce à Douala en 2023. Je sais utiliser les logiciels de bureau. Je cherche un poste d'assistant commercial. Je suis dynamique et sérieuse. »
\end{quote}

\textbf{B. Version (Chinois $\rightarrow$ Français) — Traduction littéraire soignée}
Traduis le texte suivant en français correct.
\begin{textebox}[Texte à traduire]
我叫 保罗 (Bǎoluó)， 来自 喀麦隆 布埃亚。我 在 2020年 考入 雅温得 大学 学习 计算机 科学。大学 期间， 我 刻苦 学习 中文， 考过 HSK 5 级。我 会 编程 (Python, Java)， 也 懂 网络 维护。我 性格 开朗， 团队 合作 精神 强。去年 我 在 一家 中资 科技 公司 实习 三 个月， 表现 优秀。现 申请 贵公司 IT 技术 支持 岗位， 希望 能 为 公司 贡献 力量。
\end{textebox}
\end{exercice}

\begin{exercice}[Expression écrite : Rédaction d'un CV complet (Type Bac)]
\textbf{Consigne :} Rédige ton propre CV en chinois simplifié pour une candidature à un stage d'été ou un premier emploi (réel ou fictif).
\begin{itemize}
    \item \textbf{Longueur :} Environ 150 à 200 caractères chinois (hors titres de rubriques).
    \item \textbf{Structure obligatoire (5 rubriques) :}
    \begin{enumerate}
        \item \zh{个人信息} (Nom, âge, nationalité, contact, objectif de poste)
        \item \zh{教育背景} (Diplômes, établissement, dates, spécialité, mentions/HSK)
        \item \zh{工作/实习经历} (Dates, entreprise, poste, 2-3 missions principales avec verbes d'action)
        \item \zh{技能特长} (Langues + niveaux, Informatique, Permis, autres)
        \item \zh{自我评价} (3-4 adjectifs/qualités + motivation courte)
    \end{enumerate}
    \item \textbf{Contraintes linguistiques :}
    \begin{itemize}
        \item Utilise au moins 8 des caractères complexes vus en leçon : \zh{简历, 经验, 教育, 技能, 聘, 应, 验, 历, 简, 育}.
        \item Emploie correctement les connecteurs temporels (\zh{期间}, \zh{自...以来}, \zh{从...到...}) et les adverbes de degré (\zh{熟练}, \zh{精通}, \zh{良好}).
        \item Respecte l'ordre des mots standard (Sujet + Temps + Lieu + Verbe + Objet).
        \item Soigne l'écriture des caractères (propreté, proportion).
    \end{itemize}
\end{itemize}
\begin{center}
\textit{Écris ton CV sur ta copie d'examen en respectant la mise en page en colonnes ou en liste.}
\end{center}
\end{exercice}

\newpage

\coursec{Corrigés des exercices}

\begin{corrige}[Exercice 1 — QCM : Vocabulaire et structure du CV]
\begin{enumerate}
    \item \textbf{Réponse : a.} \zh{简历} (jiǎnlì) est le terme standard pour « CV ». \zh{履历} (lǚlì) existe mais désigne plutôt le « parcours, les antécédents » (registre plus formel) ; \zh{简介} (jiǎnjiè) signifie « brève présentation » (d'un produit, d'une entreprise).
    \item \textbf{Réponse : a.} \zh{工作经验} (gōngzuò jīngyàn) = « expérience professionnelle ». \zh{学习经历} = « parcours scolaire » ; \zh{爱好特长} = « loisirs et points forts » ; \zh{个人信息} = « informations personnelles ».
    \item \textbf{Réponse : d.} Les trois formulations sont correctes : \zh{2020年—2024年} (style CV, tiret), \zh{2020年到2024年} (courant), \zh{2020年至2024年} (formel, écrit). Dans une phrase rédigée, on préfère la structure \zh{从2020年到2024年} (de... à...).
    \item \textbf{Réponse : a.} \zh{聘} (pìn) signifie « engager, embaucher ». \zh{应聘} (yìngpìn) = « postuler (répondre à une offre d'embauche) » ; c'est le mot composé qui signifie « postuler ».
\end{enumerate}
\textbf{Point de vigilance :} ne pas confondre \zh{简历} (CV) et \zh{简单} (jiǎndān, simple) : même premier caractère, sens différents.
\end{corrige}

\begin{corrige}[Exercice 2 — Vrai ou Faux : Connaissance du CV chinois]
\begin{enumerate}
    \item \textbf{Faux.} La photo est fréquente et appréciée en Chine, mais elle se place traditionnellement en \textbf{haut à droite} du CV, pas en bas ; elle n'est pas juridiquement obligatoire.
    \item \textbf{Faux.} \zh{教育背景} (jiàoyù bèijǐng) couvre tout le parcours éducatif pertinent : lycée (\zh{高中}), université, formations. Pour un élève de Terminale, le baccalauréat y figure naturellement.
    \item \textbf{Faux.} Mentionner le niveau HSK est fortement recommandé car c'est une certification objective, mais ce n'est pas une obligation : on peut écrire \zh{中文：良好} (chinois : bon niveau).
    \item \textbf{Vrai.} L'ordre chronologique inverse (le plus récent d'abord) est la norme, en Chine comme en France : le recruteur voit immédiatement l'expérience la plus actuelle.
    \item \textbf{Faux.} \zh{熟练掌握} (shúliàn zhǎngwò, « maîtriser parfaitement ») s'emploie pour les langues mais aussi pour les logiciels, les techniques : \zh{熟练掌握办公软件} « maîtriser parfaitement les logiciels de bureau ».
\end{enumerate}
\end{corrige}

\begin{corrige}[Exercice 3 — Vocabulaire : Associer caractères, pinyin et sens]
Correspondances attendues (le tableau de l'énoncé sert de référence) :
\begin{itemize}
    \item \textbf{简 — jiǎn — simple ; simplifier} (radical \zh{⺮} bambou : les lattes de bambou des documents anciens)
    \item \textbf{历 — lì — calendrier ; expérience ; passer par} (radical \zh{厂})
    \item \textbf{经 — jīng — trame ; passer par ; expérience} (radical \zh{纟} soie)
    \item \textbf{验 — yàn — vérifier ; expérience ; test} (radical \zh{马} cheval, à gauche)
    \item \textbf{教 — jiāo/jiào — enseigner ; éduquer} (radical \zh{攵} frapper/agir, à droite)
    \item \textbf{育 — yù — élever ; éduquer} (radical \zh{月} chair, en bas)
    \item \textbf{技 — jì — technique ; compétence} (radical \zh{扌} main)
    \item \textbf{能 — néng — pouvoir ; capacité}
    \item \textbf{聘 — pìn — engager ; embaucher} (radical \zh{耳} oreille)
    \item \textbf{应 — yìng — répondre ; postuler} (radical \zh{广})
\end{itemize}
\textbf{Point de vigilance :} \zh{教} a deux lectures : jiāo (\zh{教书}, enseigner) et jiào (\zh{教育}, éducation). Dans \zh{教育背景}, on lit \emph{jiào}.
\end{corrige}

\begin{corrige}[Exercice 4 — Pinyin vers Caractères : Reconnaissance auditive/visuelle]
\begin{enumerate}
    \item jiǎnlì $\rightarrow$ \zh{简历} (clés : \zh{⺮} bambou pour 简 ; \zh{厂} pour 历)
    \item jīngyàn $\rightarrow$ \zh{经验} (clés : \zh{纟} soie pour 经 ; \zh{马} cheval pour 验)
    \item jiàoyù $\rightarrow$ \zh{教育} (clés : \zh{攵} pour 教 ; \zh{月} pour 育)
    \item jìnéng $\rightarrow$ \zh{技能} (clé : \zh{扌} main pour 技)
    \item yìngpìn $\rightarrow$ \zh{应聘} (clés : \zh{广} pour 应 ; \zh{耳} oreille pour 聘)
    \item gōngzuò $\rightarrow$ \zh{工作} (clé : \zh{亻} homme pour 作)
    \item zìwǒ $\rightarrow$ \zh{自我} (attention : l'énoncé écrivait « zìwo », la lecture correcte est \emph{zìwǒ})
    \item pínggū $\rightarrow$ \zh{评估} (clés : \zh{讠} parole pour 评 ; \zh{亻} pour 估)
    \item xiàngmù $\rightarrow$ \zh{项目} (clé : \zh{木} bois pour 项)
    \item jiǎnglì $\rightarrow$ \zh{奖励} (clés : \zh{大} pour 奖 ; \zh{力} force pour 励)
\end{enumerate}
\textbf{Pièges d'homophones :} jīngyàn n'est pas \zh{惊艳} (éblouissant) ; jiǎnlì n'est pas \zh{建立} (jiànlì, établir) ; xiàngmù n'est pas \zh{项目} confondu avec \zh{项目}/\zh{象} — vérifier toujours le sens dans le contexte du CV.
\end{corrige}

\begin{corrige}[Exercice 5 — Traduction : Thème et Version (phrases isolées)]
\begin{enumerate}
    \item \zh{我叫保罗，我是喀麦隆人。} \\ \emph{Wǒ jiào Bǎoluó, wǒ shì Kāmàilóng rén.} \\ « Je m'appelle Paul, je suis camerounais. » — Structure \zh{我叫} + nom ; nationalité = pays + \zh{人}.
    \item \zh{从2020年到2024年，我在巴富萨姆中学学习。} \\ \emph{Cóng 2020 nián dào 2024 nián, wǒ zài Bāfùsàmǔ zhōngxué xuéxí.} \\ Structure \zh{从...到...} ; le complément de lieu \zh{在} + lieu se place \textbf{avant} le verbe.
    \item \zh{我有工作经验。我在一家中国公司工作过。} \\ \emph{Wǒ yǒu gōngzuò jīngyàn. Wǒ zài yì jiā Zhōngguó gōngsī gōngzuò guo.} \\ \zh{过} marque l'expérience vécue ; classificateur \zh{家} pour les entreprises.
    \item \zh{我的技能包括中文、法文和英文。} \\ \emph{Wǒ de jìnéng bāokuò Zhōngwén, Fǎwén hé Yīngwén.} \\ « Mes compétences comprennent le chinois, le français et l'anglais. » — \zh{包括} = « comprendre, inclure » ; énumération avec \zh{、} et \zh{和} avant le dernier élément.
    \item \zh{我熟练掌握办公软件和编程基础。} \\ \emph{Wǒ shúliàn zhǎngwò bàngōng ruǎnjiàn hé biānchéng jīchǔ.} \\ « Je maîtrise parfaitement les logiciels de bureau et les bases de la programmation. »
    \item \zh{我的爱好是足球和音乐。我身体很好。} \\ \emph{Wǒ de àihào shì zúqiú hé yīnyuè. Wǒ shēntǐ hěn hǎo.} \\ « Mes loisirs sont le football et la musique. Je suis en bonne santé. » — En chinois on dit « le corps est très bien » : \zh{身体很好}.
    \item \zh{应聘职位：中文助教。联系电话：6XXXXXXXX。} \\ \emph{Yìngpìn zhíwèi : Zhōngwén zhùjiào. Liánxì diànhuà : 6XXXXXXXX.} \\ « Poste sollicité : assistant de chinois. Téléphone : 6XXXXXXXX. » — \zh{职位} = poste ; \zh{助教} = assistant d'enseignement.
    \item \zh{我希望获得一次面试机会。谢谢您的关注！} \\ \emph{Wǒ xīwàng huòdé yí cì miànshì jīhuì. Xièxie nín de guānzhù !} \\ « J'espère obtenir un entretien. Merci de votre attention. » — \zh{您} (vous, poli) est de rigueur dans un courrier formel.
\end{enumerate}
\textbf{Points de vigilance :} ne jamais placer le complément de lieu après le verbe (\zh{我学习在雅温得} est faux) ; \zh{过} se colle au verbe ; utiliser \zh{您} et non \zh{你} dans un CV/une lettre.
\end{corrige}

\begin{corrige}[Exercice 6 — Discrimination de caractères : Choix du bon graphème]
\begin{enumerate}
    \item \zh{简}单 (jiǎndān, simple) — \zh{间} (jiān, entre/pièce) n'a pas ce sens.
    \item \zh{历}史 (lìshǐ, histoire) — \zh{历} = passer par le temps ; phrase complète : \zh{他有在中国工作的丰富历史。} (sens : un long parcours).
    \item \zh{厉}害 (lìhai, redoutable, impressionnant) — \zh{厉} = sévère, intense. Phrase : \zh{这个经验很厉害。}
    \item \zh{技}能 (jìnéng, compétence) — radical \zh{扌} main : l'habileté de la main.
    \item \zh{枝} (zhī, branche) — radical \zh{木} bois : \zh{他折断了学校前面树的一根枝。} (il a cassé une branche).
    \item \zh{检}查 (jiǎnchá, contrôler la qualité) : \zh{我们要检查产品质量。} — \zh{检} = inspecter, contrôler.
    \item \zh{验} (ou \zh{检}) : pour un examen médical, l'usage standard est \zh{检查身体} ; \zh{验} convient pour « tester, analyser » (\zh{验血}, prise de sang). Réponse attendue : \zh{检}查 (\zh{医生要检查我的身体。}).
\end{enumerate}
\textbf{Astuce mémotechnique :} \zh{技} (main \zh{扌}) = compétence manuelle ; \zh{枝} (bois \zh{木}) = branche d'arbre. Le radical sémantique tranche toujours entre les homophones graphiques.
\end{corrige}

\begin{corrige}[Exercice 7 — Réorganisation de mots : Structure de phrase CV]
\begin{enumerate}
    \item \zh{2020年到2024年，我在雅温得中学学习。} \\ \emph{2020 nián dào 2024 nián, wǒ zài Yǎwēndé zhōngxué xuéxí.} \\ « De 2020 à 2024, j'ai étudié au Lycée de Yaoundé. » — Temps en tête, puis Sujet + \zh{在} + lieu + verbe.
    \item \zh{我会说中文、法文和英文。} \\ \emph{Wǒ huì shuō Zhōngwén, Fǎwén hé Yīngwén.} \\ « Je sais parler chinois, français et anglais. » — \zh{会} + verbe = savoir faire.
    \item \zh{我有两年的工作经验。} \\ \emph{Wǒ yǒu liǎng nián de gōngzuò jīngyàn.} \\ « J'ai deux ans d'expérience professionnelle. » — Complément de durée \zh{两年} + \zh{的} + nom ; on dit \zh{两} (et non \zh{二}) devant un classificateur/nom mesuré.
    \item \zh{我熟练使用微软办公软件。} \\ \emph{Wǒ shúliàn shǐyòng Wēiruǎn bàngōng ruǎnjiàn.} \\ « J'utilise couramment les logiciels Microsoft Office. » — L'adverbe \zh{熟练} précède le verbe \zh{使用}.
    \item \zh{我的爱好是足球和音乐。} \\ \emph{Wǒ de àihào shì zúqiú hé yīnyuè.} \\ « Mes loisirs sont le football et la musique. » — Structure : possesseur + \zh{的} + nom + \zh{是} + attribut.
\end{enumerate}
\textbf{Règle à retenir :} Sujet + Temps + Lieu + Verbe + Objet ; les adverbes (\zh{熟练}, \zh{很}) se placent juste avant le verbe ou l'adjectif.
\end{corrige}

\begin{corrige}[Exercice 8 — Correction d'erreurs : CV d'un élève fictif]
Les 5 erreurs et leurs corrections :
\begin{enumerate}
    \item \textbf{Ordre des mots :} \zh{在一家餐厅服务员} est incorrect : il manque le verbe. Correction : \zh{在一家餐厅当服务员。} (\emph{zài yì jiā cāntīng dāng fúwùyuán}, « j'ai travaillé comme serveur dans un restaurant »). \zh{当} = « exercer la fonction de ».
    \item \textbf{Caractère / vocabulaire :} \zh{国籍：喀麦隆} — la rubrique nationalité attend \zh{喀麦隆} (pays) ou mieux \zh{喀麦隆籍} ; acceptable, mais la vraie erreur de registre est ailleurs : on retient plutôt l'erreur suivante.
    \item \textbf{Registre / précision :} \zh{英文 好} est trop vague et oral pour un CV. Correction : \zh{英文：良好} (\emph{Yīngwén : liánghǎo}, « anglais : bon niveau ») ou \zh{英语：流利}.
    \item \textbf{Grammaire (ordre des mots) :} \zh{希望给我一个机会} : le sujet implicite est ambigu et la tournure est impolie (impératif déguisé). Correction : \zh{希望贵公司给我一个机会。} ou mieux \zh{希望能获得一次机会。} (\emph{xīwàng néng huòdé yí cì jīhuì}).
    \item \textbf{Format des dates :} \zh{2020-2024} doit s'écrire en chinois : \zh{2020年—2024年} ou \zh{2020年至2024年}.
\end{enumerate}
\textbf{Texte corrigé (extrait) :} \\ \zh{姓名：Étienne；国籍：喀麦隆。教育经历：2020年至2024年，雅温得中学，高中毕业。工作经历：2022年至2023年，在一家餐厅当服务员。技能：中文HSK3级，法文母语，英文良好。爱好：看书、打篮球、听音乐。自我评价：我工作很努力，希望能获得一次机会。}
\end{corrige}

\begin{corrige}[Exercice 9 — Compréhension de l'écrit : Lettre de motivation + CV (Type Bac)]
\textbf{Barème indicatif (sur 20) :} questions de compréhension 8 pts (2 pts par question) ; exercices de langue 12 pts (4 pts par item).

\textbf{Questions de compréhension :}
\begin{enumerate}
    \item L'expéditeur est \textbf{Marco} (\zh{马可}, Mǎkě). Il a \textbf{19 ans}, il est \textbf{camerounais}, étudiant en \textbf{deuxième année à l'Université de Yaoundé I} (\zh{雅温得第一大学}), spécialité \textbf{commerce international} (\zh{国际贸易}).
    \item Objet : une \textbf{demande de stage} (\zh{实习}) dans l'entreprise, d'une durée de \textbf{deux mois}, de \textbf{juillet à août} (\zh{七月到八月}).
    \item Trois atouts : (a) les \textbf{langues} : HSK 4 en chinois, capacité de communication courante et commerciale simple ; (b) l'\textbf{informatique} : maîtrise d'Excel et PowerPoint ; (c) l'\textbf{expérience et le caractère} : organisation d'activités sur le campus, sérieux, sens des responsabilités, forte capacité d'apprentissage.
    \item La formule de fin est \zh{此致敬礼} (\emph{cǐ zhì jìng lǐ}) : « Je vous prie d'agréer l'expression de mes salutations respectueuses » (formule épistolaire formelle chinoise standard). On trouve aussi \zh{祝工作顺利} « je vous souhaite succès dans votre travail ».
\end{enumerate}

\textbf{Exercices de langue :}
\begin{enumerate}
    \setcounter{enumi}{4}
    \item « étudiant en deuxième année » = \zh{大二学生} ; « théorie du commerce international » = \zh{国际贸易理论} ; « joindre son CV à la lettre » = \zh{随信附上我的简历}.
    \item Transformation : \zh{我考过了HSK 4级。} (\emph{Wǒ kǎo guo le HSK sì jí.}) — \zh{考过} + \zh{了} marque l'expérience accomplie : « j'ai (déjà) passé et réussi le HSK 4 ».
    \item Synonyme formel en 4 caractères : \zh{一丝不苟} (\emph{yì sī bù gǒu}, « méticuleux, sans la moindre négligence ») ; on accepte aussi \zh{兢兢业业} (\emph{jīngjīng yèyè}, « consciencieux et dévoué »).
\end{enumerate}
\textbf{Points de vigilance :} \zh{贵公司} (« votre honorable entreprise ») est la désignation polie obligatoire ; \zh{随信附上} est la formule figée pour « ci-joint » ; ne pas traduire \zh{此致敬礼} mot à mot.
\end{corrige}

\begin{corrige}[Exercice 10 — Thème et Version : Textes courts (Type Bac)]
\textbf{Barème indicatif (sur 20) :} A. Thème 10 pts (vocabulaire 3, grammaire/ordre des mots 4, mots-outils 2, caractères 1) ; B. Version 10 pts (fidélité 5, qualité du français 3, vocabulaire 2).

\textbf{A. Proposition de thème :} \\
\zh{我叫玛丽，今年二十岁。从2021年到2024年，我在雅温得第二大学学习中文，获得了HSK 4级证书。我能流利地说中文和法文。2023年，我在杜阿拉一家贸易公司实习了两个月。我会使用办公软件。我应聘商务助理的职位。我性格开朗，工作认真。} \\
\emph{Wǒ jiào Mǎlì, jīnnián èrshí suì. Cóng 2021 nián dào 2024 nián, wǒ zài Yǎwēndé Dì-èr Dàxué xuéxí Zhōngwén, huòdé le HSK sì jí zhèngshū. Wǒ néng liúlì de shuō Zhōngwén hé Fǎwén. 2023 nián, wǒ zài Dù'ālā yì jiā màoyì gōngsī shíxí le liǎng ge yuè. Wǒ huì shǐyòng bàngōng ruǎnjiàn. Wǒ yìngpìn shāngwù zhùshǒu de zhíwèi. Wǒ xìnggé kāilǎng, gōngzuò rènzhēn.}

Justifications : \zh{从...到...} pour la période ; \zh{了} après \zh{获得} et \zh{实习} pour l'accompli ; durée \zh{两个月} placée après le verbe ; \zh{流利地} (ou \zh{流利地说}) avec \zh{地} devant le verbe ; \zh{应聘} + poste + \zh{的职位}.

\textbf{B. Proposition de version :} \\
« Je m'appelle Paul et je viens de Buea, au Cameroun. En 2020, j'ai été admis à l'Université de Yaoundé pour y étudier l'informatique. Pendant mes études universitaires, j'ai travaillé le chinois avec acharnement et j'ai obtenu le HSK 5. Je sais programmer (Python, Java) et je m'y connais aussi en maintenance des réseaux. J'ai un caractère ouvert et un fort esprit d'équipe. L'année dernière, j'ai effectué un stage de trois mois dans une entreprise technologique à capitaux chinois, où j'ai obtenu d'excellents résultats. Je postule actuellement à un poste de support technique informatique dans votre entreprise, et j'espère pouvoir apporter ma contribution à la société. »

\textbf{Points de vigilance :} \zh{考入} = « être admis (après concours) » ; \zh{中资公司} = « entreprise à capitaux chinois » ; \zh{贡献力量} = « apporter sa contribution » ; \zh{刻苦学习} = « étudier avec acharnement ».
\end{corrige}

\begin{corrige}[Exercice 11 — Expression écrite : Rédaction d'un CV complet (Type Bac)]
\textbf{Barème indicatif (sur 20) :} respect de la structure en 5 rubriques 4 pts ; richesse et exactitude du vocabulaire du thème (au moins 8 caractères cibles) 4 pts ; grammaire et ordre des mots 5 pts ; connecteurs temporels et adverbes de degré 3 pts ; qualité de l'écriture des caractères 2 pts ; cohérence et registre formel 2 pts.

\textbf{Proposition de rédaction modèle :}
\begin{textebox}[简历 — CV de Aïcha]
个人信息 \\
姓名：艾莎 (Àishā)；年龄：19岁；国籍：喀麦隆；电话：6XXXXXXXX \\
应聘职位：中文导游助理 \\
教育背景 \\
2020年至2024年：杜阿拉中学，高中毕业。 \\
自2022年以来学习中文，2024年获得HSK 4级证书。 \\
实习经历 \\
2023年7月到8月：在一家中资贸易公司实习，负责接待客户和翻译简单文件。工作期间，我学会了使用办公软件。 \\
技能特长 \\
中文：HSK 4级，能流利交流；法文：母语；英文：良好。 \\
熟练掌握微软办公软件；有驾驶执照。 \\
自我评价 \\
我性格开朗，做事认真负责，学习能力强。希望贵公司给我一次面试的机会，我一定努力工作，为公司贡献力量。
\end{textebox}
\emph{Pinyin (extraits clés) : Yìngpìn zhíwèi : Zhōngwén dǎoyóu zhùshǒu. — Zì 2022 nián yǐlái xuéxí Zhōngwén. — Wǒ shúliàn zhǎngwò Wēiruǎn bàngōng ruǎnjiàn.}

\textbf{Points de vigilance :}
\begin{itemize}
    \item Vérifier que les 8 caractères cibles apparaissent bien : \zh{简}历, \zh{经验}, \zh{教育}, \zh{技能}, \zh{应聘}, \zh{实习} (含\zh{习}), \zh{经历}.
    \item \zh{自...以来} exige un point de départ : \zh{自2022年以来} ; \zh{从...到...} exige deux bornes.
    \item \zh{期间} se place après le groupe nominal : \zh{工作期间}, jamais en tête de phrase seul.
    \item Utiliser \zh{贵公司} (et non \zh{你的公司}) pour désigner l'entreprise destinataire.
    \item Éviter les tournures francaises : « Je suis dynamique » se rend par \zh{我性格开朗/工作积极}, non par une calque littéral.
\end{itemize}
\end{corrige}

\newpage

\coursec{Fiche bilan}

\begin{popfigure}
\begin{tikzpicture}[mindmap, grow cyclic, every node/.style={concept, align=center, font=\small},
root concept/.append style={concept color=popblue, font=\bfseries},
level 1/.append style={level distance=3.4cm, sibling angle=60},
level 1 concept/.append style={font=\scriptsize}]
\node[root concept, align=center]{Leçon 40\\Rédiger un CV\\en chinois}
child[concept color=poporange]{node[align=center]{Rubriques du CV\\简历 jiǎnlì\\个人信息 · 教育\\经验 · 技能}}
child[concept color=popgreen]{node[align=center]{Caractères clés\\简历 经验\\教育 技能 聘}}
child[concept color=poppurple]{node[align=center]{Radicaux\\⺮(竹) 纟(丝)\\攵(文) 耳}}
child[concept color=poppink]{node[align=center]{Structures\\在 + lieu + 工作\\会 + verbe\\从…到…}}
child[concept color=popgold]{node[align=center]{Connecteurs\\首先 其次 然后\\最后 另外}}
child[concept color=popblue!60]{node[align=center]{Traduction\\thème / version\\ordre S-T-L-V\\compl. de temps}}
;
\end{tikzpicture}
\legende{Carte mentale de la leçon 40 : le CV chinois, ses rubriques, ses caractères et ses structures.}
\end{popfigure}

\begin{definition}[Lexique clé du CV (à connaître par cœur)]
\begin{center}
\begin{tabularx}{\linewidth}{|X|X|X|X|}
\hline
\rowcolor{popblueL} Caractères & Pinyin & Nature & Traduction \\
\hline
简历 & jiǎnlì & n. & CV, curriculum vitae \\
\hline
个人信息 & gèrén xìnxī & n. & informations personnelles \\
\hline
教育 & jiàoyù & n. & éducation, formation \\
\hline
经验 & jīngyàn & n. & expérience \\
\hline
技能 & jìnéng & n. & compétence, aptitude \\
\hline
应聘 & yìngpìn & v. & postuler, répondre à une offre \\
\hline
公司 & gōngsī & n. & entreprise, société \\
\hline
职位 & zhíwèi & n. & poste, fonction \\
\hline
毕业 & bìyè & v. & être diplômé, finir ses études \\
\hline
会 & huì & v. mod. & savoir (faire), être capable de \\
\hline
\end{tabularx}
\end{center}
Remarque de prononciation : dans \zh{应聘}, le caractère \zh{应} se lit \emph{yìng} (4\up{e} ton) quand il signifie « répondre à » ; il se lit \emph{yīng} dans \zh{应该} (yīnggāi, « devoir »).
\end{definition}

\begin{propriete}[Règles de grammaire et structures du CV]
\begin{center}
\begin{tabularx}{\linewidth}{|X|X|X|}
\hline
\rowcolor{popblueL} Structure & Exemple (caractères + pinyin) & Traduction \\
\hline
从 + début + 到 + fin & 我从2020年到2023年在雅温得上中学。Wǒ cóng 2020 nián dào 2023 nián zài Yǎwēndé shàng zhōngxué. & De 2020 à 2023, j'ai fait mes études secondaires à Yaoundé. \\
\hline
在 + lieu + 工作 / 学习 & 我在杜阿拉的一家公司工作。Wǒ zài Dù'ālā de yì jiā gōngsī gōngzuò. & Je travaille dans une entreprise de Douala. \\
\hline
会 + verbe (compétence) & 我会说汉语和法语。Wǒ huì shuō Hànyǔ hé Fǎyǔ. & Je sais parler chinois et français. \\
\hline
Sujet + 毕业于 + école & 我毕业于雅温得第二大学。Wǒ bìyè yú Yǎwēndé Dì-èr Dàxué. & Je suis diplômé de l'Université de Yaoundé II. \\
\hline
\end{tabularx}
\end{center}
Points de vigilance :
\begin{itemize}
\item L'ordre obligatoire est \cle{Sujet + Temps + Lieu + Verbe} : le complément de temps (\zh{从2020年到2023年}) précède le complément de lieu (\zh{在雅温得}), et les deux précèdent le verbe.
\item \zh{毕业于} (bìyè yú) : la préposition \zh{于} est figée dans cette expression ; on ne dit jamais \zh{毕业在}.
\item \zh{会} exprime une compétence \emph{acquise par l'apprentissage} (parler une langue, utiliser un logiciel) ; pour la permission ou la possibilité, on utilise \zh{可以} ou \zh{能}.
\end{itemize}
\end{propriete}

\begin{attention}[Écriture des caractères : points de vigilance]
\begin{itemize}
\item \zh{简} (jiǎn) : clé du bambou ⺮ en haut, \zh{间} en dessous ; écrire d'abord le haut, puis le bas.
\item \zh{历} (lì) : clé de la falaise 厂 à l'extérieur, \zh{力} à l'intérieur ; ne pas confondre avec \zh{厉} (lì, sévère), qui contient \zh{万}.
\item \zh{经} (jīng) : clé de la soie 纟 à gauche ; ne pas confondre avec \zh{轻} (qīng, léger) dont la clé est le char 车.
\item \zh{验} (yàn) : clé du cheval 马 à gauche ; ne pas confondre avec \zh{脸} (liǎn, visage) dont la clé est la chair 月.
\item \zh{教} (jiào) : \zh{孝} à gauche, clé du geste 攵 à droite ; écrire la gauche avant la droite.
\item \zh{技} (jì) : clé de la main 扌 à gauche ; ne pas confondre avec \zh{枝} (zhī, branche) dont la clé est le bois 木.
\item \zh{聘} (pìn) : clé de l'oreille 耳 à gauche ; dans 耳, tracer l'horizontal avant le vertical.
\end{itemize}
Règle générale d'écriture : de haut en bas, de gauche à droite, l'extérieur avant l'intérieur, l'horizontal avant le vertical.
\end{attention}

\begin{textebox}[Modèle de mini-CV commenté]
个人简历 (gèrén jiǎnlì) — « CV personnel »\\
我叫玛丽，今年二十二岁，住在雅温得。\\
Wǒ jiào Mǎlì, jīnnián èrshí'èr suì, zhù zài Yǎwēndé.\\
« Je m'appelle Marie, j'ai vingt-deux ans, j'habite à Yaoundé. » (个人信息)\\
我从2019年到2022年在雅温得上中学，毕业于雅温得第二大学。\\
Wǒ cóng 2019 nián dào 2022 nián zài Yǎwēndé shàng zhōngxué, bìyè yú Yǎwēndé Dì-èr Dàxué.\\
« De 2019 à 2022, j'ai fait mes études secondaires à Yaoundé et je suis diplômée de l'Université de Yaoundé II. » (教育经历)\\
我在杜阿拉的一家公司工作了两年。\\
Wǒ zài Dù'ālā de yì jiā gōngsī gōngzuò le liǎng nián.\\
« J'ai travaillé deux ans dans une entreprise de Douala. » (工作经验 ; \zh{了} marque ici l'accomplissement, la durée \zh{两年} se place après le verbe)\\
我会说汉语、法语和英语，也会用电脑。\\
Wǒ huì shuō Hànyǔ, Fǎyǔ hé Yīngyǔ, yě huì yòng diànnǎo.\\
« Je sais parler chinois, français et anglais, et je sais aussi utiliser l'ordinateur. » (技能)\\
我想应聘贵公司的职位。\\
Wǒ xiǎng yìngpìn guì gōngsī de zhíwèi.\\
« Je souhaite postuler à un poste dans votre entreprise. » (\zh{贵} guì est le préfixe honorifique « votre »)
\end{textebox}

\begin{methode}[Méthode express : rédiger un CV en chinois]
\begin{enumerate}
\item Titre : \zh{个人简历} (gèrén jiǎnlì, « CV personnel »).
\item Rubrique 1 : 个人信息 — nom, âge, adresse, téléphone, courriel.
\item Rubrique 2 : 教育经历 — diplômes avec 从…到… et 毕业于.
\item Rubrique 3 : 工作经验 — postes occupés avec 在 + lieu + 工作 (+ durée après le verbe).
\item Rubrique 4 : 技能 — langues et compétences avec 会 + verbe.
\item Conclusion : formule de candidature avec 应聘 + 职位.
\item Relire : ordre Sujet–Temps–Lieu–Verbe, tons, caractères complets, ponctuation chinoise ，。.
\end{enumerate}
\end{methode}

\begin{methode}[Méthode express : traduction thème et version]
\begin{enumerate}
\item Repérer le sujet, puis placer le complément de temps avant le complément de lieu, tous deux avant le verbe.
\item « savoir faire » se traduit par 会 + verbe ; « postuler à » par 应聘 ; « être diplômé de » par 毕业于.
\item En version, traduire 从…到… par « de … à … », 毕业于 par « diplômé de », et 贵公司 par « votre entreprise ».
\item Vérifier enfin que chaque caractère est complet (aucun trait manquant) et que la ponctuation est chinoise.
\end{enumerate}
\end{methode}

\begin{exoresolu}[Thème guidé]
Traduire en chinois : « De 2021 à 2023, j'ai étudié le chinois à Douala. Je sais parler chinois et je veux postuler à ce poste. »
\tcblower
我从2021年到2023年在杜阿拉学汉语。我会说汉语，我想应聘这个职位。\\
Wǒ cóng 2021 nián dào 2023 nián zài Dù'ālā xué Hànyǔ. Wǒ huì shuō Hànyǔ, wǒ xiǎng yìngpìn zhège zhíwèi.\\
Étapes : (1) sujet \zh{我} ; (2) complément de temps \zh{从2021年到2023年} ; (3) complément de lieu \zh{在杜阿拉} ; (4) verbe \zh{学汉语}. « savoir parler » = \zh{会说} ; « ce poste » = \zh{这个职位} (classificateur \zh{个}).
\end{exoresolu}

\begin{aretenir}[Checklist avant le Bac]
\begin{itemize}
\item[$\square$] Je sais écrire de mémoire les caractères 简历, 经验, 教育, 技能, 聘.
\item[$\square$] Je connais le radical de chacun de ces caractères (⺮, 纟, 马, 攵, 扌, 耳).
\item[$\square$] Je distingue 经 / 轻, 验 / 脸, 技 / 枝, 历 / 厉.
\item[$\square$] Je cite les quatre rubriques d'un CV chinois : 个人信息, 教育经历, 工作经验, 技能.
\item[$\square$] J'utilise correctement la structure 从…到… pour les périodes.
\item[$\square$] Je place le complément de temps et le complément de lieu avant le verbe.
\item[$\square$] J'exprime une compétence avec 会 + verbe (我会说汉语。).
\item[$\square$] Je sais dire « être diplômé de » avec 毕业于.
\item[$\square$] J'emploie les connecteurs 首先, 其次, 然后, 最后 dans une rédaction.
\item[$\square$] Je traduis « postuler à un poste » par 应聘 + 职位.
\item[$\square$] Je relis mon texte : tons du pinyin, ponctuation chinoise ，。, caractères sans trait manquant.
\item[$\square$] Je sais rédiger mon propre mini-CV en chinois en dix lignes maximum.
\end{itemize}
\end{aretenir}

\findecours

\end{document}
